\documentclass[10pt,letterpaper]{article}
\usepackage[margin=0.85in]{geometry}
\usepackage[T1]{fontenc}
\usepackage{lmodern,microtype,amsmath,amssymb,amsthm,mathtools}
\usepackage{booktabs,longtable,array,enumitem,xurl}
\usepackage[dvipsnames]{xcolor}
\usepackage[colorlinks=true,linkcolor=MidnightBlue,citecolor=MidnightBlue,urlcolor=MidnightBlue]{hyperref}
\usepackage{fancyhdr}
\setlength{\headheight}{14pt}
\setlength{\emergencystretch}{3em}
\pagestyle{fancy}\fancyhf{}
\fancyhead[L]{\small Pell residues and Diophantine consistency}
\fancyhead[R]{\small ProveIt research}
\fancyfoot[C]{\thepage}
\setlist{itemsep=3pt,topsep=5pt}
\newtheorem{theorem}{Theorem}[section]
\newtheorem{lemma}[theorem]{Lemma}
\newtheorem{proposition}[theorem]{Proposition}
\newtheorem{corollary}[theorem]{Corollary}
\theoremstyle{remark}\newtheorem{remark}[theorem]{Remark}
\newcommand{\N}{\mathbb N}\newcommand{\Z}{\mathbb Z}
\newcommand{\ZFC}{\mathrm{ZFC}}\newcommand{\ZF}{\mathrm{ZF}}
\newcommand{\Con}{\operatorname{Con}}\newcommand{\Prf}{\operatorname{Prf}}
\newcommand{\file}[1]{\texttt{\nolinkurl{#1}}}
\newcommand{\nvars}{395}
\newcommand{\nrows}{321}
\newcommand{\nterms}{7{,}323}
\newcommand{\nheight}{83{,}842{,}758}
\newcommand{\nlast}{320}
\newcommand{\nvarlast}{394}
\newcommand{\zvars}{663}
\newcommand{\zrows}{397}
\newcommand{\zterms}{7{,}268}
\newcommand{\zheight}{2{,}445{,}790{,}974{,}704}
\newcommand{\zlast}{396}
\newcommand{\zvarlast}{662}
\newcommand{\svars}{687}
\newcommand{\srows}{421}
\newcommand{\sterms}{5{,}602}
\newcommand{\sheight}{63{,}906{,}013}
\newcommand{\slast}{420}
\newcommand{\svarlast}{686}
\newcommand{\guardcount}{96}
\newcommand{\nreduction}{2.95}
\newcommand{\zreduction}{25.34}

\newcommand{\PN}{P_{\nvars}}\newcommand{\PZ}{Q_{\zvars}}\newcommand{\PS}{Q_{\svars}}
\hypersetup{pdftitle={Pell residues instead of power values: explicit Diophantine consistency equations},
 pdfauthor={Prepared for the ProveIt project},
 pdfsubject={Explicit ordinary polynomials, proofs, exact coefficients and reproducible reductions}}
\title{\vspace{-1em}\textbf{Pell residues instead of power values}\\[0.4em]
\large Explicit Diophantine consistency equations\\
in \nvars\ nonnegative integers and \zvars\ unrestricted integers}
\author{Prepared for Vladimir Reshetnikov's ProveIt project\\
\small AI-assisted mathematical research with reproducible exact algebra}
\date{23 October 2026}
\begin{document}
\maketitle
\begin{abstract}
We construct literal integer-coefficient quartics $\PN$ and $\PZ$ satisfying
\[
\ZFC\vdash
\bigl(\nexists x\in\N^{\nvars}\;\PN(x)=0\bigr)
\ \longleftrightarrow\ \Con(\ZFC)
\ \longleftrightarrow\
\bigl(\nexists z\in\Z^{\zvars}\;\PZ(z)=0\bigr).
\]
The natural equation has \nrows\ quadratic residuals, \nterms\ expanded
monomials, and coefficient height \nheight. A \svars-variable integer
companion retains smaller coefficients. Complete residual lists and fully
collected polynomial files are supplied.

The central new transformation replaces eleven full power values
$(q+1)^n$ by residues of indexed Pell solutions. A single congruence makes
$q+1$ a root of the common Pell characteristic polynomial modulo
$q^{H+1}$. Separating the growth lower bound from the actual Pell base
makes that congruence compatible with completeness. Exact lower digits
and binomial divisibility are then recovered without representing the
full powers. This reduces the underlying natural arithmetic graph from
690 to 655 coordinates. Combining it with previously established signed
Pell-index arguments and exact sign inference gives the integer reduction.

The comparison is against the current incoming 407/888/917-variable
direct quartics, not against obsolete milestones. These constructions
encode the calibrated numerical proof calculus $C_{222}$ directly; they
do not simulate K215 or instantiate Jones's universal polynomial.
Separate twelve-integer refinements improve the few-variable tradeoff.
We claim explicit improvements in stated measures, not global minimality
or worldwide priority. The proof is ordinary mathematics formalizable in
ZFC, composed with the identified inherited calculus calibration. No new
end-to-end Lean/Rocq proof term or emitted object-language ZFC derivation
is claimed.
\end{abstract}

\clearpage\tableofcontents\clearpage
\section{Results, comparison, and attribution}

Throughout, $\N=\{0,1,2,\ldots\}$; an integer variable has no sign
restriction. A polynomial's height is the maximum absolute coefficient
after full expansion and collection of equal monomials.

\begin{theorem}[Fixed ordinary equations]\label{thm:main}
The complete residual lists in Appendices~\ref{app:natural},
\ref{app:integer}, and~\ref{app:small} define
\[
\PN(x)=\sum_{j=0}^{\nlast}R_j(x)^2,\qquad
\PZ(z)=\sum_{j=0}^{\zlast}S_j(z)^2,\qquad
\PS(w)=\sum_{j=0}^{\slast}T_j(w)^2.
\]
Each polynomial has exact degree four. In ZFC, existence of a zero of
each polynomial over its stated domain is equivalent to existence of a
finite ZFC refutation. Consequently its no-zero sentence is equivalent,
provably in ZFC, to $\Con(\ZFC)$.
\end{theorem}

The term ``fixed'' excludes an unspecified machine index, a free source
parameter, or a coefficient chosen using an undecidable predicate. Every
coefficient and every variable is present in the delivered JSON. The
appendices specify the three principal quartics without any executable
macro. Powers, factorials, and bounded quantifiers below describe their
construction; they are absent from the final polynomial language.

\begin{table}[ht]\centering\small
\begin{tabular}{llrrrrr}
\toprule
Status & Domain & Variables & Rows & Degree & Monomials & Height\\
\midrule
Incoming & $\N$ & 407 & 323 & 4 & 8,420 & 83,843,774\\
This work & $\N$ & \nvars & \nrows & 4 & \nterms & \nheight\\
Incoming & $\Z$ & 888 & 470 & 4 & 10,156 & 2,445,790,974,704\\
This work & $\Z$ & \zvars & \zrows & 4 & \zterms & \zheight\\
Incoming & $\Z$ & 917 & 499 & 4 & 6,986 & 63,906,036\\
This work & $\Z$ & \svars & \srows & 4 & \sterms & \sheight\\
\bottomrule
\end{tabular}
\caption{Like-for-like direct quartic comparisons against the latest
inspected incoming archive. A row is a residual before squaring.}
\end{table}

The comparison snapshot is ProveIt commit
\file{ec80841c07e91be8956ce6ed4af9a9c0b782d3d9}. In particular, the incoming
\file{ConZFC_407_888_Quartics_2026-10-24.zip} was inspected directly
\cite{parent,repo}. The natural variable count decreases by
\nreduction\%; the principal integer count decreases by \zreduction\%.
The raw arithmetic graph has 655 coordinates and 581 quadratic rows,
including 12 exact power requests, three exact binomials, eight binomial
divisibility requests, and 15 independent Pell-index certificates.

\subsection{What is new here}

The new arithmetic step is the common modular-root construction in
Sections~\ref{sec:cofinal}--\ref{sec:digits}: freely increase the Pell
base above its growth lower bound, force a specified root modulo a
common power of the radix, and use Pell residues directly in digit
tests. The explicit mixed-domain integration, emitted coefficient lists,
and new twelve-integer centered circuits are also contributions of this
iteration.

The numerical calculus, its ZFC calibration, accumulator argument, CRT
compiler, factorial representation, recycled radix, additive Pell-index
reuse, and natural contraction method are inherited from the 407/888
package and its dependencies. Crucially, the signed auxiliary-index
theorem was already proved in \file{ConZFC_Shared_Pell_1873.zip}; it is
reused with credit, and independently checked here \cite{signed}.
Earlier sign-orbit and positivity arguments are also credited
\cite{signedcompiler}. Rediscovering those facts during this investigation
does not establish novelty. The archive preserves the applicable source
notices, original input ZIP, exact source hashes, and mathematical reviews.

\subsection{What ``simpler'' means in this report}

The repository already contains equations with far fewer variables,
including a 37-integer quartic using a different coefficient architecture,
and nine-positive-variable equations of high degree \cite{joint,repo}.
Our main results improve the \emph{direct quartics with moderate literal
coefficients and completely collected expansions}. They do not dominate
every construction in the repository. Section~\ref{sec:twelve} separately
improves the twelve-integer degree/operation frontier. There is no theorem
here that the global minimum number of variables, degree, or description
length has been found.

\section{The calibrated source and why it expresses refutability}
\label{sec:source}

The source is the numerical calculus $C_{222}$. Its free truth code is
0 and its false target is 1. Define
\[
 I(a,b)=2(a+b)(a+b+1)+4b,\quad
 M(a,b)=I(a,b)+1,\quad A(a,b)=I(a,b)+2,\quad E(a,b)=I(a,b)+3,
\]
and $N(a)=I(a,1)$, $X(a,b)=N(A(a,N(b)))$.
The constructors encode implication, membership, universal
quantification, equality, negation, and existential quantification in
the inherited arithmetic convention. Eight unconditional axiom families
and one guarded bundle are fixed polynomial compositions of these
constructors. Their complete definitions, the exact twelve branch rules,
and the source identity checks are contained in the preserved parent
archive, especially \file{upstream/provenance/upstream/calculus.tex}.

Each rule has two padded premises $p,q$ and one conclusion $v$. The
branches are truth padding $(0,0,0)$, the eight axioms $(0,0,c_j)$, the
guarded bundle $(0,0,c_*)$, modus ponens
$(a,I(a,b),b)$, and generalization $(b,0,A(a,b))$.
The bundle guard $c\ne a,b$ is represented by
\[
 (c-a)^2(c-b)^2=(d+1)^2.
\]
The four schema parameters $a,b,c,d$ are shared across the finite
disjunction; only the selected branch must hold.

The exact inherited semantic endpoint is
\begin{equation}\label{eq:calibration}
 C_{222}(1)\quad\longleftrightarrow\quad
 \exists p\,\Prf_{\ZF}(p,\bot)
 \quad\longleftrightarrow\quad
 \exists p\,\Prf_{\ZFC}(p,\bot).
\end{equation}
The named sources are \file{KEndpoint.c222_iff_zf_refutation} and
\file{ConstructibleChoice.zf_refutation_iff_zfc_refutation}, pinned to
commit \file{08200eb10fde2db408c69dffd840d6920a569b1f} in the parent
provenance. The second equivalence uses the constructible-universe
equiconsistency of ZF and ZFC. The first is a calculus theorem before
the machine-specific endpoint: this construction uses no transition table.
Our ordinary proof composes with these identified calibration results;
the complete transitive Lean development has not been recompiled here.

\subsection{A bounded certificate with eight columns}

For a sufficiently padded horizon $t=t_0+1$, let
\[
 M=M_0+1=\bigl((t+1)^t\bigr)^{2^{22}}.
\]
The parent label-majorant certificate shows that this carrier is large
enough for every label generated by the bounded schema parameters.
At row $i<t$, an accumulator update has the form
\begin{equation}\label{eq:flow}
 S'(1+vM)=S(1+pM)(1+qM).
\end{equation}
Project the rule's labels $p,q,v$ to their branch-specific polynomial
expressions. The only locally bounded columns are then
\[
 a,b,c,d,\ S,S',u,u'.
\]
The first four have bound $t$, the next two bound $y$, and the last two
bound $x$. The retained beta equations are
\[
 x=S+(1+(i+1)y)u,\qquad
 x=S'+(1+(i+2)y)u'.
\]
Because $S,S'\le y$, these are canonical remainder equations; consecutive
rows decode the same shared state. The endpoint rows give $S_0=M+1$ and
$S_t=1$. Completeness chooses $y$ sufficiently large and divisible by
$t!$, then chooses $x$ by CRT.

Multiplying \eqref{eq:flow} gives equality at $M$ of products of linear
factors. The carrier bound prevents carries and recovers their equality
as polynomials. Unique factorization recovers label multiplicities; the
inherited single-conclusion flow argument converts this certificate into
a finite Horn derivation of 1. Conversely a finite derivation can be
padded, scheduled, bounded, and encoded in this form. This is the
unchanged source theorem, not an inference from polynomial test samples.

\section{The arithmetic interface preserved by the new construction}
\label{sec:crt}

Put $R=2+t+x+y+M_0$ and $H=R^{2^{23}}$. The exponent is a fixed integer:
its implementation is 23 squarings. A checked monotone majorant for the
local polynomial has exponent 6,893,253, below $2^{23}$, and bounds the
local residual sum and every decoded coordinate. Let $T=H!$ and
\[
 B=\prod_{i=0}^{t-1}(1+(i+1)T).
\]
The explanatory product is eliminated by a finite arithmetic graph.
The CRT certificate has positive encoded column representatives
$a_1,\ldots,a_8$ and an index representative $C$, with
\begin{equation}\label{eq:crtcert}
 B\mid\binom{a_j}{u_j+1},\qquad
 B=1+(C+1)T,\qquad F(C,a_1,\ldots,a_8)\equiv0\pmod B.
\end{equation}
Here $F$ is the nonnegative local rule-and-beta polynomial and $u_j$ are
the heterogeneous bounds just stated.

For soundness, choose a prime divisor $p$ of $1+(i+1)T$. Since $p\nmid H!$,
we have $p>H$. The falling-factorial identity shows that
$B\mid\binom{a_j}{u_j+1}$ supplies a residue $w_j\in[0,u_j]$ modulo $p$.
The canonical index equation gives $C\equiv i\pmod p$. Thus
$F(i,w)\equiv0\pmod p$, while the majorant gives $0\le F(i,w)<p$;
therefore $F(i,w)=0$. Completeness uses the pairwise coprime moduli
$1+(i+1)T$ and finite CRT. The factorial $(u_j+1)!$ is coprime to $B$
and can be cancelled in the converse falling-factorial argument.

\subsection{Factorials and the finite product graph}

For $n\ge1$, write
\[
 \rho_n=(n+2)^{n+2},\qquad V_n=(\rho_n+1)^n,
 \qquad c_n=\binom{\rho_n+n}{n}.
\]
The inherited elementary identity is
\begin{equation}\label{eq:factorial}
 n!=\left\lceil V_n/c_n\right\rceil.
\end{equation}
Indeed $V_n/c_n=n!\prod_{i=1}^n(\rho_n+1)/(\rho_n+i)$ lies in
$(n!-1,n!]$. The estimate follows by bounding
$1-\prod(1-\alpha_i)$ by $\sum\alpha_i$, and then using
$n!n(n-1)/(2(\rho_n+1))<1$. The exact quadratic graph uses
$V_n+c_n=c_nf+r+1$ and $c_n=r+s+1$.

Set $Q=(1+(t+1)T)^t+2$ and $D=TQ-1$. The product graph uses
\begin{equation}\label{eq:product}
 B=\left(T^t t!\binom{Q+t}{t}\right)\bmod D.
\end{equation}
The numerator equals $\prod_{i=1}^t(TQ+iT)$ and hence reduces to the
desired product modulo $D$; the product is strictly smaller than $D$.
These identities explain the three exact binomial requests and two
factorial requests in the actual graph.

\subsection{Why the radix bound is not circular}

Choose the inherited recycled radix
\begin{equation}\label{eq:radix}
 q=B(V_H+1),\qquad \sum_{j=1}^8 a_j\le8B.
\end{equation}
Soundness first recovers all twelve retained exact powers from the growth
frame and indexed-Pell criteria of Section~\ref{sec:cofinal}. The modulus
row for the exact power $q^H$ forces $q>0$. Thus \eqref{eq:radix}, with
the recovered $V_H\ge1$, gives $B\ge1$ and $q\ge2$. The common Pell
residue identity in Lemma~\ref{lem:residue} now follows without using any
binomial or factorial interpretation. Only after these steps do the
following digit estimates recover the factorials and the CRT product.

Only the coefficients through the requested digit must be less than $q$.
For the factorial binomials, unimodality and the product formula give
\[
 \binom{\rho_n+n}{j}\le\binom{\rho_n+n}{n}\le V_n
 \qquad(0\le j\le n).
\]
First this proves the $H$-factorial digit and $T=H!$, then the $t$-factorial.
For the remaining binomials, use only the raw remainder inequality
$B<TQ$, without presuming \eqref{eq:product} is already valid. With
$e=t(H+1)$ one obtains
\[
 Q+t\le H^{e+1},\qquad
 a_j<8TQ\le H^{(t+2)(H+1)}.
\]
Their lower indices $k$ satisfy $k\le R$ and $k(t+2)\le H$. Hence
\[
 \binom{n}{j}\le n^k\le H^{H(H+1)}<V_H<q
 \qquad(0\le j\le k)
\]
for all remaining requests. Now the product binomial and all eight range
tests are justified. The order matters: using the intended product value
of $B$ earlier would be circular. The complete inherited estimates are
preserved in \file{proofs/inherited_radix_and_pell.txt}.

\section{A Pell base with a separately certified growth bound}
\label{sec:cofinal}

For $A\ge2$, define
\[
 X_0=1,\ Y_0=0,\qquad
 X_{j+1}=AX_j+(A^2-1)Y_j,\quad Y_{j+1}=X_j+AY_j.
\]
Then $(X_j,Y_j)$ are the nonnegative solutions of
$X^2-(A^2-1)Y^2=1$. Their recurrence, classification, and index
criterion are the standard Pell--Matiyasevich tools used by the parent
construction \cite{pell,carneiro}.

Let $\mathcal A$ be the finite set of distinct atoms occurring as a base,
exponent, or output of the \emph{12 retained exact-power requests}. Put
\begin{equation}\label{eq:lowerframe}
 W=2+\sum_{a\in\mathcal A}a,\quad Z=z+1,\qquad
 L^2-\bigl((W+1)^2-1\bigr)(WZ)^2=1,\qquad A=L+\sigma,
\end{equation}
with all displayed coordinates nonnegative. The parent put $A$ directly
in this growth norm. Introducing $L$ separates a lower bound from the
actual base.

\begin{lemma}[Growth and cofinal choice]\label{lem:growth}
Every solution of \eqref{eq:lowerframe} satisfies $A>W^W$.
Conversely, after fixing the semantic atoms, a suitable $L,Z$ exist and
every integer $A\ge L$ is admissible.
\end{lemma}
\begin{proof}
Pell classification gives $L=X_e(W+1)$ and $WZ=Y_e(W+1)$.
The recurrence modulo $W$ gives $Y_e(W+1)\equiv e\pmod W$.
Thus $W\mid e$ and $e>0$, so $e\ge W$. Consequently
$L\ge(W+1)^e\ge(W+1)^W>W^W$.
For completeness take $e=W$ and $Z=Y_W(W+1)/W$, then set $\sigma=A-L$.
\end{proof}

For a retained exact-power request $(n,k,m)$, the main indexed pair at
$A$ and modulus $T_n=2An-n^2-1$ still give
\begin{equation}\label{eq:exactpower}
 X_k+(n-A)Y_k\equiv m\pmod{T_n}.
\end{equation}
The frame implies $0\le m<T_n$ and $0<n^k<T_n$; the usual Pell
recurrence gives the same congruence with $n^k$ in place of $m$.
Therefore $m=n^k$. The modulus row with $T_n\ge0$ excludes a zero
base before any factorial or binomial interpretation is used.
All exact exponents are positive by their source expressions.

The new modular-only upper indices are \emph{not} required to be at
most $W$. They retain their explicit indexed-Pell constraints
$Y\ge k\ge1$. This distinction saves additional frame coordinates
without weakening any growth estimate used for an exact power.

\section{Replacing eleven full powers by one common residue condition}
\label{sec:digits}

Let $q\ge2$, $b=q+1$, and $M_q=q^{H+1}=q\,q^H$.
The exact denominator $q^H$ is already needed for the $H$-factorial
digit, so forming $M_q$ needs one multiplication. Add
\begin{equation}\label{eq:root}
 \boxed{\quad 2A(q+1)=(q+1)^2+1+M_qh_0.\quad}
\end{equation}

\begin{lemma}[Compatibility with completeness]
For fixed $q,H$ there are arbitrarily large integers $A$ satisfying
\eqref{eq:root}. They can be chosen above $L$ and $q+1$, with $h_0\ge0$.
\end{lemma}
\begin{proof}
Since $\gcd(q+1,M_q)=1$, the gcd of $2(q+1)$ and $M_q$ is either 1
or 2. In the even case $q+1$ is odd, so $(q+1)^2+1$ is even.
The linear congruence is therefore soluble. Its solution class contains
arbitrarily large positive representatives. Choose one large enough;
the resulting quotient is nonnegative.
\end{proof}

\begin{lemma}[Pell residue identity]\label{lem:residue}
Under \eqref{eq:root}, for every $n\in\N$,
\begin{equation}\label{eq:residue}
 \boxed{\quad \mathcal R_n=X_n(A)+(q+1-A)Y_n(A)
              \equiv(q+1)^n\pmod{q^{H+1}}.\quad}
\end{equation}
Moreover $\mathcal R_n>0$.
\end{lemma}
\begin{proof}
In $\Z[s]/(s^2-(A^2-1))$ we have
$(A+s)^n=X_n+Y_ns$. Equation~\eqref{eq:root} makes $s\mapsto q+1-A$
a homomorphism to $\Z/M_q\Z$, since
$(q+1-A)^2-(A^2-1)=(q+1)^2-2A(q+1)+1$ vanishes there.
This proves the congruence. For $n>0$, the norm identity yields
\[
 X_n^2-(A-1)^2Y_n^2=1+2(A-1)Y_n^2>0,
\]
so $X_n>(A-1)Y_n$ and $\mathcal R_n>qY_n>0$.
For $n=0$, $\mathcal R_0=1$.
\end{proof}

\subsection{Exact digits}

For a requested lower index $1\le k\le H$, set $d=q^k$.
The new exact-digit rows are
\begin{equation}\label{eq:digit}
 X_n+(q+1)Y_n=AY_n+r+dc+qd h,\qquad
 d=r+s+1,\qquad q=c+o+1.
\end{equation}
Thus $0\le r<d$ and $0\le c<q$. Because $qd\mid M_q$,
Lemma~\ref{lem:residue} identifies $r+dc$ with the least residue of
$(q+1)^n$ modulo $qd$. The binomial theorem and the lower-digit bounds
of Section~\ref{sec:crt} force $c=\binom nk$.
Conversely those correct low digits provide the row: the quotient
$(\mathcal R_n-r-dc)/(qd)$ is integral, and is nonnegative because
$\mathcal R_n>0$ and $0\le r+dc<qd$.

\subsection{Divisibility without the digit value}

For a range test only $B\mid\binom nk$ is required. Since $B\mid q$,
we have $Bd\mid M_q$. Replace the exact-digit rows by
\begin{equation}\label{eq:divdigit}
 X_n+(q+1)Y_n=AY_n+r+Bd h,\qquad d=r+s+1.
\end{equation}
The least residue of $(q+1)^n$ modulo $Bd$ is
\[
 r_0+d\bigl(\binom nk\bmod B\bigr),\qquad 0\le r_0<d.
\]
Higher digits vanish modulo $B$ because $B\mid q$. Consequently
\eqref{eq:divdigit} holds exactly when $B\mid\binom nk$.
Again the quotient is nonnegative when constructed from the true
residue, by positivity of $\mathcal R_n$.

\subsection{Why this saves witnesses}

The parent represented the entire numerator $(q+1)^n$ for each of the
three exact and eight divisibility requests, then extracted its low
digits. The new construction retains the already required indexed Pell
pairs and their exact denominators, but removes those eleven numerator
power graphs. The full-power count falls from 23 to 12. The new lower
frame and root congruence cost a small common amount, yielding the net
raw reduction $690\to655$. The final natural contraction is
$655\to\nvars$.

The order of witness selection is explicit: choose a finite source
certificate, its true arithmetic values and radix; choose $W,L$;
choose $A$ in the soluble root class above $L$; then choose common
auxiliary Pell/CRT witnesses and digit quotients. There is no equation
requiring a semantic output to be fixed before a quantity on which it
depends has been chosen.

\section{Signed auxiliary Pell witnesses: the inherited theorem}
\label{sec:signed}

We restate the sign argument used in the integer compilation, with
credit to the earlier incoming proof \cite{signed}. Its role here is
to permit a smaller integer-domain implementation of the new residue
construction.

\begin{theorem}[Signed auxiliary index recovery]\label{thm:signed}
Let $A,b\ge2$, $X,Y,V\in\N$, $1\le k\le Y$, and $U,S,R\in\Z$.
Suppose
\begin{align*}
 X^2-(A^2-1)Y^2&=1,& U^2-(A^2-1)V^2&=1,& V&>0,\\
 S^2-(b^2-1)R^2&=1,&Y^2&\mid V,\\
 b&\equiv A\pmod U,& b&\equiv1\pmod{2Y},\\
 S&\equiv X\pmod U,&R&\equiv k\pmod{2Y}.
\end{align*}
Then $(X,Y)=(X_k(A),Y_k(A))$. In particular the conclusion applies
to the delivered system, which retains the stronger moduli $4Y$.
\end{theorem}

\begin{proof}
Classify the nonnegative main pair as $(X_i(A),Y_i(A))$, where $i\ge1$,
and $(|U|,V)$ as $(X_m(A),Y_m(A))$, where $m\ge1$.
The standard divisibility argument gives $Y_i\mid m$ and $i\le m$.
For completeness, first $Y_i\mid Y_m$ implies $i\mid m$: divide the
index by $i$, use the unit $X_i$ modulo $Y_i$, and the strict increase
of the smaller remainder ordinate. Writing $m=di$, expansion modulo
$Y_i^2$ gives
\[
 Y_m\equiv dX_i^{d-1}Y_i\pmod{Y_i^2}.
\]
Since $\gcd(X_i,Y_i)=1$, the stronger condition $Y_i^2\mid Y_m$
implies $Y_i\mid d$, hence $Y_i\mid m$.

Classify $(|S|,|R|)$ at base $b$ with index $j\ge0$. Congruence of
the base and the possibly negative sign of $S$ give
\[
 X_j(A)\equiv\pm X_i(A)\pmod{X_m(A)}.
\]
The Pell abscissa congruence lemma implies
$j\equiv\pm i\pmod{2m}$. One can derive its signed version from
the usual positive congruence theorem by
$X_{j+2m}\equiv-X_j\pmod{X_m}$. Equivalently fold $j$ to $[0,m]$:
for $m\ge2$, $2X_{m-1}<X_m$ makes the smaller positive residues
distinct even up to sign. The endpoint $i=m$ has residue zero and is
separate; when $m=1$, the already known $Y\mid m$ gives $i=k=Y=1$.

Since $Y\mid m$, we obtain $j\equiv\pm i\pmod{2Y}$.
The recurrence at $b\equiv1\pmod{2Y}$ gives
$Y_j(b)\equiv j\pmod{2Y}$, so the possibly negative sign of $R$
also gives $j\equiv\pm k\pmod{2Y}$.
Thus $i\equiv\pm k\pmod{2Y}$. Both $i$ and $k$ belong to $[1,Y]$.
Equality modulo $2Y$ forces $i=k$; in the opposite-sign case
$i+k=2Y$ forces $i=k=Y$. This proves the assertion.
\end{proof}

The argument preserves the sign of the \emph{main} pair. It permits
arbitrary auxiliary signs. Replacing $S,R$ by their absolute values
on the same assignment can break their congruences, so that operation
is not used as a witness-restoration claim. The signed theorem identifies
the main indices directly. If an entirely natural parent witness is
desired, its auxiliary frame can then be reconstructed using the
inherited simultaneous-completeness theorem.

\section{Integer domains, sign certificates, and exact projection}
\label{sec:integer}

The raw graph has 655 coordinates. The mixed-domain compiler retains
nonnegativity exactly where the arithmetic proof requires it. The
auxiliary $U,S_i,R_i$, divisibility quotients, and their product wires
can remain signed. Main Pell pairs, all needed source and radix
parameters, and all genuine remainder/inequality slacks retain their
required bounds. The domain certificate checks this partition explicitly.

\subsection{Which quotient signs are unnecessary}

For an exact power, recurrence gives the identity
\[
 X_k+(n-A)Y_k-n^k
 =T_n\sum_{j=0}^{k-1}n^{k-1-j}Y_j.
\]
An arbitrary integer quotient therefore suffices for the congruence
proof. After the bounded residue is identified with $n^k$, this quotient
is the displayed nonnegative sum. The eleven modular extraction
quotients are also determined by a positive number and its least
nonnegative residue. The common root quotient is positive because
$A\ge q+1$ and $2A(q+1)-(q+1)^2-1>0$.

Three further restorations are elementary. The initial endpoint has
$x=(M_0+2)+(y+1)u$ with $0\le M_0+2<y+1$ and $x\ge0$, forcing $u\ge0$.
The product graph has $N=B+Du$, $0\le B<D$, $N\ge0$, forcing its
quotient nonnegative. Finally $F=Bv$, $F\ge0$, $B\ge1$ forces $v\ge0$.
These conclusions do not presume the product or binomial semantics
before they have been established.

\subsection{Encoding only the necessary signs}

For an integer $v$,
\begin{equation}\label{eq:ternary}
 v\ge0\quad\longleftrightarrow\quad
 \exists a,b,c\in\Z:\ v=a^2+b^2+c(c+1).
\end{equation}
The right side is nonnegative. Conversely, apply the three-square
theorem to $4v+1$, which is never of the excluded form $4^r(8s+7)$.
Exactly one square is odd modulo four; rearranging yields
\eqref{eq:ternary}. Each retained sign guard therefore costs three
integer witnesses and one quadratic row.

The final certificate needs \guardcount\ guards. Exact implications
from existing rows propagate the other required signs. They include
positive monomial sums, affine squares, and the inherited implication
$xA=B$ with a certified $A>0$ and an appropriate lower bound on $B$.
The checker validates every antecedent before admitting its conclusion;
no omitted quotient sign is used to justify a retained semantic sign.
The full finite ledger and its proof rules are included with each
integer equation.

\subsection{Projection and the single quartic}

A monic defining equation $v-f=0$ over $\Z$ allows exact elimination
of $v$. Substitution in the remaining rows and reverse restoration
$v=f$ prove existential equivalence. Over $\N$, an elimination must
additionally certify that its restoration is nonnegative; the parent
contractor uses positive monomial and rational sum-of-squares
certificates for that purpose. The independent verifier checks every
substitution as a polynomial identity and checks that residual degree
never exceeds two.

The natural contraction deliberately retains the common sum $W$.
Eliminating that one variable gives a 394-variable alternative but
greatly increases the expanded term count. The selected \nvars-variable
form has only \nterms\ terms. The integer alternatives use two support
thresholds: one reduces variables, and the other preserves smaller
coefficients. These thresholds control expression growth; they are not
mathematical assumptions.

Finally, over either domain a finite sum of integer squares is zero
exactly when each residual is zero. This yields the literal quartics
of Theorem~\ref{thm:main}. Their exact degree and coefficient counts
are obtained by collecting all products in the squared sums, not by
estimating the size of the compiler.

\section{The internal-ZFC statement and its proof boundary}
\label{sec:internal}

The desired statement is a provable equivalence, not a proof that the
polynomial has no zeros. Combining the reductions gives
\[
 \exists\hbox{ polynomial zero}
 \ \longleftrightarrow\
 \exists\hbox{ bounded arithmetic certificate}
 \ \longleftrightarrow\
 C_{222}(1)
 \ \longleftrightarrow\
 \exists p\,\Prf_{\ZFC}(p,\bot).
\]
Negating this equivalence proves the required no-zero/consistency
equivalence in ZFC.

The new arguments are uniform integer identities, induction on finite
Pell recurrences, Euclidean division, CRT, and finite certificate
transformations. They can be carried out over ZFC's own $\omega$ and
its definable integers. In particular they are not claims proved only
for each externally standard numeral; the quantified proofs themselves
are formalizable in the theory. The same uniform interpretation is
required of the inherited calculus and arithmetic calibration.

We distinguish three kinds of evidence. The ordinary mathematical proof
supplies the unbounded equivalences, with the inherited dependencies
identified. Exact executable checkers establish the finite coefficient
identities, domains, restorations, degree, and literal output binding.
Small finite arithmetic fixtures catch transcription errors in the
lemmas. Neither kind of computation replaces the ordinary proof.
This package emits no new kernel-checked end-to-end theorem and no
encoded object-language ZFC proof term. That formalization is a separate
project task; it is not silently claimed complete.

\section{Two additional twelve-integer improvements}
\label{sec:twelve}

The direct quartics optimize a different region from the few-variable
constructions. To advance both regions, the archive also supplies two
centered polynomials in twelve unrestricted integers, using the literal
parent in the valuation-and-signed-certificates archive \cite{valuation}.

\begin{center}\small
\begin{tabular}{lrrr}
\toprule
Construction & Integer variables & Exact degree & Live operations\\
\midrule
Incoming rational form & 12 & 2400 & 215\\
Incoming integral form & 12 & 2912 & 209\\
Centered integral form & 12 & \textbf{2400} & \textbf{210}\\
Centered rational form & 12 & \textbf{2368} & 216\\
\bottomrule
\end{tabular}
\end{center}

These are exact polynomial translations of an unrestricted coordinate:
$n_{\rm old}=n+QV$ or $n_{\rm old}=n+V$. The expressions $Q,V$ do not
depend on $n$, so the inverse translations are also integral and
polynomial. Every zero is preserved in both directions, with no new
positivity or parity hypothesis. The useful cancellation identities are
\begin{align*}
 (Q-1)(SV+z)^2-S^2(QV^2-1)
 &= (Q-1)z^2+2(Q-1)SVz-S^2(V^2-1),\\
 (SQV+z)^2-S^2(Q-1)(QV^2-1)
 &= z^2+2SQVz+S^2(QV^2+Q-1).
\end{align*}
The integral implementation shares $QV$ and computes $QV^2=(QV)V$.
The exact weighted-degree calculation, nonzero leading coefficients,
complete circuit identities, inherited coefficient recipe, and standalone
checkers are supplied under \file{supplemental/integer12/}.

These are incremental improvements, not the central new encoding idea.
Their source architecture has different large fixed coefficients; the
moderate-height claim for the direct quartics does not apply to them.
Operation counts refer to live binary addition, subtraction, and
multiplication nodes in a shared arithmetic circuit. They are neither
expanded monomial counts nor binary coefficient lengths.

\section{Reproduction, inspection, and remaining research}

The archive includes the LaTeX source, compiled PDF, complete residual
and expanded coefficient files, deterministic builders, independent
checkers, mathematical reviews, corruption controls, source pins, and
the original direct-construction dependency archive with its notices.
Python 3.10 or later and the standard library suffice for the arithmetic
reproduction. A conventional LaTeX installation compiles the report.

\begin{verbatim}
python3 verify.py
python3 -O verify.py
python3 reproduce.py --output /path/to/new-empty-directory
pdflatex -interaction=nonstopmode -halt-on-error report.tex
pdflatex -interaction=nonstopmode -halt-on-error report.tex
\end{verbatim}

The verifier treats delivered equation files as inputs and reconstructs
coefficient collections independently. It rejects altered sign proofs,
invalid restoration identities, and unauthorized signed roles.
Reproduction builds the new source in an isolated directory from the
pinned parent plus the new modules. No program searches for a refutation
or a zero of the final polynomials.

The new route suggests a broader optimization principle: when a
Diophantine certificate consumes only a residue of an expensive
arithmetic value, represent that residue directly and prove simultaneous
completeness for the shared parameters. Exact values are needed only
where their magnitude or arithmetic meaning is actually consumed.
Here the extra freedom in the Pell base is what makes this principle
an exact construction instead of an unsupported identification of
independently chosen witnesses.

Further reductions may exploit additional relationships among requested
indices, fewer locally bounded proof parameters, or stronger sign
inference. Such changes must preserve their complete domains and
cofinal witness construction. None is assumed in the delivered result.
Global optimality, an ultimate few-variable record, and a fresh
end-to-end formal proof certificate remain outside what has been proved.

\begin{thebibliography}{20}
\bibitem{repo} V. Reshetnikov and ProveIt contributors,
\emph{DiophantineConsistency research}, comparison snapshot
\file{ec80841c07e91be8956ce6ed4af9a9c0b782d3d9}.
\url{https://github.com/VladimirReshetnikov/ProveIt/tree/ec80841c07e91be8956ce6ed4af9a9c0b782d3d9/Computability/DiophantineConsistency}.
\bibitem{parent} ProveIt incoming contribution,
\emph{Residue certificates and nonnegative gates for the consistency of ZFC},
\file{ConZFC_407_888_Quartics_2026-10-24.zip}, same snapshot.
\url{https://github.com/VladimirReshetnikov/ProveIt/blob/ec80841c07e91be8956ce6ed4af9a9c0b782d3d9/docs/incoming/ConZFC_407_888_Quartics_2026-10-24.zip}.
The complete archive is included in \file{inputs/}.
\bibitem{signed} ProveIt incoming contribution,
\file{ConZFC_Shared_Pell_1873.zip}, especially
\file{proofs/SIGNED_INDEX_THEOREM.md} and its independent review.
\url{https://github.com/VladimirReshetnikov/ProveIt/blob/ec80841c07e91be8956ce6ed4af9a9c0b782d3d9/docs/incoming/ConZFC_Shared_Pell_1873.zip}.
\bibitem{signedcompiler} ProveIt incoming contribution,
\file{ConZFC_Direct_Quartics_3776_3671_2026-10-23.zip}, especially the
sign-orbit and guard-minimization proofs.
\url{https://github.com/VladimirReshetnikov/ProveIt/blob/ec80841c07e91be8956ce6ed4af9a9c0b782d3d9/docs/incoming/ConZFC_Direct_Quartics_3776_3671_2026-10-23.zip}.
\bibitem{joint} ProveIt incoming contribution,
\file{ZFC_Diophantine_Joint_Guards_2026-10-24.zip}, same snapshot.
\url{https://github.com/VladimirReshetnikov/ProveIt/blob/ec80841c07e91be8956ce6ed4af9a9c0b782d3d9/docs/incoming/ZFC_Diophantine_Joint_Guards_2026-10-24.zip}.
\bibitem{valuation} ProveIt incoming contribution,
\file{ConZFC_Valuation_and_Signed_Certificates_2026-10-23.zip},
especially the literal \file{integer12} subtree.
\url{https://github.com/VladimirReshetnikov/ProveIt/blob/ec80841c07e91be8956ce6ed4af9a9c0b782d3d9/docs/incoming/ConZFC_Valuation_and_Signed_Certificates_2026-10-23.zip}.
\bibitem{pell} Mathlib contributors, \emph{Pell's equation and
Matiyasevic's theorem}, especially \file{eq_pell},
\file{dvd_of_ysq_dvd}, \file{xn_modEq_x2n_add},
\file{modEq_of_xn_modEq}, and \file{matiyasevic}.
\url{https://leanprover-community.github.io/mathlib4_docs/Mathlib/NumberTheory/PellMatiyasevic.html}.
\bibitem{carneiro} M. Carneiro, \emph{A Lean formalization of
Matiyasevi\v{c}'s theorem}, 2018, arXiv:1802.01795.
\url{https://arxiv.org/abs/1802.01795}.
\end{thebibliography}

\clearpage\appendix
\section{The complete \nvars-variable natural equation}\label{app:natural}

Each $x_i$ independently ranges over $\N$, including zero. The exact
single polynomial is
\[
 \boxed{\PN(x_0,\ldots,x_{\nvarlast})=
       \sum_{j=0}^{\nlast}R_j(x_0,\ldots,x_{\nvarlast})^2=0.}
\]
The following list contains every residual and every coefficient.
The machine-readable order is \file{equations/natural395/equation.json};
the fully collected expansion is in its sibling
\file{expanded_polynomial.json}. Auxiliary names used in the construction
do not add coordinates beyond the numbered variables in this list.
\begingroup\footnotesize
\setlength{\tabcolsep}{3pt}
\begin{longtable}{r@{\hspace{0.8em}}l}
\toprule
Residual & Polynomial (set each to zero)\\
\midrule
\endfirsthead
\toprule
Residual & Polynomial (continued)\\
\midrule
\endhead
\bottomrule\endfoot
$R_{0}$ & $\begin{aligned}[t]&-x_{3}^{2}+x_{4}\end{aligned}$\\[3pt]
$R_{1}$ & $\begin{aligned}[t]&-x_{4}^{2}+x_{5}\end{aligned}$\\[3pt]
$R_{2}$ & $\begin{aligned}[t]&-x_{5}^{2}+x_{6}\end{aligned}$\\[3pt]
$R_{3}$ & $\begin{aligned}[t]&-x_{6}^{2}+x_{7}\end{aligned}$\\[3pt]
$R_{4}$ & $\begin{aligned}[t]&-x_{7}^{2}+x_{8}\end{aligned}$\\[3pt]
$R_{5}$ & $\begin{aligned}[t]&-x_{8}^{2}+x_{9}\end{aligned}$\\[3pt]
$R_{6}$ & $\begin{aligned}[t]&-x_{9}^{2}+x_{10}\end{aligned}$\\[3pt]
$R_{7}$ & $\begin{aligned}[t]&-x_{10}^{2}+x_{11}\end{aligned}$\\[3pt]
$R_{8}$ & $\begin{aligned}[t]&-x_{11}^{2}+x_{12}\end{aligned}$\\[3pt]
$R_{9}$ & $\begin{aligned}[t]&-x_{12}^{2}+x_{13}\end{aligned}$\\[3pt]
$R_{10}$ & $\begin{aligned}[t]&-x_{13}^{2}+x_{14}\end{aligned}$\\[3pt]
$R_{11}$ & $\begin{aligned}[t]&-x_{14}^{2}+x_{15}\end{aligned}$\\[3pt]
$R_{12}$ & $\begin{aligned}[t]&-x_{15}^{2}+x_{16}\end{aligned}$\\[3pt]
$R_{13}$ & $\begin{aligned}[t]&-x_{16}^{2}+x_{17}\end{aligned}$\\[3pt]
$R_{14}$ & $\begin{aligned}[t]&-x_{17}^{2}+x_{18}\end{aligned}$\\[3pt]
$R_{15}$ & $\begin{aligned}[t]&-x_{18}^{2}+x_{19}\end{aligned}$\\[3pt]
$R_{16}$ & $\begin{aligned}[t]&-x_{19}^{2}+x_{20}\end{aligned}$\\[3pt]
$R_{17}$ & $\begin{aligned}[t]&-x_{20}^{2}+x_{21}\end{aligned}$\\[3pt]
$R_{18}$ & $\begin{aligned}[t]&-x_{21}^{2}+x_{22}\end{aligned}$\\[3pt]
$R_{19}$ & $\begin{aligned}[t]&-x_{22}^{2}+x_{23}\end{aligned}$\\[3pt]
$R_{20}$ & $\begin{aligned}[t]&-x_{23}^{2}+x_{24}\end{aligned}$\\[3pt]
$R_{21}$ & $\begin{aligned}[t]&-1-x_{2}+x_{24}^{2}\end{aligned}$\\[3pt]
$R_{22}$ & $\begin{aligned}[t]&-4-2x_{0}-x_{0}x_{2}-x_{0}x_{26}-2x_{2}-2x_{26}+x_{28}\end{aligned}$\\[3pt]
$R_{23}$ & $\begin{aligned}[t]&-2+x_{1}-x_{2}-x_{2}x_{25}-3x_{25}-x_{25}x_{26}\end{aligned}$\\[3pt]
$R_{24}$ & $\begin{aligned}[t]&-1+x_{1}-x_{27}-x_{27}x_{28}\end{aligned}$\\[3pt]
$R_{25}$ & $\begin{aligned}[t]&-25-10x_{0}-x_{0}^{2}-2x_{0}x_{1}-4x_{0}x_{2}-2x_{0}x_{26}-10x_{1}-x_{1}^{2}-4x_{1}x_{2}-2x_{1}x_{26}-20x_{2}-4x_{2}^{2}-4x_{2}x_{26}\\ &-10x_{26}-x_{26}^{2}+x_{38}\end{aligned}$\\[3pt]
$R_{26}$ & $\begin{aligned}[t]&-x_{38}^{2}+x_{39}\end{aligned}$\\[3pt]
$R_{27}$ & $\begin{aligned}[t]&-x_{39}^{2}+x_{40}\end{aligned}$\\[3pt]
$R_{28}$ & $\begin{aligned}[t]&-x_{40}^{2}+x_{41}\end{aligned}$\\[3pt]
$R_{29}$ & $\begin{aligned}[t]&-x_{41}^{2}+x_{42}\end{aligned}$\\[3pt]
$R_{30}$ & $\begin{aligned}[t]&-x_{42}^{2}+x_{43}\end{aligned}$\\[3pt]
$R_{31}$ & $\begin{aligned}[t]&-x_{43}^{2}+x_{44}\end{aligned}$\\[3pt]
$R_{32}$ & $\begin{aligned}[t]&-x_{44}^{2}+x_{45}\end{aligned}$\\[3pt]
$R_{33}$ & $\begin{aligned}[t]&-x_{45}^{2}+x_{46}\end{aligned}$\\[3pt]
$R_{34}$ & $\begin{aligned}[t]&-x_{46}^{2}+x_{47}\end{aligned}$\\[3pt]
$R_{35}$ & $\begin{aligned}[t]&-x_{47}^{2}+x_{48}\end{aligned}$\\[3pt]
$R_{36}$ & $\begin{aligned}[t]&-x_{48}^{2}+x_{49}\end{aligned}$\\[3pt]
$R_{37}$ & $\begin{aligned}[t]&-x_{49}^{2}+x_{50}\end{aligned}$\\[3pt]
$R_{38}$ & $\begin{aligned}[t]&-x_{50}^{2}+x_{51}\end{aligned}$\\[3pt]
$R_{39}$ & $\begin{aligned}[t]&-x_{51}^{2}+x_{52}\end{aligned}$\\[3pt]
$R_{40}$ & $\begin{aligned}[t]&-x_{52}^{2}+x_{53}\end{aligned}$\\[3pt]
$R_{41}$ & $\begin{aligned}[t]&-x_{53}^{2}+x_{54}\end{aligned}$\\[3pt]
$R_{42}$ & $\begin{aligned}[t]&-x_{54}^{2}+x_{55}\end{aligned}$\\[3pt]
$R_{43}$ & $\begin{aligned}[t]&-x_{55}^{2}+x_{56}\end{aligned}$\\[3pt]
$R_{44}$ & $\begin{aligned}[t]&-x_{56}^{2}+x_{57}\end{aligned}$\\[3pt]
$R_{45}$ & $\begin{aligned}[t]&-x_{57}^{2}+x_{58}\end{aligned}$\\[3pt]
$R_{46}$ & $\begin{aligned}[t]&-x_{58}^{2}+x_{59}\end{aligned}$\\[3pt]
$R_{47}$ & $\begin{aligned}[t]&-x_{59}^{2}+x_{60}\end{aligned}$\\[3pt]
$R_{48}$ & $\begin{aligned}[t]&-x_{61}-x_{61}x_{64}-x_{61}x_{65}+x_{63}+x_{65}\end{aligned}$\\[3pt]
$R_{49}$ & $\begin{aligned}[t]&-x_{0}x_{61}-2x_{61}+x_{66}\end{aligned}$\\[3pt]
$R_{50}$ & $\begin{aligned}[t]&-x_{69}-x_{69}x_{72}-x_{69}x_{73}+x_{71}+x_{73}\end{aligned}$\\[3pt]
$R_{51}$ & $\begin{aligned}[t]&-x_{68}x_{69}+x_{75}\end{aligned}$\\[3pt]
$R_{52}$ & $\begin{aligned}[t]&2-2x_{61}-x_{61}x_{67}+x_{76}+x_{78}\end{aligned}$\\[3pt]
$R_{53}$ & $\begin{aligned}[t]&x_{74}x_{75}-x_{76}-x_{76}x_{77}-x_{77}-x_{77}x_{78}\end{aligned}$\\[3pt]
$R_{54}$ & $\begin{aligned}[t]&1+x_{37}x_{61}+x_{61}-x_{76}\end{aligned}$\\[3pt]
$R_{55}$ & $\begin{aligned}[t]&-6-5x_{29}-x_{29}^{2}-2x_{29}x_{30}-5x_{30}-x_{30}^{2}+x_{79}\end{aligned}$\\[3pt]
$R_{56}$ & $\begin{aligned}[t]&-72-68x_{30}-16x_{30}^{2}-16x_{30}x_{79}-34x_{79}-4x_{79}^{2}+x_{80}\end{aligned}$\\[3pt]
$R_{57}$ & $\begin{aligned}[t]&-30-11x_{29}-x_{29}^{2}-4x_{29}x_{80}-22x_{80}-4x_{80}^{2}+x_{81}\end{aligned}$\\[3pt]
$R_{58}$ & $\begin{aligned}[t]&-380-312x_{80}-64x_{80}^{2}-32x_{80}x_{81}-78x_{81}-4x_{81}^{2}+x_{82}\end{aligned}$\\[3pt]
$R_{59}$ & $\begin{aligned}[t]&-6-5x_{30}-x_{30}^{2}-2x_{30}x_{31}-5x_{31}-x_{31}^{2}+x_{83}\end{aligned}$\\[3pt]
$R_{60}$ & $\begin{aligned}[t]&-6-5x_{29}-x_{29}^{2}-2x_{29}x_{31}-5x_{31}-x_{31}^{2}+x_{84}\end{aligned}$\\[3pt]
$R_{61}$ & $\begin{aligned}[t]&-414-22x_{29}-2x_{29}^{2}-16x_{29}x_{31}-8x_{29}x_{83}-200x_{30}-32x_{30}^{2}-64x_{30}x_{31}-32x_{30}x_{79}-32x_{30}x_{84}-320x_{31}\\ &-64x_{31}^{2}-32x_{31}x_{79}-32x_{31}x_{83}-32x_{31}x_{84}-100x_{79}-8x_{79}^{2}-16x_{79}x_{84}-52x_{83}-8x_{83}^{2}-108x_{84}-8x_{84}^{2}\\ &+x_{85}\end{aligned}$\\[3pt]
$R_{62}$ & $\begin{aligned}[t]&-1344-800x_{30}-128x_{30}^{2}-256x_{30}x_{31}-128x_{30}x_{79}-128x_{30}x_{84}-864x_{31}-128x_{31}^{2}-128x_{31}x_{79}\\ &-128x_{31}x_{84}-400x_{79}-32x_{79}^{2}-64x_{79}x_{84}-432x_{84}-32x_{84}^{2}-2x_{85}-2x_{85}^{2}+x_{86}\end{aligned}$\\[3pt]
$R_{63}$ & $\begin{aligned}[t]&-247-168x_{29}-32x_{29}^{2}-64x_{29}x_{30}-32x_{29}x_{83}-32x_{29}x_{84}-188x_{30}-32x_{30}^{2}-32x_{30}x_{83}-32x_{30}x_{84}\\ &-2x_{79}-92x_{83}-8x_{83}^{2}-16x_{83}x_{84}-84x_{84}-8x_{84}^{2}+x_{87}\end{aligned}$\\[3pt]
$R_{64}$ & $\begin{aligned}[t]&-960-672x_{29}-128x_{29}^{2}-256x_{29}x_{30}-128x_{29}x_{83}-128x_{29}x_{84}-736x_{30}-128x_{30}^{2}-128x_{30}x_{83}\\ &-128x_{30}x_{84}-368x_{83}-32x_{83}^{2}-64x_{83}x_{84}-336x_{84}-32x_{84}^{2}-2x_{87}-2x_{87}^{2}+x_{88}\end{aligned}$\\[3pt]
$R_{65}$ & $\begin{aligned}[t]&-276-30x_{29}-2x_{29}^{2}-16x_{29}x_{31}-8x_{29}x_{83}-30x_{30}-2x_{30}^{2}-16x_{30}x_{31}-8x_{30}x_{84}-272x_{31}-64x_{31}^{2}\\ &-32x_{31}x_{83}-32x_{31}x_{84}-68x_{83}-8x_{83}^{2}-68x_{84}-8x_{84}^{2}+x_{89}\end{aligned}$\\[3pt]
$R_{66}$ & $\begin{aligned}[t]&-552-120x_{30}-8x_{30}^{2}-64x_{30}x_{31}-32x_{30}x_{84}-544x_{31}-128x_{31}^{2}-128x_{31}x_{84}-272x_{84}-32x_{84}^{2}-2x_{89}\\ &-2x_{89}^{2}+x_{90}\end{aligned}$\\[3pt]
$R_{67}$ & $\begin{aligned}[t]&-247-4x_{29}-352x_{30}-128x_{30}^{2}-64x_{30}x_{79}-64x_{30}x_{83}-92x_{79}-8x_{79}^{2}-16x_{79}x_{83}-84x_{83}-8x_{83}^{2}-2x_{84}\\ &+x_{91}\end{aligned}$\\[3pt]
$R_{68}$ & $\begin{aligned}[t]&-960-1408x_{30}-512x_{30}^{2}-256x_{30}x_{79}-256x_{30}x_{83}-368x_{79}-32x_{79}^{2}-64x_{79}x_{83}-336x_{83}-32x_{83}^{2}\\ &-2x_{91}-2x_{91}^{2}+x_{92}\end{aligned}$\\[3pt]
$R_{69}$ & $\begin{aligned}[t]&-56-60x_{30}-16x_{30}^{2}-16x_{30}x_{79}-30x_{79}-4x_{79}^{2}+x_{93}\end{aligned}$\\[3pt]
$R_{70}$ & $\begin{aligned}[t]&-30-11x_{29}-x_{29}^{2}-4x_{29}x_{93}-22x_{93}-4x_{93}^{2}+x_{94}\end{aligned}$\\[3pt]
$R_{71}$ & $\begin{aligned}[t]&-380-312x_{93}-64x_{93}^{2}-32x_{93}x_{94}-78x_{94}-4x_{94}^{2}+x_{95}\end{aligned}$\\[3pt]
$R_{72}$ & $\begin{aligned}[t]&-30-11x_{30}-x_{30}^{2}-4x_{30}x_{95}-22x_{95}-4x_{95}^{2}+x_{96}\end{aligned}$\\[3pt]
$R_{73}$ & $\begin{aligned}[t]&-349-420x_{30}-128x_{30}^{2}-128x_{30}x_{79}-208x_{79}-32x_{79}^{2}-2x_{83}+x_{97}\end{aligned}$\\[3pt]
$R_{74}$ & $\begin{aligned}[t]&-1368-1664x_{30}-512x_{30}^{2}-512x_{30}x_{79}-832x_{79}-128x_{79}^{2}-2x_{97}-2x_{97}^{2}+x_{98}\end{aligned}$\\[3pt]
$R_{75}$ & $\begin{aligned}[t]&-56-60x_{29}-16x_{29}^{2}-16x_{29}x_{84}-8x_{29}x_{98}-30x_{84}-4x_{84}^{2}-4x_{84}x_{98}-15x_{98}-x_{98}^{2}+x_{99}\end{aligned}$\\[3pt]
$R_{76}$ & $\begin{aligned}[t]&-6-5x_{31}-x_{31}^{2}-2x_{31}x_{32}-5x_{32}-x_{32}^{2}+x_{100}\end{aligned}$\\[3pt]
$R_{77}$ & $\begin{aligned}[t]&-224-26x_{29}-2x_{29}^{2}-16x_{29}x_{32}-8x_{29}x_{100}-288x_{31}-128x_{31}^{2}-64x_{31}x_{83}-64x_{31}x_{84}-88x_{32}-32x_{32}^{2}\\ &-32x_{32}x_{100}-76x_{83}-8x_{83}^{2}-16x_{83}x_{84}-68x_{84}-8x_{84}^{2}-44x_{100}-8x_{100}^{2}+x_{101}\end{aligned}$\\[3pt]
$R_{78}$ & $\begin{aligned}[t]&-640-1152x_{31}-512x_{31}^{2}-256x_{31}x_{83}-256x_{31}x_{84}-304x_{83}-32x_{83}^{2}-64x_{83}x_{84}-272x_{84}-32x_{84}^{2}\\ &-2x_{101}-2x_{101}^{2}+x_{102}\end{aligned}$\\[3pt]
$R_{79}$ & $\begin{aligned}[t]&-241-168x_{29}-32x_{29}^{2}-64x_{29}x_{30}-32x_{29}x_{83}-32x_{29}x_{84}-184x_{30}-32x_{30}^{2}-32x_{30}x_{83}-32x_{30}x_{84}-x_{31}\\ &-92x_{83}-8x_{83}^{2}-16x_{83}x_{84}-84x_{84}-8x_{84}^{2}+x_{103}\end{aligned}$\\[3pt]
$R_{80}$ & $\begin{aligned}[t]&-42-52x_{31}-16x_{31}^{2}-16x_{31}x_{83}-26x_{83}-4x_{83}^{2}+x_{104}\end{aligned}$\\[3pt]
$R_{81}$ & $\begin{aligned}[t]&-30-11x_{31}-x_{31}^{2}-4x_{31}x_{104}-22x_{104}-4x_{104}^{2}+x_{105}\end{aligned}$\\[3pt]
$R_{82}$ & $\begin{aligned}[t]&-306-35x_{30}-x_{30}^{2}-16x_{30}x_{104}-4x_{30}x_{105}-280x_{104}-64x_{104}^{2}-32x_{104}x_{105}-70x_{105}-4x_{105}^{2}+x_{106}\end{aligned}$\\[3pt]
$R_{83}$ & $\begin{aligned}[t]&-4556-4320x_{104}-1024x_{104}^{2}-512x_{104}x_{105}-128x_{104}x_{106}-1080x_{105}-64x_{105}^{2}-32x_{105}x_{106}-270x_{106}\\ &-4x_{106}^{2}+x_{107}\end{aligned}$\\[3pt]
$R_{84}$ & $\begin{aligned}[t]&-30-11x_{31}-x_{31}^{2}-4x_{31}x_{107}-22x_{107}-4x_{107}^{2}+x_{108}\end{aligned}$\\[3pt]
$R_{85}$ & $\begin{aligned}[t]&-306-35x_{30}-x_{30}^{2}-16x_{30}x_{107}-4x_{30}x_{108}-280x_{107}-64x_{107}^{2}-32x_{107}x_{108}-70x_{108}-4x_{108}^{2}+x_{109}\end{aligned}$\\[3pt]
$R_{86}$ & $\begin{aligned}[t]&-1028-672x_{29}-128x_{29}^{2}-256x_{29}x_{30}-128x_{29}x_{83}-128x_{29}x_{84}-736x_{30}-128x_{30}^{2}-128x_{30}x_{83}\\ &-128x_{30}x_{84}-368x_{83}-32x_{83}^{2}-64x_{83}x_{84}-336x_{84}-32x_{84}^{2}-2x_{103}-2x_{103}^{2}-32x_{107}-8x_{108}-2x_{109}\\ &+x_{110}\end{aligned}$\\[3pt]
$R_{87}$ & $\begin{aligned}[t]&-264-128x_{107}-32x_{108}-8x_{109}-2x_{110}-2x_{110}^{2}+x_{111}\end{aligned}$\\[3pt]
$R_{88}$ & $\begin{aligned}[t]&-2-3x_{29}-x_{29}^{2}-2x_{29}x_{111}-3x_{111}-x_{111}^{2}+x_{112}\end{aligned}$\\[3pt]
$R_{89}$ & $\begin{aligned}[t]&-12-28x_{111}-16x_{111}^{2}-16x_{111}x_{112}-14x_{112}-4x_{112}^{2}+x_{113}\end{aligned}$\\[3pt]
$R_{90}$ & $\begin{aligned}[t]&-30-11x_{31}-x_{31}^{2}-4x_{31}x_{113}-22x_{113}-4x_{113}^{2}+x_{114}\end{aligned}$\\[3pt]
$R_{91}$ & $\begin{aligned}[t]&-306-35x_{29}-x_{29}^{2}-16x_{29}x_{113}-4x_{29}x_{114}-280x_{113}-64x_{113}^{2}-32x_{113}x_{114}-70x_{114}-4x_{114}^{2}+x_{115}\end{aligned}$\\[3pt]
$R_{92}$ & $\begin{aligned}[t]&-4556-4320x_{113}-1024x_{113}^{2}-512x_{113}x_{114}-128x_{113}x_{115}-1080x_{114}-64x_{114}^{2}-32x_{114}x_{115}-270x_{115}\\ &-4x_{115}^{2}+x_{116}\end{aligned}$\\[3pt]
$R_{93}$ & $\begin{aligned}[t]&-90-76x_{31}-16x_{31}^{2}-16x_{31}x_{83}-16x_{31}x_{116}-38x_{83}-4x_{83}^{2}-8x_{83}x_{116}-38x_{116}-4x_{116}^{2}+x_{117}\end{aligned}$\\[3pt]
$R_{94}$ & $\begin{aligned}[t]&-4556-4320x_{107}-1024x_{107}^{2}-512x_{107}x_{108}-128x_{107}x_{109}-1080x_{108}-64x_{108}^{2}-32x_{108}x_{109}-270x_{109}\\ &-4x_{109}^{2}+x_{118}\end{aligned}$\\[3pt]
$R_{95}$ & $\begin{aligned}[t]&-682-70x_{30}-2x_{30}^{2}-32x_{30}x_{116}-8x_{30}x_{117}-592x_{116}-128x_{116}^{2}-64x_{116}x_{117}-148x_{117}-8x_{117}^{2}-2x_{118}\\ &+x_{119}\end{aligned}$\\[3pt]
$R_{96}$ & $\begin{aligned}[t]&-x_{119}-x_{119}^{2}+x_{120}\end{aligned}$\\[3pt]
$R_{97}$ & $\begin{aligned}[t]&-42-52x_{30}-16x_{30}^{2}-8x_{30}x_{31}-16x_{30}x_{79}-13x_{31}-x_{31}^{2}-4x_{31}x_{79}-26x_{79}-4x_{79}^{2}+x_{121}\end{aligned}$\\[3pt]
$R_{98}$ & $\begin{aligned}[t]&-2278-2264x_{30}-800x_{30}^{2}-800x_{30}x_{79}-160x_{30}x_{121}-70x_{31}-2x_{31}^{2}-32x_{31}x_{118}-8x_{31}x_{120}-1132x_{79}\\ &-200x_{79}^{2}-80x_{79}x_{121}-592x_{118}-128x_{118}^{2}-64x_{118}x_{120}-148x_{120}-8x_{120}^{2}-228x_{121}-8x_{121}^{2}+x_{122}\end{aligned}$\\[3pt]
$R_{99}$ & $\begin{aligned}[t]&-2712-280x_{31}-8x_{31}^{2}-128x_{31}x_{118}-32x_{31}x_{120}-2368x_{118}-512x_{118}^{2}-256x_{118}x_{120}-592x_{120}\\ &-32x_{120}^{2}-2x_{122}-2x_{122}^{2}+x_{123}\end{aligned}$\\[3pt]
$R_{100}$ & $\begin{aligned}[t]&-2-3x_{123}-x_{123}^{2}+x_{124}\end{aligned}$\\[3pt]
$R_{101}$ & $\begin{aligned}[t]&-1-x_{2}-x_{2}x_{34}-x_{34}+x_{125}\end{aligned}$\\[3pt]
$R_{102}$ & $\begin{aligned}[t]&-x_{33}^{2}+2x_{33}x_{34}-x_{34}^{2}+x_{126}\end{aligned}$\\[3pt]
$R_{103}$ & $\begin{aligned}[t]&-2x_{82}x_{125}-4x_{125}+x_{127}\end{aligned}$\\[3pt]
$R_{104}$ & $\begin{aligned}[t]&-x_{33}^{2}+2x_{33}x_{34}+2x_{33}x_{127}-x_{34}^{2}-2x_{34}x_{127}-x_{127}^{2}+x_{128}\end{aligned}$\\[3pt]
$R_{105}$ & $\begin{aligned}[t]&-x_{86}x_{125}+x_{129}\end{aligned}$\\[3pt]
$R_{106}$ & $\begin{aligned}[t]&-x_{33}^{2}+2x_{33}x_{34}+2x_{33}x_{129}-x_{34}^{2}-2x_{34}x_{129}-x_{129}^{2}+x_{130}\end{aligned}$\\[3pt]
$R_{107}$ & $\begin{aligned}[t]&-x_{88}x_{125}+x_{131}\end{aligned}$\\[3pt]
$R_{108}$ & $\begin{aligned}[t]&-x_{33}^{2}+2x_{33}x_{34}+2x_{33}x_{131}-x_{34}^{2}-2x_{34}x_{131}-x_{131}^{2}+x_{132}\end{aligned}$\\[3pt]
$R_{109}$ & $\begin{aligned}[t]&-x_{90}x_{125}+x_{133}\end{aligned}$\\[3pt]
$R_{110}$ & $\begin{aligned}[t]&-x_{33}^{2}+2x_{33}x_{34}+2x_{33}x_{133}-x_{34}^{2}-2x_{34}x_{133}-x_{133}^{2}+x_{134}\end{aligned}$\\[3pt]
$R_{111}$ & $\begin{aligned}[t]&-x_{92}x_{125}+x_{135}\end{aligned}$\\[3pt]
$R_{112}$ & $\begin{aligned}[t]&-x_{33}^{2}+2x_{33}x_{34}+2x_{33}x_{135}-x_{34}^{2}-2x_{34}x_{135}-x_{135}^{2}+x_{136}\end{aligned}$\\[3pt]
$R_{113}$ & $\begin{aligned}[t]&-4x_{30}x_{125}-2x_{96}x_{125}-4x_{125}+x_{137}\end{aligned}$\\[3pt]
$R_{114}$ & $\begin{aligned}[t]&-x_{33}^{2}+2x_{33}x_{34}+2x_{33}x_{137}-x_{34}^{2}-2x_{34}x_{137}-x_{137}^{2}+x_{138}\end{aligned}$\\[3pt]
$R_{115}$ & $\begin{aligned}[t]&-4x_{98}x_{125}-2x_{99}x_{125}+x_{139}\end{aligned}$\\[3pt]
$R_{116}$ & $\begin{aligned}[t]&-x_{33}^{2}+2x_{33}x_{34}+2x_{33}x_{139}-x_{34}^{2}-2x_{34}x_{139}-x_{139}^{2}+x_{140}\end{aligned}$\\[3pt]
$R_{117}$ & $\begin{aligned}[t]&-x_{102}x_{125}+x_{141}\end{aligned}$\\[3pt]
$R_{118}$ & $\begin{aligned}[t]&-x_{33}^{2}+2x_{33}x_{34}+2x_{33}x_{141}-x_{34}^{2}-2x_{34}x_{141}-x_{141}^{2}+x_{142}\end{aligned}$\\[3pt]
$R_{119}$ & $\begin{aligned}[t]&-x_{29}^{2}+2x_{29}x_{31}-x_{31}^{2}+x_{143}\end{aligned}$\\[3pt]
$R_{120}$ & $\begin{aligned}[t]&-x_{30}^{2}+2x_{30}x_{31}-x_{31}^{2}+x_{144}\end{aligned}$\\[3pt]
$R_{121}$ & $\begin{aligned}[t]&-x_{143}x_{144}+x_{145}\end{aligned}$\\[3pt]
$R_{122}$ & $\begin{aligned}[t]&-4-4x_{32}-x_{32}^{2}+x_{146}\end{aligned}$\\[3pt]
$R_{123}$ & $\begin{aligned}[t]&-x_{145}^{2}+2x_{145}x_{146}-x_{146}^{2}+x_{147}\end{aligned}$\\[3pt]
$R_{124}$ & $\begin{aligned}[t]&-2x_{124}x_{125}-4x_{125}+x_{148}\end{aligned}$\\[3pt]
$R_{125}$ & $\begin{aligned}[t]&-x_{33}^{2}+2x_{33}x_{34}+2x_{33}x_{148}-x_{34}^{2}-2x_{34}x_{148}-x_{148}^{2}+x_{149}\end{aligned}$\\[3pt]
$R_{126}$ & $\begin{aligned}[t]&-1-x_{2}-x_{2}x_{30}-x_{30}+x_{150}\end{aligned}$\\[3pt]
$R_{127}$ & $\begin{aligned}[t]&-1-x_{34}-x_{34}x_{150}-x_{150}+x_{151}\end{aligned}$\\[3pt]
$R_{128}$ & $\begin{aligned}[t]&-1-x_{2}-x_{2}x_{29}-x_{29}+x_{152}\end{aligned}$\\[3pt]
$R_{129}$ & $\begin{aligned}[t]&-4-4x_{2}-4x_{2}x_{30}-2x_{2}x_{79}-4x_{30}-2x_{79}+x_{153}\end{aligned}$\\[3pt]
$R_{130}$ & $\begin{aligned}[t]&-1-x_{33}-x_{33}x_{152}-x_{152}+x_{154}\end{aligned}$\\[3pt]
$R_{131}$ & $\begin{aligned}[t]&-x_{153}x_{154}-x_{154}+x_{155}\end{aligned}$\\[3pt]
$R_{132}$ & $\begin{aligned}[t]&-x_{151}^{2}+2x_{151}x_{155}-x_{155}^{2}+x_{156}\end{aligned}$\\[3pt]
$R_{133}$ & $\begin{aligned}[t]&-6-6x_{2}-4x_{2}x_{30}-2x_{2}x_{79}-4x_{30}-2x_{79}+x_{157}\end{aligned}$\\[3pt]
$R_{134}$ & $\begin{aligned}[t]&-1-x_{34}-x_{34}x_{157}-x_{157}+x_{158}\end{aligned}$\\[3pt]
$R_{135}$ & $\begin{aligned}[t]&-1-x_{33}-x_{33}x_{150}-x_{150}+x_{159}\end{aligned}$\\[3pt]
$R_{136}$ & $\begin{aligned}[t]&-x_{158}^{2}+2x_{158}x_{159}-x_{159}^{2}+x_{160}\end{aligned}$\\[3pt]
$R_{137}$ & $\begin{aligned}[t]&-x_{126}x_{128}+x_{161}\end{aligned}$\\[3pt]
$R_{138}$ & $\begin{aligned}[t]&-x_{130}x_{161}+x_{162}\end{aligned}$\\[3pt]
$R_{139}$ & $\begin{aligned}[t]&-x_{132}x_{162}+x_{163}\end{aligned}$\\[3pt]
$R_{140}$ & $\begin{aligned}[t]&-x_{134}x_{163}+x_{164}\end{aligned}$\\[3pt]
$R_{141}$ & $\begin{aligned}[t]&-x_{136}x_{164}+x_{165}\end{aligned}$\\[3pt]
$R_{142}$ & $\begin{aligned}[t]&-x_{138}x_{165}+x_{166}\end{aligned}$\\[3pt]
$R_{143}$ & $\begin{aligned}[t]&-x_{140}x_{166}+x_{167}\end{aligned}$\\[3pt]
$R_{144}$ & $\begin{aligned}[t]&-x_{142}x_{167}+x_{168}\end{aligned}$\\[3pt]
$R_{145}$ & $\begin{aligned}[t]&-x_{147}x_{168}-x_{149}x_{168}+x_{169}\end{aligned}$\\[3pt]
$R_{146}$ & $\begin{aligned}[t]&-x_{156}x_{169}+x_{170}\end{aligned}$\\[3pt]
$R_{147}$ & $\begin{aligned}[t]&-2-x_{2}-x_{2}x_{37}-x_{26}-x_{26}x_{37}-2x_{37}+x_{171}\end{aligned}$\\[3pt]
$R_{148}$ & $\begin{aligned}[t]&-1-x_{35}-x_{35}x_{171}-x_{171}+x_{172}\end{aligned}$\\[3pt]
$R_{149}$ & $\begin{aligned}[t]&-3-x_{2}-x_{2}x_{36}-x_{26}-x_{26}x_{36}-3x_{36}-x_{36}x_{171}-x_{171}+x_{173}\end{aligned}$\\[3pt]
$R_{150}$ & $\begin{aligned}[t]&2-4x_{1}+2x_{1}^{2}-2x_{1}x_{33}-2x_{1}x_{34}-2x_{1}x_{172}-2x_{1}x_{173}+2x_{33}+x_{33}^{2}+2x_{33}x_{172}+2x_{34}+x_{34}^{2}+2x_{34}x_{173}\\ &-x_{76}x_{174}+x_{160}x_{170}+2x_{172}+x_{172}^{2}+2x_{173}+x_{173}^{2}\end{aligned}$\\[3pt]
$R_{151}$ & $\begin{aligned}[t]&-8-x_{29}-x_{30}-x_{31}-x_{32}-x_{33}-x_{34}-x_{35}-x_{36}+8x_{76}-x_{175}\end{aligned}$\\[3pt]
$R_{152}$ & $\begin{aligned}[t]&2-x_{63}x_{76}+x_{64}+x_{65}-x_{76}+x_{263}\end{aligned}$\\[3pt]
$R_{153}$ & $\begin{aligned}[t]&-2-x_{64}-x_{64}x_{261}-x_{64}x_{262}-x_{65}-x_{65}x_{261}-x_{65}x_{262}+x_{176}-2x_{261}-x_{261}x_{263}-2x_{262}-x_{262}x_{263}-x_{263}\end{aligned}$\\[3pt]
$R_{154}$ & $\begin{aligned}[t]&-x_{177}-x_{177}x_{180}+x_{181}\end{aligned}$\\[3pt]
$R_{155}$ & $\begin{aligned}[t]&-x_{181}^{2}+x_{182}\end{aligned}$\\[3pt]
$R_{156}$ & $\begin{aligned}[t]&-1-2x_{177}-x_{177}^{2}+x_{183}\end{aligned}$\\[3pt]
$R_{157}$ & $\begin{aligned}[t]&-1+x_{178}^{2}+x_{182}-x_{182}x_{183}\end{aligned}$\\[3pt]
$R_{158}$ & $\begin{aligned}[t]&-10-6x_{64}-x_{64}^{2}-2x_{64}x_{65}+2x_{64}x_{178}+2x_{64}x_{179}-2x_{64}x_{263}-6x_{65}-x_{65}^{2}+2x_{65}x_{178}+2x_{65}x_{179}\\ &-2x_{65}x_{263}-x_{176}x_{184}+6x_{178}+2x_{178}x_{263}+6x_{179}+2x_{179}x_{263}-6x_{263}-x_{263}^{2}\end{aligned}$\\[3pt]
$R_{159}$ & $\begin{aligned}[t]&-x_{186}^{2}+x_{189}\end{aligned}$\\[3pt]
$R_{160}$ & $\begin{aligned}[t]&1-x_{178}^{2}-2x_{178}x_{179}-x_{179}^{2}+x_{211}\end{aligned}$\\[3pt]
$R_{161}$ & $\begin{aligned}[t]&-1+x_{185}^{2}-x_{189}x_{211}\end{aligned}$\\[3pt]
$R_{162}$ & $\begin{aligned}[t]&-4x_{186}-4x_{186}x_{187}+x_{190}\end{aligned}$\\[3pt]
$R_{163}$ & $\begin{aligned}[t]&1-x_{178}-x_{179}-x_{185}x_{188}+x_{190}\end{aligned}$\\[3pt]
$R_{164}$ & $\begin{aligned}[t]&-1-2x_{0}-x_{0}^{2}-2x_{0}x_{194}-2x_{194}-x_{194}^{2}+x_{195}\end{aligned}$\\[3pt]
$R_{165}$ & $\begin{aligned}[t]&-1+x_{191}^{2}-x_{195}x_{211}\end{aligned}$\\[3pt]
$R_{166}$ & $\begin{aligned}[t]&-x_{193}^{2}+x_{196}\end{aligned}$\\[3pt]
$R_{167}$ & $\begin{aligned}[t]&-1-2x_{190}-x_{190}^{2}+x_{197}\end{aligned}$\\[3pt]
$R_{168}$ & $\begin{aligned}[t]&-1+x_{192}^{2}+x_{196}-x_{196}x_{197}\end{aligned}$\\[3pt]
$R_{169}$ & $\begin{aligned}[t]&-x_{185}x_{198}-x_{191}+x_{192}\end{aligned}$\\[3pt]
$R_{170}$ & $\begin{aligned}[t]&-1-x_{0}-4x_{0}x_{199}+x_{193}-4x_{194}x_{199}-4x_{199}\end{aligned}$\\[3pt]
$R_{171}$ & $\begin{aligned}[t]&-5-4x_{0}-x_{0}^{2}+2x_{0}x_{178}+2x_{0}x_{179}+4x_{178}+4x_{179}-x_{200}\end{aligned}$\\[3pt]
$R_{172}$ & $\begin{aligned}[t]&-x_{0}x_{178}-x_{0}x_{179}-x_{178}-x_{178}x_{194}-x_{179}-x_{179}x_{194}+x_{202}\end{aligned}$\\[3pt]
$R_{173}$ & $\begin{aligned}[t]&2+3x_{0}+x_{0}^{2}+x_{0}x_{194}-x_{3}+x_{191}+2x_{194}-x_{200}x_{201}-x_{202}\end{aligned}$\\[3pt]
$R_{174}$ & $\begin{aligned}[t]&-x_{60}^{2}-2x_{60}x_{206}-x_{206}^{2}+x_{207}\end{aligned}$\\[3pt]
$R_{175}$ & $\begin{aligned}[t]&-1+x_{203}^{2}-x_{207}x_{211}\end{aligned}$\\[3pt]
$R_{176}$ & $\begin{aligned}[t]&-x_{205}^{2}+x_{208}\end{aligned}$\\[3pt]
$R_{177}$ & $\begin{aligned}[t]&-1-x_{197}x_{208}+x_{204}^{2}+x_{208}\end{aligned}$\\[3pt]
$R_{178}$ & $\begin{aligned}[t]&-x_{185}x_{209}-x_{203}+x_{204}\end{aligned}$\\[3pt]
$R_{179}$ & $\begin{aligned}[t]&-x_{60}-4x_{60}x_{210}+x_{205}-4x_{206}x_{210}\end{aligned}$\\[3pt]
$R_{180}$ & $\begin{aligned}[t]&-2x_{60}x_{178}-2x_{60}x_{179}-2x_{178}x_{206}-2x_{179}x_{206}+x_{212}\end{aligned}$\\[3pt]
$R_{181}$ & $\begin{aligned}[t]&-2x_{178}x_{203}-2x_{179}x_{203}+x_{213}\end{aligned}$\\[3pt]
$R_{182}$ & $\begin{aligned}[t]&-x_{60}-2x_{60}x_{211}-x_{206}-2x_{206}x_{211}+x_{214}\end{aligned}$\\[3pt]
$R_{183}$ & $\begin{aligned}[t]&-5-4x_{60}-x_{60}^{2}+2x_{60}x_{178}+2x_{60}x_{179}+4x_{178}+4x_{179}-x_{215}\end{aligned}$\\[3pt]
$R_{184}$ & $\begin{aligned}[t]&x_{60}x_{213}+x_{60}x_{214}-x_{62}-x_{178}x_{213}-x_{178}x_{214}-x_{179}x_{213}-x_{179}x_{214}+x_{203}+2x_{203}x_{211}+x_{211}x_{212}+2x_{213}\\ &+2x_{214}-x_{215}x_{216}\end{aligned}$\\[3pt]
$R_{185}$ & $\begin{aligned}[t]&-2-2x_{62}-x_{62}^{2}+2x_{62}x_{178}+2x_{62}x_{179}+2x_{178}+2x_{179}-x_{217}\end{aligned}$\\[3pt]
$R_{186}$ & $\begin{aligned}[t]&x_{60}+x_{60}x_{62}-x_{60}x_{178}-x_{60}x_{179}+x_{62}x_{206}-x_{63}-x_{178}x_{206}-x_{179}x_{206}+x_{203}+x_{206}-x_{217}x_{218}\end{aligned}$\\[3pt]
$R_{187}$ & $\begin{aligned}[t]&-2-2x_{66}-x_{66}^{2}+2x_{66}x_{178}+2x_{66}x_{179}+2x_{178}+2x_{179}-x_{219}\end{aligned}$\\[3pt]
$R_{188}$ & $\begin{aligned}[t]&1+x_{0}+x_{0}x_{66}+x_{66}+x_{66}x_{194}-x_{67}+x_{191}+x_{194}-x_{202}-x_{219}x_{220}\end{aligned}$\\[3pt]
$R_{189}$ & $\begin{aligned}[t]&-1-x_{61}^{2}+2x_{61}x_{178}+2x_{61}x_{179}-x_{221}\end{aligned}$\\[3pt]
$R_{190}$ & $\begin{aligned}[t]&x_{0}x_{61}+x_{61}+x_{61}x_{194}-x_{68}+x_{191}-x_{202}-x_{221}x_{222}\end{aligned}$\\[3pt]
$R_{191}$ & $\begin{aligned}[t]&-2x_{0}x_{178}-2x_{0}x_{179}-2x_{178}-2x_{178}x_{194}-2x_{179}-2x_{179}x_{194}+x_{223}\end{aligned}$\\[3pt]
$R_{192}$ & $\begin{aligned}[t]&-2x_{178}x_{191}-2x_{179}x_{191}+x_{224}\end{aligned}$\\[3pt]
$R_{193}$ & $\begin{aligned}[t]&-1-x_{0}-2x_{0}x_{211}-x_{194}-2x_{194}x_{211}-2x_{211}+x_{225}\end{aligned}$\\[3pt]
$R_{194}$ & $\begin{aligned}[t]&-10-6x_{0}-x_{0}^{2}+2x_{0}x_{178}+2x_{0}x_{179}+6x_{178}+6x_{179}-x_{226}\end{aligned}$\\[3pt]
$R_{195}$ & $\begin{aligned}[t]&x_{0}x_{224}+x_{0}x_{225}-x_{70}-x_{178}x_{224}-x_{178}x_{225}-x_{179}x_{224}-x_{179}x_{225}+x_{191}+2x_{191}x_{211}+x_{211}x_{223}+3x_{224}\\ &+3x_{225}-x_{226}x_{227}\end{aligned}$\\[3pt]
$R_{196}$ & $\begin{aligned}[t]&-2-2x_{70}-x_{70}^{2}+2x_{70}x_{178}+2x_{70}x_{179}+2x_{178}+2x_{179}-x_{228}\end{aligned}$\\[3pt]
$R_{197}$ & $\begin{aligned}[t]&1+x_{0}+x_{0}x_{70}+x_{70}+x_{70}x_{194}-x_{71}+x_{191}+x_{194}-x_{202}-x_{228}x_{229}\end{aligned}$\\[3pt]
$R_{198}$ & $\begin{aligned}[t]&-5-4x_{64}-x_{64}^{2}-2x_{64}x_{65}+2x_{64}x_{178}+2x_{64}x_{179}-2x_{64}x_{263}-4x_{65}-x_{65}^{2}+2x_{65}x_{178}+2x_{65}x_{179}\\ &-2x_{65}x_{263}+4x_{178}+2x_{178}x_{263}+4x_{179}+2x_{179}x_{263}-x_{230}-4x_{263}-x_{263}^{2}\end{aligned}$\\[3pt]
$R_{199}$ & $\begin{aligned}[t]&-1+2x_{60}+x_{60}x_{64}+x_{60}x_{65}-x_{60}x_{178}-x_{60}x_{179}+x_{60}x_{263}+x_{64}x_{206}+x_{65}x_{206}-x_{178}x_{206}-x_{179}x_{206}\\ &+x_{203}+2x_{206}+x_{206}x_{263}-x_{230}x_{231}-x_{261}-x_{262}\end{aligned}$\\[3pt]
$R_{200}$ & $\begin{aligned}[t]&1+2x_{0}+x_{0}x_{64}+x_{0}x_{65}+x_{0}x_{263}+x_{64}+x_{64}x_{194}+x_{65}+x_{65}x_{194}+x_{191}+2x_{194}+x_{194}x_{263}-x_{202}\\ &-x_{230}x_{232}+x_{263}-x_{273}-x_{274}\end{aligned}$\\[3pt]
$R_{201}$ & $\begin{aligned}[t]&-1+x_{0}x_{211}+x_{64}x_{191}+x_{64}x_{202}+x_{65}x_{191}+x_{65}x_{202}-x_{178}x_{202}-x_{179}x_{202}+2x_{191}+x_{191}x_{263}+x_{194}x_{211}\\ &+2x_{202}+x_{202}x_{263}+x_{211}-x_{230}x_{233}-x_{298}-x_{299}\end{aligned}$\\[3pt]
$R_{202}$ & $\begin{aligned}[t]&-9-6x_{2}-x_{2}^{2}-2x_{2}x_{26}-2x_{2}x_{237}-6x_{26}-x_{26}^{2}-2x_{26}x_{237}-6x_{237}-x_{237}^{2}+x_{238}\end{aligned}$\\[3pt]
$R_{203}$ & $\begin{aligned}[t]&-1-x_{211}x_{238}+x_{234}^{2}\end{aligned}$\\[3pt]
$R_{204}$ & $\begin{aligned}[t]&-x_{236}^{2}+x_{239}\end{aligned}$\\[3pt]
$R_{205}$ & $\begin{aligned}[t]&-1-x_{197}x_{239}+x_{235}^{2}+x_{239}\end{aligned}$\\[3pt]
$R_{206}$ & $\begin{aligned}[t]&-x_{185}x_{240}-x_{234}+x_{235}\end{aligned}$\\[3pt]
$R_{207}$ & $\begin{aligned}[t]&-3-x_{2}-4x_{2}x_{241}-x_{26}-4x_{26}x_{241}+x_{236}-4x_{237}x_{241}-12x_{241}\end{aligned}$\\[3pt]
$R_{208}$ & $\begin{aligned}[t]&5+2x_{2}+x_{2}x_{64}+x_{2}x_{65}-x_{2}x_{178}-x_{2}x_{179}+x_{2}x_{263}+2x_{26}+x_{26}x_{64}+x_{26}x_{65}-x_{26}x_{178}-x_{26}x_{179}\\ &+x_{26}x_{263}+3x_{64}+x_{64}x_{237}+3x_{65}+x_{65}x_{237}-3x_{178}-x_{178}x_{237}-3x_{179}-x_{179}x_{237}-x_{230}x_{242}+x_{234}+2x_{237}\\ &+x_{237}x_{263}+3x_{263}-x_{343}-x_{344}\end{aligned}$\\[3pt]
$R_{209}$ & $\begin{aligned}[t]&-1-2x_{1}-x_{1}^{2}-2x_{1}x_{246}-2x_{246}-x_{246}^{2}+x_{247}\end{aligned}$\\[3pt]
$R_{210}$ & $\begin{aligned}[t]&-1-x_{211}x_{247}+x_{243}^{2}\end{aligned}$\\[3pt]
$R_{211}$ & $\begin{aligned}[t]&-x_{245}^{2}+x_{248}\end{aligned}$\\[3pt]
$R_{212}$ & $\begin{aligned}[t]&-1-x_{197}x_{248}+x_{244}^{2}+x_{248}\end{aligned}$\\[3pt]
$R_{213}$ & $\begin{aligned}[t]&-x_{185}x_{249}-x_{243}+x_{244}\end{aligned}$\\[3pt]
$R_{214}$ & $\begin{aligned}[t]&-1-x_{1}-4x_{1}x_{250}+x_{245}-4x_{246}x_{250}-4x_{250}\end{aligned}$\\[3pt]
$R_{215}$ & $\begin{aligned}[t]&1+2x_{1}+x_{1}x_{64}+x_{1}x_{65}-x_{1}x_{178}-x_{1}x_{179}+x_{1}x_{263}+x_{64}+x_{64}x_{246}+x_{65}+x_{65}x_{246}-x_{178}-x_{178}x_{246}\\ &-x_{179}-x_{179}x_{246}-x_{230}x_{251}+x_{243}+2x_{246}+x_{246}x_{263}+x_{263}-x_{366}-x_{367}\end{aligned}$\\[3pt]
$R_{216}$ & $\begin{aligned}[t]&-x_{60}^{2}-2x_{60}x_{62}-2x_{60}x_{255}-x_{62}^{2}-2x_{62}x_{255}-x_{255}^{2}+x_{256}\end{aligned}$\\[3pt]
$R_{217}$ & $\begin{aligned}[t]&-1-x_{211}x_{256}+x_{252}^{2}\end{aligned}$\\[3pt]
$R_{218}$ & $\begin{aligned}[t]&-x_{254}^{2}+x_{257}\end{aligned}$\\[3pt]
$R_{219}$ & $\begin{aligned}[t]&-1-x_{197}x_{257}+x_{253}^{2}+x_{257}\end{aligned}$\\[3pt]
$R_{220}$ & $\begin{aligned}[t]&-x_{185}x_{258}-x_{252}+x_{253}\end{aligned}$\\[3pt]
$R_{221}$ & $\begin{aligned}[t]&-x_{60}-4x_{60}x_{259}-x_{62}-4x_{62}x_{259}+x_{254}-4x_{255}x_{259}\end{aligned}$\\[3pt]
$R_{222}$ & $\begin{aligned}[t]&-1+3x_{60}+x_{60}x_{64}+x_{60}x_{65}-x_{60}x_{178}-x_{60}x_{179}+x_{60}x_{263}+3x_{62}+x_{62}x_{64}+x_{62}x_{65}-x_{62}x_{178}-x_{62}x_{179}\\ &+x_{62}x_{263}-x_{64}+x_{64}x_{255}-x_{64}x_{261}-x_{64}x_{262}-x_{65}+x_{65}x_{255}-x_{65}x_{261}-x_{65}x_{262}-x_{176}x_{260}-x_{178}x_{255}\\ &-x_{179}x_{255}+x_{252}+3x_{255}+x_{255}x_{263}-2x_{261}-x_{262}\end{aligned}$\\[3pt]
$R_{223}$ & $\begin{aligned}[t]&-1-2x_{0}-x_{0}^{2}-2x_{0}x_{70}-2x_{0}x_{267}-2x_{70}-x_{70}^{2}-2x_{70}x_{267}-2x_{267}-x_{267}^{2}+x_{268}\end{aligned}$\\[3pt]
$R_{224}$ & $\begin{aligned}[t]&-1-x_{211}x_{268}+x_{264}^{2}\end{aligned}$\\[3pt]
$R_{225}$ & $\begin{aligned}[t]&-x_{266}^{2}+x_{269}\end{aligned}$\\[3pt]
$R_{226}$ & $\begin{aligned}[t]&-1-x_{197}x_{269}+x_{265}^{2}+x_{269}\end{aligned}$\\[3pt]
$R_{227}$ & $\begin{aligned}[t]&-x_{185}x_{270}-x_{264}+x_{265}\end{aligned}$\\[3pt]
$R_{228}$ & $\begin{aligned}[t]&-1-x_{0}-4x_{0}x_{271}-x_{70}-4x_{70}x_{271}+x_{266}-4x_{267}x_{271}-4x_{271}\end{aligned}$\\[3pt]
$R_{229}$ & $\begin{aligned}[t]&-2-x_{64}-x_{64}x_{273}-x_{64}x_{274}-x_{65}-x_{65}x_{273}-x_{65}x_{274}-x_{263}-x_{263}x_{273}-x_{263}x_{274}-2x_{273}-2x_{274}+x_{275}\end{aligned}$\\[3pt]
$R_{230}$ & $\begin{aligned}[t]&x_{64}+x_{65}-x_{72}-x_{73}+x_{263}-x_{276}\end{aligned}$\\[3pt]
$R_{231}$ & $\begin{aligned}[t]&2+3x_{0}+x_{0}x_{64}+x_{0}x_{65}-x_{0}x_{178}-x_{0}x_{179}+x_{0}x_{263}+x_{64}+x_{64}x_{70}+x_{64}x_{267}+x_{65}+x_{65}x_{70}+x_{65}x_{267}\\ &+3x_{70}-x_{70}x_{178}-x_{70}x_{179}+x_{70}x_{263}-x_{72}-x_{72}x_{273}-x_{72}x_{274}-x_{73}-x_{73}x_{273}-x_{73}x_{274}-x_{178}-x_{178}x_{267}\\ &-x_{179}-x_{179}x_{267}+x_{263}+x_{263}x_{267}+x_{264}+3x_{267}-x_{272}x_{275}-2x_{273}-x_{274}\end{aligned}$\\[3pt]
$R_{232}$ & $\begin{aligned}[t]&-9-6x_{0}-x_{0}^{2}-2x_{0}x_{67}-2x_{0}x_{280}-6x_{67}-x_{67}^{2}-2x_{67}x_{280}-6x_{280}-x_{280}^{2}+x_{281}\end{aligned}$\\[3pt]
$R_{233}$ & $\begin{aligned}[t]&-1-x_{211}x_{281}+x_{277}^{2}\end{aligned}$\\[3pt]
$R_{234}$ & $\begin{aligned}[t]&-x_{279}^{2}+x_{282}\end{aligned}$\\[3pt]
$R_{235}$ & $\begin{aligned}[t]&-1-x_{197}x_{282}+x_{278}^{2}+x_{282}\end{aligned}$\\[3pt]
$R_{236}$ & $\begin{aligned}[t]&-x_{185}x_{283}-x_{277}+x_{278}\end{aligned}$\\[3pt]
$R_{237}$ & $\begin{aligned}[t]&-3-x_{0}-4x_{0}x_{284}-x_{67}-4x_{67}x_{284}+x_{279}-4x_{280}x_{284}-12x_{284}\end{aligned}$\\[3pt]
$R_{238}$ & $\begin{aligned}[t]&1+x_{64}+x_{65}-x_{74}+x_{263}-x_{288}\end{aligned}$\\[3pt]
$R_{239}$ & $\begin{aligned}[t]&9+3x_{0}+x_{0}x_{64}+x_{0}x_{65}-x_{0}x_{178}-x_{0}x_{179}+x_{0}x_{263}+3x_{64}+x_{64}x_{67}+x_{64}x_{280}+3x_{65}+x_{65}x_{67}+x_{65}x_{280}\\ &+3x_{67}-x_{67}x_{178}-x_{67}x_{179}+x_{67}x_{263}-x_{74}-x_{74}x_{273}-x_{74}x_{274}-3x_{178}-x_{178}x_{280}-3x_{179}-x_{179}x_{280}\\ &+3x_{263}+x_{263}x_{280}-x_{275}x_{285}+x_{277}+3x_{280}-x_{286}\end{aligned}$\\[3pt]
$R_{240}$ & $\begin{aligned}[t]&x_{273}+x_{274}-x_{286}-x_{287}\end{aligned}$\\[3pt]
$R_{241}$ & $\begin{aligned}[t]&-1-2x_{29}-x_{29}^{2}-2x_{29}x_{292}-2x_{292}-x_{292}^{2}+x_{293}\end{aligned}$\\[3pt]
$R_{242}$ & $\begin{aligned}[t]&-1-x_{211}x_{293}+x_{289}^{2}\end{aligned}$\\[3pt]
$R_{243}$ & $\begin{aligned}[t]&-x_{291}^{2}+x_{294}\end{aligned}$\\[3pt]
$R_{244}$ & $\begin{aligned}[t]&-1-x_{197}x_{294}+x_{290}^{2}+x_{294}\end{aligned}$\\[3pt]
$R_{245}$ & $\begin{aligned}[t]&-x_{185}x_{295}-x_{289}+x_{290}\end{aligned}$\\[3pt]
$R_{246}$ & $\begin{aligned}[t]&-1-x_{29}-4x_{29}x_{296}+x_{291}-4x_{292}x_{296}-4x_{296}\end{aligned}$\\[3pt]
$R_{247}$ & $\begin{aligned}[t]&-x_{76}-x_{76}x_{298}-x_{76}x_{299}+x_{300}\end{aligned}$\\[3pt]
$R_{248}$ & $\begin{aligned}[t]&3+3x_{29}+x_{29}x_{64}+x_{29}x_{65}-x_{29}x_{178}-x_{29}x_{179}+x_{29}x_{263}+x_{64}+x_{64}x_{292}+x_{65}+x_{65}x_{292}-x_{178}-x_{178}x_{292}\\ &-x_{179}-x_{179}x_{292}+x_{263}+x_{263}x_{292}+x_{289}+3x_{292}-x_{297}x_{300}-x_{298}\end{aligned}$\\[3pt]
$R_{249}$ & $\begin{aligned}[t]&-1-2x_{30}-x_{30}^{2}-2x_{30}x_{304}-2x_{304}-x_{304}^{2}+x_{305}\end{aligned}$\\[3pt]
$R_{250}$ & $\begin{aligned}[t]&-1-x_{211}x_{305}+x_{301}^{2}\end{aligned}$\\[3pt]
$R_{251}$ & $\begin{aligned}[t]&-x_{303}^{2}+x_{306}\end{aligned}$\\[3pt]
$R_{252}$ & $\begin{aligned}[t]&-1-x_{197}x_{306}+x_{302}^{2}+x_{306}\end{aligned}$\\[3pt]
$R_{253}$ & $\begin{aligned}[t]&-x_{185}x_{307}-x_{301}+x_{302}\end{aligned}$\\[3pt]
$R_{254}$ & $\begin{aligned}[t]&-1-x_{30}-4x_{30}x_{308}+x_{303}-4x_{304}x_{308}-4x_{308}\end{aligned}$\\[3pt]
$R_{255}$ & $\begin{aligned}[t]&3+3x_{30}+x_{30}x_{64}+x_{30}x_{65}-x_{30}x_{178}-x_{30}x_{179}+x_{30}x_{263}+x_{64}+x_{64}x_{304}+x_{65}+x_{65}x_{304}-x_{178}-x_{178}x_{304}\\ &-x_{179}-x_{179}x_{304}+x_{263}+x_{263}x_{304}-x_{300}x_{309}+x_{301}+3x_{304}-x_{310}\end{aligned}$\\[3pt]
$R_{256}$ & $\begin{aligned}[t]&x_{298}+x_{299}-x_{310}-x_{311}\end{aligned}$\\[3pt]
$R_{257}$ & $\begin{aligned}[t]&-1-2x_{31}-x_{31}^{2}-2x_{31}x_{315}-2x_{315}-x_{315}^{2}+x_{316}\end{aligned}$\\[3pt]
$R_{258}$ & $\begin{aligned}[t]&-1-x_{211}x_{316}+x_{312}^{2}\end{aligned}$\\[3pt]
$R_{259}$ & $\begin{aligned}[t]&-x_{314}^{2}+x_{317}\end{aligned}$\\[3pt]
$R_{260}$ & $\begin{aligned}[t]&-1-x_{197}x_{317}+x_{313}^{2}+x_{317}\end{aligned}$\\[3pt]
$R_{261}$ & $\begin{aligned}[t]&-x_{185}x_{318}-x_{312}+x_{313}\end{aligned}$\\[3pt]
$R_{262}$ & $\begin{aligned}[t]&-1-x_{31}-4x_{31}x_{319}+x_{314}-4x_{315}x_{319}-4x_{319}\end{aligned}$\\[3pt]
$R_{263}$ & $\begin{aligned}[t]&3+3x_{31}+x_{31}x_{64}+x_{31}x_{65}-x_{31}x_{178}-x_{31}x_{179}+x_{31}x_{263}+x_{64}+x_{64}x_{315}+x_{65}+x_{65}x_{315}-x_{178}-x_{178}x_{315}\\ &-x_{179}-x_{179}x_{315}+x_{263}+x_{263}x_{315}-x_{300}x_{320}+x_{312}+3x_{315}-x_{321}\end{aligned}$\\[3pt]
$R_{264}$ & $\begin{aligned}[t]&x_{298}+x_{299}-x_{321}-x_{322}\end{aligned}$\\[3pt]
$R_{265}$ & $\begin{aligned}[t]&-1-2x_{32}-x_{32}^{2}-2x_{32}x_{326}-2x_{326}-x_{326}^{2}+x_{327}\end{aligned}$\\[3pt]
$R_{266}$ & $\begin{aligned}[t]&-1-x_{211}x_{327}+x_{323}^{2}\end{aligned}$\\[3pt]
$R_{267}$ & $\begin{aligned}[t]&-x_{325}^{2}+x_{328}\end{aligned}$\\[3pt]
$R_{268}$ & $\begin{aligned}[t]&-1-x_{197}x_{328}+x_{324}^{2}+x_{328}\end{aligned}$\\[3pt]
$R_{269}$ & $\begin{aligned}[t]&-x_{185}x_{329}-x_{323}+x_{324}\end{aligned}$\\[3pt]
$R_{270}$ & $\begin{aligned}[t]&-1-x_{32}-4x_{32}x_{330}+x_{325}-4x_{326}x_{330}-4x_{330}\end{aligned}$\\[3pt]
$R_{271}$ & $\begin{aligned}[t]&3+3x_{32}+x_{32}x_{64}+x_{32}x_{65}-x_{32}x_{178}-x_{32}x_{179}+x_{32}x_{263}+x_{64}+x_{64}x_{326}+x_{65}+x_{65}x_{326}-x_{178}-x_{178}x_{326}\\ &-x_{179}-x_{179}x_{326}+x_{263}+x_{263}x_{326}-x_{300}x_{331}+x_{323}+3x_{326}-x_{332}\end{aligned}$\\[3pt]
$R_{272}$ & $\begin{aligned}[t]&x_{298}+x_{299}-x_{332}-x_{333}\end{aligned}$\\[3pt]
$R_{273}$ & $\begin{aligned}[t]&-1-2x_{33}-x_{33}^{2}-2x_{33}x_{337}-2x_{337}-x_{337}^{2}+x_{338}\end{aligned}$\\[3pt]
$R_{274}$ & $\begin{aligned}[t]&-1-x_{211}x_{338}+x_{334}^{2}\end{aligned}$\\[3pt]
$R_{275}$ & $\begin{aligned}[t]&-x_{336}^{2}+x_{339}\end{aligned}$\\[3pt]
$R_{276}$ & $\begin{aligned}[t]&-1-x_{197}x_{339}+x_{335}^{2}+x_{339}\end{aligned}$\\[3pt]
$R_{277}$ & $\begin{aligned}[t]&-x_{185}x_{340}-x_{334}+x_{335}\end{aligned}$\\[3pt]
$R_{278}$ & $\begin{aligned}[t]&-1-x_{33}-4x_{33}x_{341}+x_{336}-4x_{337}x_{341}-4x_{341}\end{aligned}$\\[3pt]
$R_{279}$ & $\begin{aligned}[t]&-x_{76}-x_{76}x_{343}-x_{76}x_{344}+x_{345}\end{aligned}$\\[3pt]
$R_{280}$ & $\begin{aligned}[t]&3+3x_{33}+x_{33}x_{64}+x_{33}x_{65}-x_{33}x_{178}-x_{33}x_{179}+x_{33}x_{263}+x_{64}+x_{64}x_{337}+x_{65}+x_{65}x_{337}-x_{178}-x_{178}x_{337}\\ &-x_{179}-x_{179}x_{337}+x_{263}+x_{263}x_{337}+x_{334}+3x_{337}-x_{342}x_{345}-x_{343}\end{aligned}$\\[3pt]
$R_{281}$ & $\begin{aligned}[t]&-1-2x_{34}-x_{34}^{2}-2x_{34}x_{349}-2x_{349}-x_{349}^{2}+x_{350}\end{aligned}$\\[3pt]
$R_{282}$ & $\begin{aligned}[t]&-1-x_{211}x_{350}+x_{346}^{2}\end{aligned}$\\[3pt]
$R_{283}$ & $\begin{aligned}[t]&-x_{348}^{2}+x_{351}\end{aligned}$\\[3pt]
$R_{284}$ & $\begin{aligned}[t]&-1-x_{197}x_{351}+x_{347}^{2}+x_{351}\end{aligned}$\\[3pt]
$R_{285}$ & $\begin{aligned}[t]&-x_{185}x_{352}-x_{346}+x_{347}\end{aligned}$\\[3pt]
$R_{286}$ & $\begin{aligned}[t]&-1-x_{34}-4x_{34}x_{353}+x_{348}-4x_{349}x_{353}-4x_{353}\end{aligned}$\\[3pt]
$R_{287}$ & $\begin{aligned}[t]&3+3x_{34}+x_{34}x_{64}+x_{34}x_{65}-x_{34}x_{178}-x_{34}x_{179}+x_{34}x_{263}+x_{64}+x_{64}x_{349}+x_{65}+x_{65}x_{349}-x_{178}-x_{178}x_{349}\\ &-x_{179}-x_{179}x_{349}+x_{263}+x_{263}x_{349}-x_{345}x_{354}+x_{346}+3x_{349}-x_{355}\end{aligned}$\\[3pt]
$R_{288}$ & $\begin{aligned}[t]&x_{343}+x_{344}-x_{355}-x_{356}\end{aligned}$\\[3pt]
$R_{289}$ & $\begin{aligned}[t]&-1-2x_{35}-x_{35}^{2}-2x_{35}x_{360}-2x_{360}-x_{360}^{2}+x_{361}\end{aligned}$\\[3pt]
$R_{290}$ & $\begin{aligned}[t]&-1-x_{211}x_{361}+x_{357}^{2}\end{aligned}$\\[3pt]
$R_{291}$ & $\begin{aligned}[t]&-x_{359}^{2}+x_{362}\end{aligned}$\\[3pt]
$R_{292}$ & $\begin{aligned}[t]&-1-x_{197}x_{362}+x_{358}^{2}+x_{362}\end{aligned}$\\[3pt]
$R_{293}$ & $\begin{aligned}[t]&-x_{185}x_{363}-x_{357}+x_{358}\end{aligned}$\\[3pt]
$R_{294}$ & $\begin{aligned}[t]&-1-x_{35}-4x_{35}x_{364}+x_{359}-4x_{360}x_{364}-4x_{364}\end{aligned}$\\[3pt]
$R_{295}$ & $\begin{aligned}[t]&-x_{76}-x_{76}x_{366}-x_{76}x_{367}+x_{368}\end{aligned}$\\[3pt]
$R_{296}$ & $\begin{aligned}[t]&3+3x_{35}+x_{35}x_{64}+x_{35}x_{65}-x_{35}x_{178}-x_{35}x_{179}+x_{35}x_{263}+x_{64}+x_{64}x_{360}+x_{65}+x_{65}x_{360}-x_{178}-x_{178}x_{360}\\ &-x_{179}-x_{179}x_{360}+x_{263}+x_{263}x_{360}+x_{357}+3x_{360}-x_{365}x_{368}-x_{366}\end{aligned}$\\[3pt]
$R_{297}$ & $\begin{aligned}[t]&-1-2x_{36}-x_{36}^{2}-2x_{36}x_{372}-2x_{372}-x_{372}^{2}+x_{373}\end{aligned}$\\[3pt]
$R_{298}$ & $\begin{aligned}[t]&-1-x_{211}x_{373}+x_{369}^{2}\end{aligned}$\\[3pt]
$R_{299}$ & $\begin{aligned}[t]&-x_{371}^{2}+x_{374}\end{aligned}$\\[3pt]
$R_{300}$ & $\begin{aligned}[t]&-1-x_{197}x_{374}+x_{370}^{2}+x_{374}\end{aligned}$\\[3pt]
$R_{301}$ & $\begin{aligned}[t]&-x_{185}x_{375}-x_{369}+x_{370}\end{aligned}$\\[3pt]
$R_{302}$ & $\begin{aligned}[t]&-1-x_{36}-4x_{36}x_{376}+x_{371}-4x_{372}x_{376}-4x_{376}\end{aligned}$\\[3pt]
$R_{303}$ & $\begin{aligned}[t]&3+3x_{36}+x_{36}x_{64}+x_{36}x_{65}-x_{36}x_{178}-x_{36}x_{179}+x_{36}x_{263}+x_{64}+x_{64}x_{372}+x_{65}+x_{65}x_{372}-x_{178}-x_{178}x_{372}\\ &-x_{179}-x_{179}x_{372}+x_{263}+x_{263}x_{372}-x_{368}x_{377}+x_{369}+3x_{372}-x_{378}\end{aligned}$\\[3pt]
$R_{304}$ & $\begin{aligned}[t]&x_{366}+x_{367}-x_{378}-x_{379}\end{aligned}$\\[3pt]
$R_{305}$ & $\begin{aligned}[t]&-x_{195}x_{207}+x_{380}\end{aligned}$\\[3pt]
$R_{306}$ & $\begin{aligned}[t]&-x_{238}x_{380}+x_{381}\end{aligned}$\\[3pt]
$R_{307}$ & $\begin{aligned}[t]&-x_{247}x_{381}+x_{382}\end{aligned}$\\[3pt]
$R_{308}$ & $\begin{aligned}[t]&-x_{256}x_{382}+x_{383}\end{aligned}$\\[3pt]
$R_{309}$ & $\begin{aligned}[t]&-x_{268}x_{383}+x_{384}\end{aligned}$\\[3pt]
$R_{310}$ & $\begin{aligned}[t]&-x_{281}x_{384}+x_{385}\end{aligned}$\\[3pt]
$R_{311}$ & $\begin{aligned}[t]&-x_{293}x_{385}+x_{386}\end{aligned}$\\[3pt]
$R_{312}$ & $\begin{aligned}[t]&-x_{305}x_{386}+x_{387}\end{aligned}$\\[3pt]
$R_{313}$ & $\begin{aligned}[t]&-x_{316}x_{387}+x_{388}\end{aligned}$\\[3pt]
$R_{314}$ & $\begin{aligned}[t]&-x_{327}x_{388}+x_{389}\end{aligned}$\\[3pt]
$R_{315}$ & $\begin{aligned}[t]&-x_{338}x_{389}+x_{390}\end{aligned}$\\[3pt]
$R_{316}$ & $\begin{aligned}[t]&-x_{350}x_{390}+x_{391}\end{aligned}$\\[3pt]
$R_{317}$ & $\begin{aligned}[t]&-x_{361}x_{391}+x_{392}\end{aligned}$\\[3pt]
$R_{318}$ & $\begin{aligned}[t]&-x_{373}x_{392}+x_{393}\end{aligned}$\\[3pt]
$R_{319}$ & $\begin{aligned}[t]&x_{186}-x_{393}-x_{393}x_{394}\end{aligned}$\\[3pt]
$R_{320}$ & $\begin{aligned}[t]&-24-3x_{0}-x_{1}-x_{2}-x_{3}-x_{26}-2x_{60}-x_{61}-2x_{62}-x_{63}-x_{64}-x_{65}-x_{66}-x_{67}-x_{68}-2x_{70}-x_{71}+x_{177}\\ &-x_{261}-x_{262}-x_{263}-x_{273}-x_{274}-x_{298}-x_{299}-x_{343}-x_{344}-x_{366}-x_{367}\end{aligned}$\\[3pt]
\end{longtable}
\endgroup


\clearpage
\section{The complete \zvars-variable integer equation}\label{app:integer}

Here all coordinates range independently over $\Z$. There are no
additional inequalities or omitted domain predicates:
\[
 \boxed{\PZ(x_0,\ldots,x_{\zvarlast})=
       \sum_{j=0}^{\zlast}S_j(x_0,\ldots,x_{\zvarlast})^2=0.}
\]
The sign conditions required by the proof have already been compiled
into these residuals. The variable indices in this appendix refer only
to this equation.
\begingroup\footnotesize
\setlength{\tabcolsep}{3pt}
\begin{longtable}{r@{\hspace{0.8em}}l}
\toprule
Residual & Polynomial (set each to zero)\\
\midrule
\endfirsthead
\toprule
Residual & Polynomial (continued)\\
\midrule
\endhead
\bottomrule\endfoot
$S_{0}$ & $\begin{aligned}[t]&-x_{4}^{2}+x_{5}\end{aligned}$\\[3pt]
$S_{1}$ & $\begin{aligned}[t]&-x_{5}^{2}+x_{6}\end{aligned}$\\[3pt]
$S_{2}$ & $\begin{aligned}[t]&-x_{6}^{2}+x_{7}\end{aligned}$\\[3pt]
$S_{3}$ & $\begin{aligned}[t]&-x_{7}^{2}+x_{8}\end{aligned}$\\[3pt]
$S_{4}$ & $\begin{aligned}[t]&-x_{8}^{2}+x_{9}\end{aligned}$\\[3pt]
$S_{5}$ & $\begin{aligned}[t]&-x_{9}^{2}+x_{10}\end{aligned}$\\[3pt]
$S_{6}$ & $\begin{aligned}[t]&-x_{10}^{2}+x_{11}\end{aligned}$\\[3pt]
$S_{7}$ & $\begin{aligned}[t]&-x_{11}^{2}+x_{12}\end{aligned}$\\[3pt]
$S_{8}$ & $\begin{aligned}[t]&-x_{12}^{2}+x_{13}\end{aligned}$\\[3pt]
$S_{9}$ & $\begin{aligned}[t]&-x_{13}^{2}+x_{14}\end{aligned}$\\[3pt]
$S_{10}$ & $\begin{aligned}[t]&-x_{14}^{2}+x_{15}\end{aligned}$\\[3pt]
$S_{11}$ & $\begin{aligned}[t]&-x_{15}^{2}+x_{16}\end{aligned}$\\[3pt]
$S_{12}$ & $\begin{aligned}[t]&-x_{16}^{2}+x_{17}\end{aligned}$\\[3pt]
$S_{13}$ & $\begin{aligned}[t]&-x_{17}^{2}+x_{18}\end{aligned}$\\[3pt]
$S_{14}$ & $\begin{aligned}[t]&-x_{18}^{2}+x_{19}\end{aligned}$\\[3pt]
$S_{15}$ & $\begin{aligned}[t]&-x_{19}^{2}+x_{20}\end{aligned}$\\[3pt]
$S_{16}$ & $\begin{aligned}[t]&-x_{20}^{2}+x_{21}\end{aligned}$\\[3pt]
$S_{17}$ & $\begin{aligned}[t]&-x_{21}^{2}+x_{22}\end{aligned}$\\[3pt]
$S_{18}$ & $\begin{aligned}[t]&-x_{22}^{2}+x_{23}\end{aligned}$\\[3pt]
$S_{19}$ & $\begin{aligned}[t]&-x_{23}^{2}+x_{24}\end{aligned}$\\[3pt]
$S_{20}$ & $\begin{aligned}[t]&-x_{24}^{2}+x_{25}\end{aligned}$\\[3pt]
$S_{21}$ & $\begin{aligned}[t]&1+x_{3}-x_{25}^{2}\end{aligned}$\\[3pt]
$S_{22}$ & $\begin{aligned}[t]&-x_{0}x_{2}-2x_{2}+x_{28}\end{aligned}$\\[3pt]
$S_{23}$ & $\begin{aligned}[t]&-2+x_{1}-x_{2}x_{26}-x_{3}-x_{26}\end{aligned}$\\[3pt]
$S_{24}$ & $\begin{aligned}[t]&-1+x_{1}-x_{27}-x_{27}x_{28}\end{aligned}$\\[3pt]
$S_{25}$ & $\begin{aligned}[t]&-9-6x_{0}-x_{0}^{2}-2x_{0}x_{1}-2x_{0}x_{2}-2x_{0}x_{3}-6x_{1}-x_{1}^{2}-2x_{1}x_{2}-2x_{1}x_{3}-6x_{2}-x_{2}^{2}-2x_{2}x_{3}-6x_{3}-x_{3}^{2}\\ &+x_{38}\end{aligned}$\\[3pt]
$S_{26}$ & $\begin{aligned}[t]&-x_{38}^{2}+x_{39}\end{aligned}$\\[3pt]
$S_{27}$ & $\begin{aligned}[t]&-x_{39}^{2}+x_{40}\end{aligned}$\\[3pt]
$S_{28}$ & $\begin{aligned}[t]&-x_{40}^{2}+x_{41}\end{aligned}$\\[3pt]
$S_{29}$ & $\begin{aligned}[t]&-x_{41}^{2}+x_{42}\end{aligned}$\\[3pt]
$S_{30}$ & $\begin{aligned}[t]&-x_{42}^{2}+x_{43}\end{aligned}$\\[3pt]
$S_{31}$ & $\begin{aligned}[t]&-x_{43}^{2}+x_{44}\end{aligned}$\\[3pt]
$S_{32}$ & $\begin{aligned}[t]&-x_{44}^{2}+x_{45}\end{aligned}$\\[3pt]
$S_{33}$ & $\begin{aligned}[t]&-x_{45}^{2}+x_{46}\end{aligned}$\\[3pt]
$S_{34}$ & $\begin{aligned}[t]&-x_{46}^{2}+x_{47}\end{aligned}$\\[3pt]
$S_{35}$ & $\begin{aligned}[t]&-x_{47}^{2}+x_{48}\end{aligned}$\\[3pt]
$S_{36}$ & $\begin{aligned}[t]&-x_{48}^{2}+x_{49}\end{aligned}$\\[3pt]
$S_{37}$ & $\begin{aligned}[t]&-x_{49}^{2}+x_{50}\end{aligned}$\\[3pt]
$S_{38}$ & $\begin{aligned}[t]&-x_{50}^{2}+x_{51}\end{aligned}$\\[3pt]
$S_{39}$ & $\begin{aligned}[t]&-x_{51}^{2}+x_{52}\end{aligned}$\\[3pt]
$S_{40}$ & $\begin{aligned}[t]&-x_{52}^{2}+x_{53}\end{aligned}$\\[3pt]
$S_{41}$ & $\begin{aligned}[t]&-x_{53}^{2}+x_{54}\end{aligned}$\\[3pt]
$S_{42}$ & $\begin{aligned}[t]&-x_{54}^{2}+x_{55}\end{aligned}$\\[3pt]
$S_{43}$ & $\begin{aligned}[t]&-x_{55}^{2}+x_{56}\end{aligned}$\\[3pt]
$S_{44}$ & $\begin{aligned}[t]&-x_{56}^{2}+x_{57}\end{aligned}$\\[3pt]
$S_{45}$ & $\begin{aligned}[t]&-x_{57}^{2}+x_{58}\end{aligned}$\\[3pt]
$S_{46}$ & $\begin{aligned}[t]&-x_{58}^{2}+x_{59}\end{aligned}$\\[3pt]
$S_{47}$ & $\begin{aligned}[t]&-x_{59}^{2}+x_{207}-x_{210}\end{aligned}$\\[3pt]
$S_{48}$ & $\begin{aligned}[t]&-x_{0}x_{60}-2x_{60}+x_{64}\end{aligned}$\\[3pt]
$S_{49}$ & $\begin{aligned}[t]&-x_{66}x_{67}+x_{71}\end{aligned}$\\[3pt]
$S_{50}$ & $\begin{aligned}[t]&-x_{70}x_{71}+x_{72}x_{74}+x_{73}\end{aligned}$\\[3pt]
$S_{51}$ & $\begin{aligned}[t]&1-2x_{60}-x_{60}x_{65}+x_{72}\end{aligned}$\\[3pt]
$S_{52}$ & $\begin{aligned}[t]&-1-x_{37}x_{60}-x_{60}+x_{73}\end{aligned}$\\[3pt]
$S_{53}$ & $\begin{aligned}[t]&-6-5x_{29}-x_{29}^{2}-2x_{29}x_{30}-5x_{30}-x_{30}^{2}+x_{75}\end{aligned}$\\[3pt]
$S_{54}$ & $\begin{aligned}[t]&-72-68x_{30}-16x_{30}^{2}-16x_{30}x_{75}-34x_{75}-4x_{75}^{2}+x_{76}\end{aligned}$\\[3pt]
$S_{55}$ & $\begin{aligned}[t]&-30-11x_{29}-x_{29}^{2}-4x_{29}x_{76}-22x_{76}-4x_{76}^{2}+x_{77}\end{aligned}$\\[3pt]
$S_{56}$ & $\begin{aligned}[t]&-380-312x_{76}-64x_{76}^{2}-32x_{76}x_{77}-78x_{77}-4x_{77}^{2}+x_{78}\end{aligned}$\\[3pt]
$S_{57}$ & $\begin{aligned}[t]&-6-5x_{30}-x_{30}^{2}-2x_{30}x_{31}-5x_{31}-x_{31}^{2}+x_{79}\end{aligned}$\\[3pt]
$S_{58}$ & $\begin{aligned}[t]&-30-11x_{29}-x_{29}^{2}-8x_{29}x_{31}-4x_{29}x_{79}-44x_{31}-16x_{31}^{2}-16x_{31}x_{79}-22x_{79}-4x_{79}^{2}+x_{80}\end{aligned}$\\[3pt]
$S_{59}$ & $\begin{aligned}[t]&-6-5x_{29}-x_{29}^{2}-2x_{29}x_{31}-5x_{31}-x_{31}^{2}+x_{81}\end{aligned}$\\[3pt]
$S_{60}$ & $\begin{aligned}[t]&-156-100x_{30}-16x_{30}^{2}-32x_{30}x_{31}-16x_{30}x_{75}-16x_{30}x_{81}-100x_{31}-16x_{31}^{2}-16x_{31}x_{75}-16x_{31}x_{81}\\ &-50x_{75}-4x_{75}^{2}-8x_{75}x_{81}-50x_{81}-4x_{81}^{2}+x_{82}\end{aligned}$\\[3pt]
$S_{61}$ & $\begin{aligned}[t]&-1806-2720x_{31}-1024x_{31}^{2}-512x_{31}x_{79}-128x_{31}x_{80}-512x_{31}x_{81}-128x_{31}x_{82}-680x_{79}-64x_{79}^{2}\\ &-32x_{79}x_{80}-128x_{79}x_{81}-32x_{79}x_{82}-170x_{80}-4x_{80}^{2}-32x_{80}x_{81}-8x_{80}x_{82}-680x_{81}-64x_{81}^{2}-32x_{81}x_{82}\\ &-170x_{82}-4x_{82}^{2}+x_{83}\end{aligned}$\\[3pt]
$S_{62}$ & $\begin{aligned}[t]&-110-84x_{29}-16x_{29}^{2}-32x_{29}x_{30}-16x_{29}x_{79}-16x_{29}x_{81}-84x_{30}-16x_{30}^{2}-16x_{30}x_{79}-16x_{30}x_{81}-42x_{79}\\ &-4x_{79}^{2}-8x_{79}x_{81}-42x_{81}-4x_{81}^{2}+x_{84}\end{aligned}$\\[3pt]
$S_{63}$ & $\begin{aligned}[t]&-756-1100x_{30}-400x_{30}^{2}-80x_{30}x_{75}-320x_{30}x_{79}-80x_{30}x_{84}-110x_{75}-4x_{75}^{2}-32x_{75}x_{79}-8x_{75}x_{84}\\ &-440x_{79}-64x_{79}^{2}-32x_{79}x_{84}-110x_{84}-4x_{84}^{2}+x_{85}\end{aligned}$\\[3pt]
$S_{64}$ & $\begin{aligned}[t]&-56-15x_{29}-x_{29}^{2}-8x_{29}x_{31}-4x_{29}x_{79}-60x_{31}-16x_{31}^{2}-16x_{31}x_{79}-30x_{79}-4x_{79}^{2}+x_{86}\end{aligned}$\\[3pt]
$S_{65}$ & $\begin{aligned}[t]&-56-15x_{30}-x_{30}^{2}-8x_{30}x_{31}-4x_{30}x_{81}-60x_{31}-16x_{31}^{2}-16x_{31}x_{81}-30x_{81}-4x_{81}^{2}+x_{87}\end{aligned}$\\[3pt]
$S_{66}$ & $\begin{aligned}[t]&-2756-3360x_{31}-1024x_{31}^{2}-512x_{31}x_{79}-512x_{31}x_{81}-128x_{31}x_{86}-128x_{31}x_{87}-840x_{79}-64x_{79}^{2}\\ &-128x_{79}x_{81}-32x_{79}x_{86}-32x_{79}x_{87}-840x_{81}-64x_{81}^{2}-32x_{81}x_{86}-32x_{81}x_{87}-210x_{86}-4x_{86}^{2}-8x_{86}x_{87}\\ &-210x_{87}-4x_{87}^{2}+x_{88}\end{aligned}$\\[3pt]
$S_{67}$ & $\begin{aligned}[t]&-110-168x_{30}-64x_{30}^{2}-32x_{30}x_{75}-32x_{30}x_{79}-42x_{75}-4x_{75}^{2}-8x_{75}x_{79}-42x_{79}-4x_{79}^{2}+x_{89}\end{aligned}$\\[3pt]
$S_{68}$ & $\begin{aligned}[t]&-756-220x_{29}-16x_{29}^{2}-128x_{29}x_{30}-64x_{29}x_{75}-16x_{29}x_{81}-16x_{29}x_{89}-880x_{30}-256x_{30}^{2}-256x_{30}x_{75}\\ &-64x_{30}x_{81}-64x_{30}x_{89}-440x_{75}-64x_{75}^{2}-32x_{75}x_{81}-32x_{75}x_{89}-110x_{81}-4x_{81}^{2}-8x_{81}x_{89}-110x_{89}\\ &-4x_{89}^{2}+x_{90}\end{aligned}$\\[3pt]
$S_{69}$ & $\begin{aligned}[t]&-56-60x_{30}-16x_{30}^{2}-16x_{30}x_{75}-30x_{75}-4x_{75}^{2}+x_{91}\end{aligned}$\\[3pt]
$S_{70}$ & $\begin{aligned}[t]&-30-11x_{29}-x_{29}^{2}-4x_{29}x_{91}-22x_{91}-4x_{91}^{2}+x_{92}\end{aligned}$\\[3pt]
$S_{71}$ & $\begin{aligned}[t]&-380-312x_{91}-64x_{91}^{2}-32x_{91}x_{92}-78x_{92}-4x_{92}^{2}+x_{93}\end{aligned}$\\[3pt]
$S_{72}$ & $\begin{aligned}[t]&-30-11x_{30}-x_{30}^{2}-4x_{30}x_{93}-22x_{93}-4x_{93}^{2}+x_{94}\end{aligned}$\\[3pt]
$S_{73}$ & $\begin{aligned}[t]&-156-200x_{30}-64x_{30}^{2}-64x_{30}x_{75}-100x_{75}-16x_{75}^{2}+x_{95}\end{aligned}$\\[3pt]
$S_{74}$ & $\begin{aligned}[t]&-1406-1500x_{30}-400x_{30}^{2}-320x_{30}x_{75}-80x_{30}x_{79}-80x_{30}x_{95}-600x_{75}-64x_{75}^{2}-32x_{75}x_{79}-32x_{75}x_{95}\\ &-150x_{79}-4x_{79}^{2}-8x_{79}x_{95}-150x_{95}-4x_{95}^{2}+x_{96}\end{aligned}$\\[3pt]
$S_{75}$ & $\begin{aligned}[t]&-16256-1020x_{29}-16x_{29}^{2}-512x_{29}x_{30}-256x_{29}x_{75}-16x_{29}x_{81}-64x_{29}x_{95}-16x_{29}x_{96}-16320x_{30}\\ &-4096x_{30}^{2}-4096x_{30}x_{75}-256x_{30}x_{81}-1024x_{30}x_{95}-256x_{30}x_{96}-8160x_{75}-1024x_{75}^{2}-128x_{75}x_{81}\\ &-512x_{75}x_{95}-128x_{75}x_{96}-510x_{81}-4x_{81}^{2}-32x_{81}x_{95}-8x_{81}x_{96}-2040x_{95}-64x_{95}^{2}-32x_{95}x_{96}-510x_{96}\\ &-4x_{96}^{2}+x_{97}\end{aligned}$\\[3pt]
$S_{76}$ & $\begin{aligned}[t]&-6-5x_{31}-x_{31}^{2}-2x_{31}x_{32}-5x_{32}-x_{32}^{2}+x_{98}\end{aligned}$\\[3pt]
$S_{77}$ & $\begin{aligned}[t]&-30-11x_{29}-x_{29}^{2}-8x_{29}x_{32}-4x_{29}x_{98}-44x_{32}-16x_{32}^{2}-16x_{32}x_{98}-22x_{98}-4x_{98}^{2}+x_{99}\end{aligned}$\\[3pt]
$S_{78}$ & $\begin{aligned}[t]&-72-136x_{31}-64x_{31}^{2}-32x_{31}x_{79}-32x_{31}x_{81}-34x_{79}-4x_{79}^{2}-8x_{79}x_{81}-34x_{81}-4x_{81}^{2}+x_{100}\end{aligned}$\\[3pt]
$S_{79}$ & $\begin{aligned}[t]&-420-164x_{29}-16x_{29}^{2}-128x_{29}x_{31}-64x_{29}x_{79}-16x_{29}x_{99}-16x_{29}x_{100}-656x_{31}-256x_{31}^{2}-256x_{31}x_{79}\\ &-64x_{31}x_{99}-64x_{31}x_{100}-328x_{79}-64x_{79}^{2}-32x_{79}x_{99}-32x_{79}x_{100}-82x_{99}-4x_{99}^{2}-8x_{99}x_{100}-82x_{100}\\ &-4x_{100}^{2}+x_{101}\end{aligned}$\\[3pt]
$S_{80}$ & $\begin{aligned}[t]&-462-688x_{30}-256x_{30}^{2}-32x_{30}x_{31}-256x_{30}x_{79}-64x_{30}x_{84}-43x_{31}-x_{31}^{2}-16x_{31}x_{79}-4x_{31}x_{84}\\ &-344x_{79}-64x_{79}^{2}-32x_{79}x_{84}-86x_{84}-4x_{84}^{2}+x_{102}\end{aligned}$\\[3pt]
$S_{81}$ & $\begin{aligned}[t]&-42-52x_{31}-16x_{31}^{2}-16x_{31}x_{79}-26x_{79}-4x_{79}^{2}+x_{103}\end{aligned}$\\[3pt]
$S_{82}$ & $\begin{aligned}[t]&-30-11x_{31}-x_{31}^{2}-4x_{31}x_{103}-22x_{103}-4x_{103}^{2}+x_{104}\end{aligned}$\\[3pt]
$S_{83}$ & $\begin{aligned}[t]&-306-35x_{30}-x_{30}^{2}-16x_{30}x_{103}-4x_{30}x_{104}-280x_{103}-64x_{103}^{2}-32x_{103}x_{104}-70x_{104}-4x_{104}^{2}+x_{105}\end{aligned}$\\[3pt]
$S_{84}$ & $\begin{aligned}[t]&-4556-4320x_{103}-1024x_{103}^{2}-512x_{103}x_{104}-128x_{103}x_{105}-1080x_{104}-64x_{104}^{2}-32x_{104}x_{105}-270x_{105}\\ &-4x_{105}^{2}+x_{106}\end{aligned}$\\[3pt]
$S_{85}$ & $\begin{aligned}[t]&-30-11x_{31}-x_{31}^{2}-4x_{31}x_{106}-22x_{106}-4x_{106}^{2}+x_{107}\end{aligned}$\\[3pt]
$S_{86}$ & $\begin{aligned}[t]&-306-35x_{30}-x_{30}^{2}-16x_{30}x_{106}-4x_{30}x_{107}-280x_{106}-64x_{106}^{2}-32x_{106}x_{107}-70x_{107}-4x_{107}^{2}+x_{108}\end{aligned}$\\[3pt]
$S_{87}$ & $\begin{aligned}[t]&-22052-19008x_{30}-4096x_{30}^{2}-4096x_{30}x_{79}-1024x_{30}x_{84}-256x_{30}x_{102}-4096x_{30}x_{106}-1024x_{30}x_{107}\\ &-256x_{30}x_{108}-9504x_{79}-1024x_{79}^{2}-512x_{79}x_{84}-128x_{79}x_{102}-2048x_{79}x_{106}-512x_{79}x_{107}-128x_{79}x_{108}\\ &-2376x_{84}-64x_{84}^{2}-32x_{84}x_{102}-512x_{84}x_{106}-128x_{84}x_{107}-32x_{84}x_{108}-594x_{102}-4x_{102}^{2}-128x_{102}x_{106}\\ &-32x_{102}x_{107}-8x_{102}x_{108}-9504x_{106}-1024x_{106}^{2}-512x_{106}x_{107}-128x_{106}x_{108}-2376x_{107}-64x_{107}^{2}\\ &-32x_{107}x_{108}-594x_{108}-4x_{108}^{2}+x_{109}\end{aligned}$\\[3pt]
$S_{88}$ & $\begin{aligned}[t]&-70490-531x_{29}-x_{29}^{2}-256x_{29}x_{106}-64x_{29}x_{107}-16x_{29}x_{108}-4x_{29}x_{109}-67968x_{106}-16384x_{106}^{2}\\ &-8192x_{106}x_{107}-2048x_{106}x_{108}-512x_{106}x_{109}-16992x_{107}-1024x_{107}^{2}-512x_{107}x_{108}-128x_{107}x_{109}\\ &-4248x_{108}-64x_{108}^{2}-32x_{108}x_{109}-1062x_{109}-4x_{109}^{2}+x_{110}\end{aligned}$\\[3pt]
$S_{89}$ & $\begin{aligned}[t]&-1122540-1084928x_{106}-262144x_{106}^{2}-131072x_{106}x_{107}-32768x_{106}x_{108}-8192x_{106}x_{109}\\ &-2048x_{106}x_{110}-271232x_{107}-16384x_{107}^{2}-8192x_{107}x_{108}-2048x_{107}x_{109}-512x_{107}x_{110}-67808x_{108}\\ &-1024x_{108}^{2}-512x_{108}x_{109}-128x_{108}x_{110}-16952x_{109}-64x_{109}^{2}-32x_{109}x_{110}-4238x_{110}-4x_{110}^{2}+x_{111}\end{aligned}$\\[3pt]
$S_{90}$ & $\begin{aligned}[t]&-30-11x_{31}-x_{31}^{2}-4x_{31}x_{111}-22x_{111}-4x_{111}^{2}+x_{112}\end{aligned}$\\[3pt]
$S_{91}$ & $\begin{aligned}[t]&-306-35x_{29}-x_{29}^{2}-16x_{29}x_{111}-4x_{29}x_{112}-280x_{111}-64x_{111}^{2}-32x_{111}x_{112}-70x_{112}-4x_{112}^{2}+x_{113}\end{aligned}$\\[3pt]
$S_{92}$ & $\begin{aligned}[t]&-4556-4320x_{111}-1024x_{111}^{2}-512x_{111}x_{112}-128x_{111}x_{113}-1080x_{112}-64x_{112}^{2}-32x_{112}x_{113}-270x_{113}\\ &-4x_{113}^{2}+x_{114}\end{aligned}$\\[3pt]
$S_{93}$ & $\begin{aligned}[t]&-90-76x_{31}-16x_{31}^{2}-16x_{31}x_{79}-16x_{31}x_{114}-38x_{79}-4x_{79}^{2}-8x_{79}x_{114}-38x_{114}-4x_{114}^{2}+x_{115}\end{aligned}$\\[3pt]
$S_{94}$ & $\begin{aligned}[t]&-306-35x_{30}-x_{30}^{2}-16x_{30}x_{114}-4x_{30}x_{115}-280x_{114}-64x_{114}^{2}-32x_{114}x_{115}-70x_{115}-4x_{115}^{2}+x_{116}\end{aligned}$\\[3pt]
$S_{95}$ & $\begin{aligned}[t]&-4556-4320x_{106}-1024x_{106}^{2}-512x_{106}x_{107}-128x_{106}x_{108}-1080x_{107}-64x_{107}^{2}-32x_{107}x_{108}-270x_{108}\\ &-4x_{108}^{2}+x_{117}\end{aligned}$\\[3pt]
$S_{96}$ & $\begin{aligned}[t]&-4970-4512x_{114}-1024x_{114}^{2}-512x_{114}x_{115}-128x_{114}x_{116}-128x_{114}x_{117}-1128x_{115}-64x_{115}^{2}\\ &-32x_{115}x_{116}-32x_{115}x_{117}-282x_{116}-4x_{116}^{2}-8x_{116}x_{117}-282x_{117}-4x_{117}^{2}+x_{118}\end{aligned}$\\[3pt]
$S_{97}$ & $\begin{aligned}[t]&-306-35x_{31}-x_{31}^{2}-16x_{31}x_{117}-4x_{31}x_{118}-280x_{117}-64x_{117}^{2}-32x_{117}x_{118}-70x_{118}-4x_{118}^{2}+x_{119}\end{aligned}$\\[3pt]
$S_{98}$ & $\begin{aligned}[t]&-42-52x_{30}-16x_{30}^{2}-8x_{30}x_{31}-16x_{30}x_{75}-13x_{31}-x_{31}^{2}-4x_{31}x_{75}-26x_{75}-4x_{75}^{2}+x_{120}\end{aligned}$\\[3pt]
$S_{99}$ & $\begin{aligned}[t]&-756-1100x_{30}-400x_{30}^{2}-400x_{30}x_{75}-80x_{30}x_{120}-550x_{75}-100x_{75}^{2}-40x_{75}x_{120}-110x_{120}-4x_{120}^{2}\\ &+x_{121}\end{aligned}$\\[3pt]
$S_{100}$ & $\begin{aligned}[t]&-23870-19776x_{30}-4096x_{30}^{2}-4096x_{30}x_{75}-4096x_{30}x_{117}-1024x_{30}x_{118}-256x_{30}x_{119}-1024x_{30}x_{120}\\ &-256x_{30}x_{121}-9888x_{75}-1024x_{75}^{2}-2048x_{75}x_{117}-512x_{75}x_{118}-128x_{75}x_{119}-512x_{75}x_{120}-128x_{75}x_{121}\\ &-9888x_{117}-1024x_{117}^{2}-512x_{117}x_{118}-128x_{117}x_{119}-512x_{117}x_{120}-128x_{117}x_{121}-2472x_{118}-64x_{118}^{2}\\ &-32x_{118}x_{119}-128x_{118}x_{120}-32x_{118}x_{121}-618x_{119}-4x_{119}^{2}-32x_{119}x_{120}-8x_{119}x_{121}-2472x_{120}\\ &-64x_{120}^{2}-32x_{120}x_{121}-618x_{121}-4x_{121}^{2}+x_{122}\end{aligned}$\\[3pt]
$S_{101}$ & $\begin{aligned}[t]&-70490-67968x_{117}-16384x_{117}^{2}-8192x_{117}x_{118}-2048x_{117}x_{119}-512x_{117}x_{122}-16992x_{118}-1024x_{118}^{2}\\ &-512x_{118}x_{119}-128x_{118}x_{122}-4248x_{119}-64x_{119}^{2}-32x_{119}x_{122}-1062x_{122}-4x_{122}^{2}+x_{123}\end{aligned}$\\[3pt]
$S_{102}$ & $\begin{aligned}[t]&-1-x_{3}-x_{3}x_{34}-x_{34}+x_{124}\end{aligned}$\\[3pt]
$S_{103}$ & $\begin{aligned}[t]&-x_{33}^{2}+2x_{33}x_{34}-x_{34}^{2}+x_{125}\end{aligned}$\\[3pt]
$S_{104}$ & $\begin{aligned}[t]&-2x_{78}x_{124}-4x_{124}+x_{126}\end{aligned}$\\[3pt]
$S_{105}$ & $\begin{aligned}[t]&-x_{33}^{2}+2x_{33}x_{34}+2x_{33}x_{126}-x_{34}^{2}-2x_{34}x_{126}-x_{126}^{2}+x_{127}\end{aligned}$\\[3pt]
$S_{106}$ & $\begin{aligned}[t]&-64x_{31}x_{124}-32x_{81}x_{124}-8x_{82}x_{124}-2x_{83}x_{124}-96x_{124}+x_{128}\end{aligned}$\\[3pt]
$S_{107}$ & $\begin{aligned}[t]&-x_{33}^{2}+2x_{33}x_{34}+2x_{33}x_{128}-x_{34}^{2}-2x_{34}x_{128}-x_{128}^{2}+x_{129}\end{aligned}$\\[3pt]
$S_{108}$ & $\begin{aligned}[t]&-64x_{30}x_{124}-32x_{79}x_{124}-8x_{84}x_{124}-2x_{85}x_{124}-80x_{124}+x_{130}\end{aligned}$\\[3pt]
$S_{109}$ & $\begin{aligned}[t]&-x_{33}^{2}+2x_{33}x_{34}+2x_{33}x_{130}-x_{34}^{2}-2x_{34}x_{130}-x_{130}^{2}+x_{131}\end{aligned}$\\[3pt]
$S_{110}$ & $\begin{aligned}[t]&-64x_{31}x_{124}-32x_{81}x_{124}-8x_{87}x_{124}-2x_{88}x_{124}-104x_{124}+x_{132}\end{aligned}$\\[3pt]
$S_{111}$ & $\begin{aligned}[t]&-x_{33}^{2}+2x_{33}x_{34}+2x_{33}x_{132}-x_{34}^{2}-2x_{34}x_{132}-x_{132}^{2}+x_{133}\end{aligned}$\\[3pt]
$S_{112}$ & $\begin{aligned}[t]&-64x_{30}x_{124}-32x_{75}x_{124}-8x_{89}x_{124}-2x_{90}x_{124}-80x_{124}+x_{134}\end{aligned}$\\[3pt]
$S_{113}$ & $\begin{aligned}[t]&-x_{33}^{2}+2x_{33}x_{34}+2x_{33}x_{134}-x_{34}^{2}-2x_{34}x_{134}-x_{134}^{2}+x_{135}\end{aligned}$\\[3pt]
$S_{114}$ & $\begin{aligned}[t]&-4x_{30}x_{124}-2x_{94}x_{124}-4x_{124}+x_{136}\end{aligned}$\\[3pt]
$S_{115}$ & $\begin{aligned}[t]&-x_{33}^{2}+2x_{33}x_{34}+2x_{33}x_{136}-x_{34}^{2}-2x_{34}x_{136}-x_{136}^{2}+x_{137}\end{aligned}$\\[3pt]
$S_{116}$ & $\begin{aligned}[t]&-256x_{30}x_{124}-128x_{75}x_{124}-32x_{95}x_{124}-8x_{96}x_{124}-2x_{97}x_{124}-480x_{124}+x_{138}\end{aligned}$\\[3pt]
$S_{117}$ & $\begin{aligned}[t]&-x_{33}^{2}+2x_{33}x_{34}+2x_{33}x_{138}-x_{34}^{2}-2x_{34}x_{138}-x_{138}^{2}+x_{139}\end{aligned}$\\[3pt]
$S_{118}$ & $\begin{aligned}[t]&-64x_{31}x_{124}-32x_{79}x_{124}-8x_{100}x_{124}-2x_{101}x_{124}-64x_{124}+x_{140}\end{aligned}$\\[3pt]
$S_{119}$ & $\begin{aligned}[t]&-x_{33}^{2}+2x_{33}x_{34}+2x_{33}x_{140}-x_{34}^{2}-2x_{34}x_{140}-x_{140}^{2}+x_{141}\end{aligned}$\\[3pt]
$S_{120}$ & $\begin{aligned}[t]&-x_{29}^{2}+2x_{29}x_{31}-x_{31}^{2}+x_{142}\end{aligned}$\\[3pt]
$S_{121}$ & $\begin{aligned}[t]&-x_{30}^{2}+2x_{30}x_{31}-x_{31}^{2}+x_{143}\end{aligned}$\\[3pt]
$S_{122}$ & $\begin{aligned}[t]&-x_{142}x_{143}+x_{144}\end{aligned}$\\[3pt]
$S_{123}$ & $\begin{aligned}[t]&-4-4x_{32}-x_{32}^{2}+x_{145}\end{aligned}$\\[3pt]
$S_{124}$ & $\begin{aligned}[t]&-x_{144}^{2}+2x_{144}x_{145}-x_{145}^{2}+x_{146}\end{aligned}$\\[3pt]
$S_{125}$ & $\begin{aligned}[t]&-2x_{123}x_{124}-4x_{124}+x_{147}\end{aligned}$\\[3pt]
$S_{126}$ & $\begin{aligned}[t]&-x_{33}^{2}+2x_{33}x_{34}+2x_{33}x_{147}-x_{34}^{2}-2x_{34}x_{147}-x_{147}^{2}+x_{148}\end{aligned}$\\[3pt]
$S_{127}$ & $\begin{aligned}[t]&-1-x_{3}-x_{3}x_{30}-x_{30}+x_{149}\end{aligned}$\\[3pt]
$S_{128}$ & $\begin{aligned}[t]&-1-x_{34}-x_{34}x_{149}-x_{149}+x_{150}\end{aligned}$\\[3pt]
$S_{129}$ & $\begin{aligned}[t]&-1-x_{3}-x_{3}x_{29}-x_{29}+x_{151}\end{aligned}$\\[3pt]
$S_{130}$ & $\begin{aligned}[t]&-4-4x_{3}-4x_{3}x_{30}-2x_{3}x_{75}-4x_{30}-2x_{75}+x_{152}\end{aligned}$\\[3pt]
$S_{131}$ & $\begin{aligned}[t]&-1-x_{33}-x_{33}x_{151}-x_{151}+x_{153}\end{aligned}$\\[3pt]
$S_{132}$ & $\begin{aligned}[t]&-x_{152}x_{153}-x_{153}+x_{154}\end{aligned}$\\[3pt]
$S_{133}$ & $\begin{aligned}[t]&-x_{150}^{2}+2x_{150}x_{154}-x_{154}^{2}+x_{155}\end{aligned}$\\[3pt]
$S_{134}$ & $\begin{aligned}[t]&-6-6x_{3}-4x_{3}x_{30}-2x_{3}x_{75}-4x_{30}-2x_{75}+x_{156}\end{aligned}$\\[3pt]
$S_{135}$ & $\begin{aligned}[t]&-1-x_{34}-x_{34}x_{156}-x_{156}+x_{157}\end{aligned}$\\[3pt]
$S_{136}$ & $\begin{aligned}[t]&-1-x_{33}-x_{33}x_{149}-x_{149}+x_{158}\end{aligned}$\\[3pt]
$S_{137}$ & $\begin{aligned}[t]&-x_{157}^{2}+2x_{157}x_{158}-x_{158}^{2}+x_{159}\end{aligned}$\\[3pt]
$S_{138}$ & $\begin{aligned}[t]&-x_{125}x_{127}+x_{160}\end{aligned}$\\[3pt]
$S_{139}$ & $\begin{aligned}[t]&-x_{129}x_{160}+x_{161}\end{aligned}$\\[3pt]
$S_{140}$ & $\begin{aligned}[t]&-x_{131}x_{161}+x_{162}\end{aligned}$\\[3pt]
$S_{141}$ & $\begin{aligned}[t]&-x_{133}x_{162}+x_{163}\end{aligned}$\\[3pt]
$S_{142}$ & $\begin{aligned}[t]&-x_{135}x_{163}+x_{164}\end{aligned}$\\[3pt]
$S_{143}$ & $\begin{aligned}[t]&-x_{137}x_{164}+x_{165}\end{aligned}$\\[3pt]
$S_{144}$ & $\begin{aligned}[t]&-x_{139}x_{165}+x_{166}\end{aligned}$\\[3pt]
$S_{145}$ & $\begin{aligned}[t]&-x_{141}x_{166}+x_{167}\end{aligned}$\\[3pt]
$S_{146}$ & $\begin{aligned}[t]&-x_{146}x_{167}-x_{148}x_{167}+x_{168}\end{aligned}$\\[3pt]
$S_{147}$ & $\begin{aligned}[t]&-x_{155}x_{168}+x_{169}\end{aligned}$\\[3pt]
$S_{148}$ & $\begin{aligned}[t]&-x_{2}-x_{2}x_{37}+x_{170}\end{aligned}$\\[3pt]
$S_{149}$ & $\begin{aligned}[t]&-1-x_{35}-x_{35}x_{170}-x_{170}+x_{171}\end{aligned}$\\[3pt]
$S_{150}$ & $\begin{aligned}[t]&-2+4x_{1}-2x_{1}^{2}+2x_{1}x_{33}+2x_{1}x_{34}+2x_{1}x_{171}+2x_{1}x_{172}-2x_{33}-x_{33}^{2}-2x_{33}x_{171}-2x_{34}-x_{34}^{2}-2x_{34}x_{172}\\ &+x_{73}x_{173}-x_{159}x_{169}-2x_{171}-x_{171}^{2}-2x_{172}-x_{172}^{2}\end{aligned}$\\[3pt]
$S_{151}$ & $\begin{aligned}[t]&-1-x_{2}-x_{2}x_{36}-x_{36}-x_{36}x_{170}-x_{170}+x_{172}\end{aligned}$\\[3pt]
$S_{152}$ & $\begin{aligned}[t]&1+x_{62}-x_{63}x_{73}-x_{73}+x_{265}\end{aligned}$\\[3pt]
$S_{153}$ & $\begin{aligned}[t]&-x_{62}x_{174}-x_{174}-x_{174}x_{265}+x_{179}\end{aligned}$\\[3pt]
$S_{154}$ & $\begin{aligned}[t]&-x_{180}-x_{180}x_{183}+x_{184}\end{aligned}$\\[3pt]
$S_{155}$ & $\begin{aligned}[t]&-x_{184}^{2}+x_{185}\end{aligned}$\\[3pt]
$S_{156}$ & $\begin{aligned}[t]&1-x_{181}^{2}-x_{185}+x_{185}x_{186}\end{aligned}$\\[3pt]
$S_{157}$ & $\begin{aligned}[t]&-1-2x_{180}-x_{180}^{2}+x_{186}\end{aligned}$\\[3pt]
$S_{158}$ & $\begin{aligned}[t]&5+4x_{62}+x_{62}^{2}-2x_{62}x_{181}-2x_{62}x_{182}+2x_{62}x_{265}+x_{179}x_{187}-4x_{181}-2x_{181}x_{265}-4x_{182}-2x_{182}x_{265}\\ &+4x_{265}+x_{265}^{2}\end{aligned}$\\[3pt]
$S_{159}$ & $\begin{aligned}[t]&-x_{189}^{2}+x_{192}\end{aligned}$\\[3pt]
$S_{160}$ & $\begin{aligned}[t]&1-x_{188}^{2}+x_{192}x_{215}\end{aligned}$\\[3pt]
$S_{161}$ & $\begin{aligned}[t]&1-x_{181}^{2}-2x_{181}x_{182}-x_{182}^{2}+x_{215}\end{aligned}$\\[3pt]
$S_{162}$ & $\begin{aligned}[t]&-4x_{189}-4x_{189}x_{190}+x_{193}\end{aligned}$\\[3pt]
$S_{163}$ & $\begin{aligned}[t]&1-x_{181}-x_{182}-x_{188}x_{191}+x_{193}\end{aligned}$\\[3pt]
$S_{164}$ & $\begin{aligned}[t]&-x_{195}^{2}+x_{198}\end{aligned}$\\[3pt]
$S_{165}$ & $\begin{aligned}[t]&1-x_{194}^{2}+x_{198}x_{215}\end{aligned}$\\[3pt]
$S_{166}$ & $\begin{aligned}[t]&-x_{197}^{2}+x_{199}\end{aligned}$\\[3pt]
$S_{167}$ & $\begin{aligned}[t]&1-x_{196}^{2}-x_{199}+x_{199}x_{200}\end{aligned}$\\[3pt]
$S_{168}$ & $\begin{aligned}[t]&-1-2x_{193}-x_{193}^{2}+x_{200}\end{aligned}$\\[3pt]
$S_{169}$ & $\begin{aligned}[t]&-x_{188}x_{201}-x_{194}+x_{196}\end{aligned}$\\[3pt]
$S_{170}$ & $\begin{aligned}[t]&-1-x_{0}-4x_{195}x_{202}+x_{197}\end{aligned}$\\[3pt]
$S_{171}$ & $\begin{aligned}[t]&5+4x_{0}+x_{0}^{2}-2x_{0}x_{181}-2x_{0}x_{182}-4x_{181}-4x_{182}+x_{203}\end{aligned}$\\[3pt]
$S_{172}$ & $\begin{aligned}[t]&-x_{0}x_{195}+x_{4}-x_{194}-2x_{195}+x_{203}x_{204}+x_{205}\end{aligned}$\\[3pt]
$S_{173}$ & $\begin{aligned}[t]&-x_{181}x_{195}-x_{182}x_{195}+x_{205}\end{aligned}$\\[3pt]
$S_{174}$ & $\begin{aligned}[t]&-x_{207}^{2}+x_{211}\end{aligned}$\\[3pt]
$S_{175}$ & $\begin{aligned}[t]&1-x_{206}^{2}+x_{211}x_{215}\end{aligned}$\\[3pt]
$S_{176}$ & $\begin{aligned}[t]&-x_{209}^{2}+x_{212}\end{aligned}$\\[3pt]
$S_{177}$ & $\begin{aligned}[t]&1+x_{200}x_{212}-x_{208}^{2}-x_{212}\end{aligned}$\\[3pt]
$S_{178}$ & $\begin{aligned}[t]&-x_{188}x_{213}-x_{206}+x_{208}\end{aligned}$\\[3pt]
$S_{179}$ & $\begin{aligned}[t]&-x_{207}-4x_{207}x_{214}+x_{209}+x_{210}\end{aligned}$\\[3pt]
$S_{180}$ & $\begin{aligned}[t]&x_{61}+x_{181}x_{217}+x_{181}x_{218}+x_{182}x_{217}+x_{182}x_{218}-x_{206}-2x_{206}x_{215}-x_{207}x_{217}-x_{207}x_{218}+x_{210}x_{217}\\ &+x_{210}x_{218}-x_{215}x_{216}-2x_{217}-2x_{218}+x_{219}x_{220}\end{aligned}$\\[3pt]
$S_{181}$ & $\begin{aligned}[t]&-2x_{181}x_{207}-2x_{182}x_{207}+x_{216}\end{aligned}$\\[3pt]
$S_{182}$ & $\begin{aligned}[t]&-2x_{181}x_{206}-2x_{182}x_{206}+x_{217}\end{aligned}$\\[3pt]
$S_{183}$ & $\begin{aligned}[t]&-x_{207}-2x_{207}x_{215}+x_{218}\end{aligned}$\\[3pt]
$S_{184}$ & $\begin{aligned}[t]&5-4x_{181}-2x_{181}x_{207}+2x_{181}x_{210}-4x_{182}-2x_{182}x_{207}+2x_{182}x_{210}+4x_{207}+x_{207}^{2}-2x_{207}x_{210}-4x_{210}\\ &+x_{210}^{2}+x_{219}\end{aligned}$\\[3pt]
$S_{185}$ & $\begin{aligned}[t]&2+2x_{61}+x_{61}^{2}-2x_{61}x_{181}-2x_{61}x_{182}-2x_{181}-2x_{182}+x_{221}\end{aligned}$\\[3pt]
$S_{186}$ & $\begin{aligned}[t]&-x_{61}x_{207}+x_{63}+x_{181}x_{207}+x_{182}x_{207}-x_{206}-x_{207}+x_{221}x_{222}\end{aligned}$\\[3pt]
$S_{187}$ & $\begin{aligned}[t]&2+2x_{64}+x_{64}^{2}-2x_{64}x_{181}-2x_{64}x_{182}-2x_{181}-2x_{182}+x_{223}\end{aligned}$\\[3pt]
$S_{188}$ & $\begin{aligned}[t]&-x_{64}x_{195}+x_{65}-x_{194}-x_{195}+x_{205}+x_{223}x_{224}\end{aligned}$\\[3pt]
$S_{189}$ & $\begin{aligned}[t]&1+x_{60}^{2}-2x_{60}x_{181}-2x_{60}x_{182}+x_{225}\end{aligned}$\\[3pt]
$S_{190}$ & $\begin{aligned}[t]&-x_{60}x_{195}+x_{66}-x_{194}+x_{205}+x_{225}x_{226}\end{aligned}$\\[3pt]
$S_{191}$ & $\begin{aligned}[t]&-x_{0}x_{228}-x_{0}x_{229}+x_{68}+x_{181}x_{228}+x_{181}x_{229}+x_{182}x_{228}+x_{182}x_{229}-x_{194}-2x_{194}x_{215}-x_{215}x_{227}-3x_{228}\\ &-3x_{229}+x_{230}x_{231}\end{aligned}$\\[3pt]
$S_{192}$ & $\begin{aligned}[t]&-2x_{181}x_{195}-2x_{182}x_{195}+x_{227}\end{aligned}$\\[3pt]
$S_{193}$ & $\begin{aligned}[t]&-2x_{181}x_{194}-2x_{182}x_{194}+x_{228}\end{aligned}$\\[3pt]
$S_{194}$ & $\begin{aligned}[t]&-x_{195}-2x_{195}x_{215}+x_{229}\end{aligned}$\\[3pt]
$S_{195}$ & $\begin{aligned}[t]&10+6x_{0}+x_{0}^{2}-2x_{0}x_{181}-2x_{0}x_{182}-6x_{181}-6x_{182}+x_{230}\end{aligned}$\\[3pt]
$S_{196}$ & $\begin{aligned}[t]&2+2x_{68}+x_{68}^{2}-2x_{68}x_{181}-2x_{68}x_{182}-2x_{181}-2x_{182}+x_{232}\end{aligned}$\\[3pt]
$S_{197}$ & $\begin{aligned}[t]&-x_{68}x_{195}-x_{194}-x_{195}+x_{205}+x_{232}x_{233}+x_{441}^{2}+x_{442}^{2}+x_{443}+x_{443}^{2}\end{aligned}$\\[3pt]
$S_{198}$ & $\begin{aligned}[t]&2+2x_{62}+x_{62}^{2}-2x_{62}x_{181}-2x_{62}x_{182}+2x_{62}x_{265}-2x_{181}-2x_{181}x_{265}-2x_{182}-2x_{182}x_{265}+x_{234}+2x_{265}\\ &+x_{265}^{2}\end{aligned}$\\[3pt]
$S_{199}$ & $\begin{aligned}[t]&-x_{62}x_{207}+x_{174}+x_{181}x_{207}+x_{182}x_{207}-x_{206}-x_{207}-x_{207}x_{265}+x_{234}x_{235}\end{aligned}$\\[3pt]
$S_{200}$ & $\begin{aligned}[t]&-x_{62}x_{195}+x_{175}-x_{194}-x_{195}-x_{195}x_{265}+x_{205}+x_{234}x_{236}\end{aligned}$\\[3pt]
$S_{201}$ & $\begin{aligned}[t]&-x_{62}x_{194}-x_{62}x_{205}+x_{176}+x_{181}x_{205}+x_{182}x_{205}-x_{194}-x_{194}x_{265}-x_{195}x_{215}-x_{205}-x_{205}x_{265}+x_{234}x_{237}\end{aligned}$\\[3pt]
$S_{202}$ & $\begin{aligned}[t]&-x_{239}^{2}+x_{242}\end{aligned}$\\[3pt]
$S_{203}$ & $\begin{aligned}[t]&1+x_{215}x_{242}-x_{238}^{2}\end{aligned}$\\[3pt]
$S_{204}$ & $\begin{aligned}[t]&-x_{241}^{2}+x_{243}\end{aligned}$\\[3pt]
$S_{205}$ & $\begin{aligned}[t]&1+x_{200}x_{243}-x_{240}^{2}-x_{243}\end{aligned}$\\[3pt]
$S_{206}$ & $\begin{aligned}[t]&-x_{188}x_{244}-x_{238}+x_{240}\end{aligned}$\\[3pt]
$S_{207}$ & $\begin{aligned}[t]&-1-x_{2}-4x_{239}x_{245}+x_{241}\end{aligned}$\\[3pt]
$S_{208}$ & $\begin{aligned}[t]&-x_{62}x_{239}+x_{177}+x_{181}x_{239}+x_{182}x_{239}+x_{234}x_{246}-x_{238}-x_{239}-x_{239}x_{265}\end{aligned}$\\[3pt]
$S_{209}$ & $\begin{aligned}[t]&-x_{248}^{2}+x_{251}\end{aligned}$\\[3pt]
$S_{210}$ & $\begin{aligned}[t]&1+x_{215}x_{251}-x_{247}^{2}\end{aligned}$\\[3pt]
$S_{211}$ & $\begin{aligned}[t]&-x_{250}^{2}+x_{252}\end{aligned}$\\[3pt]
$S_{212}$ & $\begin{aligned}[t]&1+x_{200}x_{252}-x_{249}^{2}-x_{252}\end{aligned}$\\[3pt]
$S_{213}$ & $\begin{aligned}[t]&-x_{188}x_{253}-x_{247}+x_{249}\end{aligned}$\\[3pt]
$S_{214}$ & $\begin{aligned}[t]&-1-x_{1}-4x_{248}x_{254}+x_{250}\end{aligned}$\\[3pt]
$S_{215}$ & $\begin{aligned}[t]&-x_{62}x_{248}+x_{178}+x_{181}x_{248}+x_{182}x_{248}+x_{234}x_{255}-x_{247}-x_{248}-x_{248}x_{265}\end{aligned}$\\[3pt]
$S_{216}$ & $\begin{aligned}[t]&-x_{257}^{2}+x_{260}\end{aligned}$\\[3pt]
$S_{217}$ & $\begin{aligned}[t]&1+x_{215}x_{260}-x_{256}^{2}\end{aligned}$\\[3pt]
$S_{218}$ & $\begin{aligned}[t]&-x_{259}^{2}+x_{261}\end{aligned}$\\[3pt]
$S_{219}$ & $\begin{aligned}[t]&1+x_{200}x_{261}-x_{258}^{2}-x_{261}\end{aligned}$\\[3pt]
$S_{220}$ & $\begin{aligned}[t]&-x_{188}x_{262}-x_{256}+x_{258}\end{aligned}$\\[3pt]
$S_{221}$ & $\begin{aligned}[t]&-x_{61}-x_{207}+x_{210}-4x_{257}x_{263}+x_{259}\end{aligned}$\\[3pt]
$S_{222}$ & $\begin{aligned}[t]&-x_{62}x_{174}+x_{62}x_{257}-x_{179}x_{264}-x_{181}x_{257}-x_{182}x_{257}+x_{256}+2x_{257}+x_{257}x_{265}-x_{525}^{2}-x_{526}^{2}-x_{527}\\ &-x_{527}^{2}\end{aligned}$\\[3pt]
$S_{223}$ & $\begin{aligned}[t]&-x_{267}^{2}+x_{270}\end{aligned}$\\[3pt]
$S_{224}$ & $\begin{aligned}[t]&1+x_{215}x_{270}-x_{266}^{2}\end{aligned}$\\[3pt]
$S_{225}$ & $\begin{aligned}[t]&-x_{269}^{2}+x_{271}\end{aligned}$\\[3pt]
$S_{226}$ & $\begin{aligned}[t]&1+x_{200}x_{271}-x_{268}^{2}-x_{271}\end{aligned}$\\[3pt]
$S_{227}$ & $\begin{aligned}[t]&-x_{188}x_{272}-x_{266}+x_{268}\end{aligned}$\\[3pt]
$S_{228}$ & $\begin{aligned}[t]&-1-x_{0}-x_{68}-4x_{267}x_{273}+x_{269}\end{aligned}$\\[3pt]
$S_{229}$ & $\begin{aligned}[t]&-x_{62}x_{175}-x_{175}-x_{175}x_{265}+x_{275}\end{aligned}$\\[3pt]
$S_{230}$ & $\begin{aligned}[t]&x_{62}x_{267}-x_{69}x_{175}-x_{181}x_{267}-x_{182}x_{267}+x_{265}x_{267}+x_{266}+2x_{267}-x_{274}x_{275}-x_{540}^{2}-x_{541}^{2}-x_{542}-x_{542}^{2}\end{aligned}$\\[3pt]
$S_{231}$ & $\begin{aligned}[t]&-x_{277}^{2}+x_{280}\end{aligned}$\\[3pt]
$S_{232}$ & $\begin{aligned}[t]&1+x_{215}x_{280}-x_{276}^{2}\end{aligned}$\\[3pt]
$S_{233}$ & $\begin{aligned}[t]&-x_{279}^{2}+x_{281}\end{aligned}$\\[3pt]
$S_{234}$ & $\begin{aligned}[t]&1+x_{200}x_{281}-x_{278}^{2}-x_{281}\end{aligned}$\\[3pt]
$S_{235}$ & $\begin{aligned}[t]&-x_{188}x_{282}-x_{276}+x_{278}\end{aligned}$\\[3pt]
$S_{236}$ & $\begin{aligned}[t]&-3-x_{0}-x_{65}-4x_{277}x_{283}+x_{279}\end{aligned}$\\[3pt]
$S_{237}$ & $\begin{aligned}[t]&x_{62}x_{277}-x_{70}x_{175}-x_{181}x_{277}-x_{182}x_{277}+x_{265}x_{277}-x_{275}x_{284}+x_{276}+2x_{277}-x_{555}^{2}-x_{556}^{2}-x_{557}-x_{557}^{2}\end{aligned}$\\[3pt]
$S_{238}$ & $\begin{aligned}[t]&-x_{286}^{2}+x_{289}\end{aligned}$\\[3pt]
$S_{239}$ & $\begin{aligned}[t]&1+x_{215}x_{289}-x_{285}^{2}\end{aligned}$\\[3pt]
$S_{240}$ & $\begin{aligned}[t]&-x_{288}^{2}+x_{290}\end{aligned}$\\[3pt]
$S_{241}$ & $\begin{aligned}[t]&1+x_{200}x_{290}-x_{287}^{2}-x_{290}\end{aligned}$\\[3pt]
$S_{242}$ & $\begin{aligned}[t]&-x_{188}x_{291}-x_{285}+x_{287}\end{aligned}$\\[3pt]
$S_{243}$ & $\begin{aligned}[t]&-1-x_{29}-4x_{286}x_{292}+x_{288}\end{aligned}$\\[3pt]
$S_{244}$ & $\begin{aligned}[t]&-x_{73}x_{176}+x_{294}\end{aligned}$\\[3pt]
$S_{245}$ & $\begin{aligned}[t]&-x_{62}x_{286}+x_{181}x_{286}+x_{182}x_{286}-x_{265}x_{286}-x_{285}-2x_{286}+x_{293}x_{294}+x_{570}^{2}+x_{571}^{2}+x_{572}+x_{572}^{2}\end{aligned}$\\[3pt]
$S_{246}$ & $\begin{aligned}[t]&-x_{296}^{2}+x_{299}\end{aligned}$\\[3pt]
$S_{247}$ & $\begin{aligned}[t]&1+x_{215}x_{299}-x_{295}^{2}\end{aligned}$\\[3pt]
$S_{248}$ & $\begin{aligned}[t]&-x_{298}^{2}+x_{300}\end{aligned}$\\[3pt]
$S_{249}$ & $\begin{aligned}[t]&1+x_{200}x_{300}-x_{297}^{2}-x_{300}\end{aligned}$\\[3pt]
$S_{250}$ & $\begin{aligned}[t]&-x_{188}x_{301}-x_{295}+x_{297}\end{aligned}$\\[3pt]
$S_{251}$ & $\begin{aligned}[t]&-1-x_{30}-4x_{296}x_{302}+x_{298}\end{aligned}$\\[3pt]
$S_{252}$ & $\begin{aligned}[t]&-x_{62}x_{296}+x_{181}x_{296}+x_{182}x_{296}-x_{265}x_{296}+x_{294}x_{303}-x_{295}-2x_{296}+x_{582}^{2}+x_{583}^{2}+x_{584}+x_{584}^{2}\end{aligned}$\\[3pt]
$S_{253}$ & $\begin{aligned}[t]&-x_{305}^{2}+x_{308}\end{aligned}$\\[3pt]
$S_{254}$ & $\begin{aligned}[t]&1+x_{215}x_{308}-x_{304}^{2}\end{aligned}$\\[3pt]
$S_{255}$ & $\begin{aligned}[t]&-x_{307}^{2}+x_{309}\end{aligned}$\\[3pt]
$S_{256}$ & $\begin{aligned}[t]&1+x_{200}x_{309}-x_{306}^{2}-x_{309}\end{aligned}$\\[3pt]
$S_{257}$ & $\begin{aligned}[t]&-x_{188}x_{310}-x_{304}+x_{306}\end{aligned}$\\[3pt]
$S_{258}$ & $\begin{aligned}[t]&-1-x_{31}-4x_{305}x_{311}+x_{307}\end{aligned}$\\[3pt]
$S_{259}$ & $\begin{aligned}[t]&-x_{62}x_{305}+x_{181}x_{305}+x_{182}x_{305}-x_{265}x_{305}+x_{294}x_{312}-x_{304}-2x_{305}+x_{594}^{2}+x_{595}^{2}+x_{596}+x_{596}^{2}\end{aligned}$\\[3pt]
$S_{260}$ & $\begin{aligned}[t]&-x_{314}^{2}+x_{317}\end{aligned}$\\[3pt]
$S_{261}$ & $\begin{aligned}[t]&1+x_{215}x_{317}-x_{313}^{2}\end{aligned}$\\[3pt]
$S_{262}$ & $\begin{aligned}[t]&-x_{316}^{2}+x_{318}\end{aligned}$\\[3pt]
$S_{263}$ & $\begin{aligned}[t]&1+x_{200}x_{318}-x_{315}^{2}-x_{318}\end{aligned}$\\[3pt]
$S_{264}$ & $\begin{aligned}[t]&-x_{188}x_{319}-x_{313}+x_{315}\end{aligned}$\\[3pt]
$S_{265}$ & $\begin{aligned}[t]&-1-x_{32}-4x_{314}x_{320}+x_{316}\end{aligned}$\\[3pt]
$S_{266}$ & $\begin{aligned}[t]&-x_{62}x_{314}+x_{181}x_{314}+x_{182}x_{314}-x_{265}x_{314}+x_{294}x_{321}-x_{313}-2x_{314}+x_{606}^{2}+x_{607}^{2}+x_{608}+x_{608}^{2}\end{aligned}$\\[3pt]
$S_{267}$ & $\begin{aligned}[t]&-x_{323}^{2}+x_{326}\end{aligned}$\\[3pt]
$S_{268}$ & $\begin{aligned}[t]&1+x_{215}x_{326}-x_{322}^{2}\end{aligned}$\\[3pt]
$S_{269}$ & $\begin{aligned}[t]&-x_{325}^{2}+x_{327}\end{aligned}$\\[3pt]
$S_{270}$ & $\begin{aligned}[t]&1+x_{200}x_{327}-x_{324}^{2}-x_{327}\end{aligned}$\\[3pt]
$S_{271}$ & $\begin{aligned}[t]&-x_{188}x_{328}-x_{322}+x_{324}\end{aligned}$\\[3pt]
$S_{272}$ & $\begin{aligned}[t]&-1-x_{33}-4x_{323}x_{329}+x_{325}\end{aligned}$\\[3pt]
$S_{273}$ & $\begin{aligned}[t]&-x_{73}x_{177}+x_{331}\end{aligned}$\\[3pt]
$S_{274}$ & $\begin{aligned}[t]&-x_{62}x_{323}+x_{181}x_{323}+x_{182}x_{323}-x_{265}x_{323}-x_{322}-2x_{323}+x_{330}x_{331}+x_{618}^{2}+x_{619}^{2}+x_{620}+x_{620}^{2}\end{aligned}$\\[3pt]
$S_{275}$ & $\begin{aligned}[t]&-x_{333}^{2}+x_{336}\end{aligned}$\\[3pt]
$S_{276}$ & $\begin{aligned}[t]&1+x_{215}x_{336}-x_{332}^{2}\end{aligned}$\\[3pt]
$S_{277}$ & $\begin{aligned}[t]&-x_{335}^{2}+x_{337}\end{aligned}$\\[3pt]
$S_{278}$ & $\begin{aligned}[t]&1+x_{200}x_{337}-x_{334}^{2}-x_{337}\end{aligned}$\\[3pt]
$S_{279}$ & $\begin{aligned}[t]&-x_{188}x_{338}-x_{332}+x_{334}\end{aligned}$\\[3pt]
$S_{280}$ & $\begin{aligned}[t]&-1-x_{34}-4x_{333}x_{339}+x_{335}\end{aligned}$\\[3pt]
$S_{281}$ & $\begin{aligned}[t]&-x_{62}x_{333}+x_{181}x_{333}+x_{182}x_{333}-x_{265}x_{333}+x_{331}x_{340}-x_{332}-2x_{333}+x_{630}^{2}+x_{631}^{2}+x_{632}+x_{632}^{2}\end{aligned}$\\[3pt]
$S_{282}$ & $\begin{aligned}[t]&-x_{342}^{2}+x_{345}\end{aligned}$\\[3pt]
$S_{283}$ & $\begin{aligned}[t]&1+x_{215}x_{345}-x_{341}^{2}\end{aligned}$\\[3pt]
$S_{284}$ & $\begin{aligned}[t]&-x_{344}^{2}+x_{346}\end{aligned}$\\[3pt]
$S_{285}$ & $\begin{aligned}[t]&1+x_{200}x_{346}-x_{343}^{2}-x_{346}\end{aligned}$\\[3pt]
$S_{286}$ & $\begin{aligned}[t]&-x_{188}x_{347}-x_{341}+x_{343}\end{aligned}$\\[3pt]
$S_{287}$ & $\begin{aligned}[t]&-1-x_{35}-4x_{342}x_{348}+x_{344}\end{aligned}$\\[3pt]
$S_{288}$ & $\begin{aligned}[t]&-x_{73}x_{178}+x_{350}\end{aligned}$\\[3pt]
$S_{289}$ & $\begin{aligned}[t]&-x_{62}x_{342}+x_{181}x_{342}+x_{182}x_{342}-x_{265}x_{342}-x_{341}-2x_{342}+x_{349}x_{350}+x_{642}^{2}+x_{643}^{2}+x_{644}+x_{644}^{2}\end{aligned}$\\[3pt]
$S_{290}$ & $\begin{aligned}[t]&-x_{352}^{2}+x_{355}\end{aligned}$\\[3pt]
$S_{291}$ & $\begin{aligned}[t]&1+x_{215}x_{355}-x_{351}^{2}\end{aligned}$\\[3pt]
$S_{292}$ & $\begin{aligned}[t]&-x_{354}^{2}+x_{356}\end{aligned}$\\[3pt]
$S_{293}$ & $\begin{aligned}[t]&1+x_{200}x_{356}-x_{353}^{2}-x_{356}\end{aligned}$\\[3pt]
$S_{294}$ & $\begin{aligned}[t]&-x_{188}x_{357}-x_{351}+x_{353}\end{aligned}$\\[3pt]
$S_{295}$ & $\begin{aligned}[t]&-1-x_{36}-4x_{352}x_{358}+x_{354}\end{aligned}$\\[3pt]
$S_{296}$ & $\begin{aligned}[t]&-x_{62}x_{352}+x_{181}x_{352}+x_{182}x_{352}-x_{265}x_{352}+x_{350}x_{359}-x_{351}-2x_{352}+x_{654}^{2}+x_{655}^{2}+x_{656}+x_{656}^{2}\end{aligned}$\\[3pt]
$S_{297}$ & $\begin{aligned}[t]&-x_{198}x_{211}+x_{360}\end{aligned}$\\[3pt]
$S_{298}$ & $\begin{aligned}[t]&-x_{242}x_{360}+x_{361}\end{aligned}$\\[3pt]
$S_{299}$ & $\begin{aligned}[t]&-x_{251}x_{361}+x_{362}\end{aligned}$\\[3pt]
$S_{300}$ & $\begin{aligned}[t]&-x_{260}x_{362}+x_{363}\end{aligned}$\\[3pt]
$S_{301}$ & $\begin{aligned}[t]&-x_{270}x_{363}+x_{364}\end{aligned}$\\[3pt]
$S_{302}$ & $\begin{aligned}[t]&-x_{280}x_{364}+x_{365}\end{aligned}$\\[3pt]
$S_{303}$ & $\begin{aligned}[t]&-x_{289}x_{365}+x_{366}\end{aligned}$\\[3pt]
$S_{304}$ & $\begin{aligned}[t]&-x_{299}x_{366}+x_{367}\end{aligned}$\\[3pt]
$S_{305}$ & $\begin{aligned}[t]&-x_{308}x_{367}+x_{368}\end{aligned}$\\[3pt]
$S_{306}$ & $\begin{aligned}[t]&-x_{317}x_{368}+x_{369}\end{aligned}$\\[3pt]
$S_{307}$ & $\begin{aligned}[t]&-x_{326}x_{369}+x_{370}\end{aligned}$\\[3pt]
$S_{308}$ & $\begin{aligned}[t]&-x_{336}x_{370}+x_{371}\end{aligned}$\\[3pt]
$S_{309}$ & $\begin{aligned}[t]&-x_{345}x_{371}+x_{372}\end{aligned}$\\[3pt]
$S_{310}$ & $\begin{aligned}[t]&-x_{355}x_{372}+x_{373}\end{aligned}$\\[3pt]
$S_{311}$ & $\begin{aligned}[t]&x_{189}-x_{373}-x_{373}x_{374}\end{aligned}$\\[3pt]
$S_{312}$ & $\begin{aligned}[t]&-16-3x_{0}-x_{1}-x_{2}-x_{4}-x_{60}-2x_{61}-x_{62}-x_{63}-x_{64}-x_{65}-x_{66}-2x_{68}-x_{174}-x_{175}-x_{176}-x_{177}-x_{178}\\ &+x_{180}-2x_{207}+2x_{210}-x_{265}-x_{441}^{2}-x_{442}^{2}-x_{443}-x_{443}^{2}\end{aligned}$\\[3pt]
$S_{313}$ & $\begin{aligned}[t]&x_{0}-x_{375}^{2}-x_{376}^{2}-x_{377}-x_{377}^{2}\end{aligned}$\\[3pt]
$S_{314}$ & $\begin{aligned}[t]&x_{3}-x_{378}^{2}-x_{379}^{2}-x_{380}-x_{380}^{2}\end{aligned}$\\[3pt]
$S_{315}$ & $\begin{aligned}[t]&x_{4}-x_{381}^{2}-x_{382}^{2}-x_{383}-x_{383}^{2}\end{aligned}$\\[3pt]
$S_{316}$ & $\begin{aligned}[t]&-2+x_{2}-x_{3}-x_{384}^{2}-x_{385}^{2}-x_{386}-x_{386}^{2}\end{aligned}$\\[3pt]
$S_{317}$ & $\begin{aligned}[t]&x_{27}-x_{387}^{2}-x_{388}^{2}-x_{389}-x_{389}^{2}\end{aligned}$\\[3pt]
$S_{318}$ & $\begin{aligned}[t]&x_{29}-x_{390}^{2}-x_{391}^{2}-x_{392}-x_{392}^{2}\end{aligned}$\\[3pt]
$S_{319}$ & $\begin{aligned}[t]&x_{30}-x_{393}^{2}-x_{394}^{2}-x_{395}-x_{395}^{2}\end{aligned}$\\[3pt]
$S_{320}$ & $\begin{aligned}[t]&x_{31}-x_{396}^{2}-x_{397}^{2}-x_{398}-x_{398}^{2}\end{aligned}$\\[3pt]
$S_{321}$ & $\begin{aligned}[t]&x_{32}-x_{399}^{2}-x_{400}^{2}-x_{401}-x_{401}^{2}\end{aligned}$\\[3pt]
$S_{322}$ & $\begin{aligned}[t]&x_{33}-x_{402}^{2}-x_{403}^{2}-x_{404}-x_{404}^{2}\end{aligned}$\\[3pt]
$S_{323}$ & $\begin{aligned}[t]&x_{34}-x_{405}^{2}-x_{406}^{2}-x_{407}-x_{407}^{2}\end{aligned}$\\[3pt]
$S_{324}$ & $\begin{aligned}[t]&x_{35}-x_{408}^{2}-x_{409}^{2}-x_{410}-x_{410}^{2}\end{aligned}$\\[3pt]
$S_{325}$ & $\begin{aligned}[t]&x_{36}-x_{411}^{2}-x_{412}^{2}-x_{413}-x_{413}^{2}\end{aligned}$\\[3pt]
$S_{326}$ & $\begin{aligned}[t]&x_{37}-x_{414}^{2}-x_{415}^{2}-x_{416}-x_{416}^{2}\end{aligned}$\\[3pt]
$S_{327}$ & $\begin{aligned}[t]&x_{61}-x_{417}^{2}-x_{418}^{2}-x_{419}-x_{419}^{2}\end{aligned}$\\[3pt]
$S_{328}$ & $\begin{aligned}[t]&x_{63}-x_{420}^{2}-x_{421}^{2}-x_{422}-x_{422}^{2}\end{aligned}$\\[3pt]
$S_{329}$ & $\begin{aligned}[t]&-1-x_{60}x_{62}+x_{62}+x_{63}-x_{423}^{2}-x_{424}^{2}-x_{425}-x_{425}^{2}\end{aligned}$\\[3pt]
$S_{330}$ & $\begin{aligned}[t]&x_{60}x_{62}-x_{63}-x_{426}^{2}-x_{427}^{2}-x_{428}-x_{428}^{2}\end{aligned}$\\[3pt]
$S_{331}$ & $\begin{aligned}[t]&x_{65}-x_{429}^{2}-x_{430}^{2}-x_{431}-x_{431}^{2}\end{aligned}$\\[3pt]
$S_{332}$ & $\begin{aligned}[t]&x_{66}-x_{432}^{2}-x_{433}^{2}-x_{434}-x_{434}^{2}\end{aligned}$\\[3pt]
$S_{333}$ & $\begin{aligned}[t]&x_{67}-x_{435}^{2}-x_{436}^{2}-x_{437}-x_{437}^{2}\end{aligned}$\\[3pt]
$S_{334}$ & $\begin{aligned}[t]&x_{68}-x_{438}^{2}-x_{439}^{2}-x_{440}-x_{440}^{2}\end{aligned}$\\[3pt]
$S_{335}$ & $\begin{aligned}[t]&-1-x_{67}x_{69}+x_{69}+x_{441}^{2}+x_{442}^{2}+x_{443}+x_{443}^{2}-x_{444}^{2}-x_{445}^{2}-x_{446}-x_{446}^{2}\end{aligned}$\\[3pt]
$S_{336}$ & $\begin{aligned}[t]&x_{67}x_{69}-x_{441}^{2}-x_{442}^{2}-x_{443}-x_{443}^{2}-x_{447}^{2}-x_{448}^{2}-x_{449}-x_{449}^{2}\end{aligned}$\\[3pt]
$S_{337}$ & $\begin{aligned}[t]&x_{70}-x_{450}^{2}-x_{451}^{2}-x_{452}-x_{452}^{2}\end{aligned}$\\[3pt]
$S_{338}$ & $\begin{aligned}[t]&-1+x_{72}-x_{73}-x_{453}^{2}-x_{454}^{2}-x_{455}-x_{455}^{2}\end{aligned}$\\[3pt]
$S_{339}$ & $\begin{aligned}[t]&-8-x_{29}-x_{30}-x_{31}-x_{32}-x_{33}-x_{34}-x_{35}-x_{36}+8x_{73}-x_{456}^{2}-x_{457}^{2}-x_{458}-x_{458}^{2}\end{aligned}$\\[3pt]
$S_{340}$ & $\begin{aligned}[t]&x_{181}-x_{459}^{2}-x_{460}^{2}-x_{461}-x_{461}^{2}\end{aligned}$\\[3pt]
$S_{341}$ & $\begin{aligned}[t]&x_{182}-x_{462}^{2}-x_{463}^{2}-x_{464}-x_{464}^{2}\end{aligned}$\\[3pt]
$S_{342}$ & $\begin{aligned}[t]&x_{183}-x_{465}^{2}-x_{466}^{2}-x_{467}-x_{467}^{2}\end{aligned}$\\[3pt]
$S_{343}$ & $\begin{aligned}[t]&x_{190}-x_{468}^{2}-x_{469}^{2}-x_{470}-x_{470}^{2}\end{aligned}$\\[3pt]
$S_{344}$ & $\begin{aligned}[t]&x_{194}-x_{471}^{2}-x_{472}^{2}-x_{473}-x_{473}^{2}\end{aligned}$\\[3pt]
$S_{345}$ & $\begin{aligned}[t]&-1-x_{0}+x_{195}-x_{474}^{2}-x_{475}^{2}-x_{476}-x_{476}^{2}\end{aligned}$\\[3pt]
$S_{346}$ & $\begin{aligned}[t]&x_{203}-x_{477}^{2}-x_{478}^{2}-x_{479}-x_{479}^{2}\end{aligned}$\\[3pt]
$S_{347}$ & $\begin{aligned}[t]&x_{206}-x_{480}^{2}-x_{481}^{2}-x_{482}-x_{482}^{2}\end{aligned}$\\[3pt]
$S_{348}$ & $\begin{aligned}[t]&x_{210}-x_{483}^{2}-x_{484}^{2}-x_{485}-x_{485}^{2}\end{aligned}$\\[3pt]
$S_{349}$ & $\begin{aligned}[t]&x_{219}-x_{486}^{2}-x_{487}^{2}-x_{488}-x_{488}^{2}\end{aligned}$\\[3pt]
$S_{350}$ & $\begin{aligned}[t]&x_{221}-x_{489}^{2}-x_{490}^{2}-x_{491}-x_{491}^{2}\end{aligned}$\\[3pt]
$S_{351}$ & $\begin{aligned}[t]&x_{223}-x_{492}^{2}-x_{493}^{2}-x_{494}-x_{494}^{2}\end{aligned}$\\[3pt]
$S_{352}$ & $\begin{aligned}[t]&x_{225}-x_{495}^{2}-x_{496}^{2}-x_{497}-x_{497}^{2}\end{aligned}$\\[3pt]
$S_{353}$ & $\begin{aligned}[t]&x_{230}-x_{498}^{2}-x_{499}^{2}-x_{500}-x_{500}^{2}\end{aligned}$\\[3pt]
$S_{354}$ & $\begin{aligned}[t]&x_{232}-x_{501}^{2}-x_{502}^{2}-x_{503}-x_{503}^{2}\end{aligned}$\\[3pt]
$S_{355}$ & $\begin{aligned}[t]&x_{234}-x_{504}^{2}-x_{505}^{2}-x_{506}-x_{506}^{2}\end{aligned}$\\[3pt]
$S_{356}$ & $\begin{aligned}[t]&x_{238}-x_{507}^{2}-x_{508}^{2}-x_{509}-x_{509}^{2}\end{aligned}$\\[3pt]
$S_{357}$ & $\begin{aligned}[t]&-1-x_{2}+x_{239}-x_{510}^{2}-x_{511}^{2}-x_{512}-x_{512}^{2}\end{aligned}$\\[3pt]
$S_{358}$ & $\begin{aligned}[t]&x_{247}-x_{513}^{2}-x_{514}^{2}-x_{515}-x_{515}^{2}\end{aligned}$\\[3pt]
$S_{359}$ & $\begin{aligned}[t]&-1-x_{1}+x_{248}-x_{516}^{2}-x_{517}^{2}-x_{518}-x_{518}^{2}\end{aligned}$\\[3pt]
$S_{360}$ & $\begin{aligned}[t]&x_{256}-x_{519}^{2}-x_{520}^{2}-x_{521}-x_{521}^{2}\end{aligned}$\\[3pt]
$S_{361}$ & $\begin{aligned}[t]&-x_{61}-x_{207}+x_{210}+x_{257}-x_{522}^{2}-x_{523}^{2}-x_{524}-x_{524}^{2}\end{aligned}$\\[3pt]
$S_{362}$ & $\begin{aligned}[t]&-1+x_{174}-x_{525}^{2}-x_{526}^{2}-x_{527}-x_{527}^{2}-x_{528}^{2}-x_{529}^{2}-x_{530}-x_{530}^{2}\end{aligned}$\\[3pt]
$S_{363}$ & $\begin{aligned}[t]&x_{265}-x_{531}^{2}-x_{532}^{2}-x_{533}-x_{533}^{2}\end{aligned}$\\[3pt]
$S_{364}$ & $\begin{aligned}[t]&x_{266}-x_{534}^{2}-x_{535}^{2}-x_{536}-x_{536}^{2}\end{aligned}$\\[3pt]
$S_{365}$ & $\begin{aligned}[t]&-1-x_{0}-x_{68}+x_{267}-x_{537}^{2}-x_{538}^{2}-x_{539}-x_{539}^{2}\end{aligned}$\\[3pt]
$S_{366}$ & $\begin{aligned}[t]&-1+x_{175}-x_{540}^{2}-x_{541}^{2}-x_{542}-x_{542}^{2}-x_{543}^{2}-x_{544}^{2}-x_{545}-x_{545}^{2}\end{aligned}$\\[3pt]
$S_{367}$ & $\begin{aligned}[t]&x_{62}-x_{69}+x_{265}-x_{546}^{2}-x_{547}^{2}-x_{548}-x_{548}^{2}\end{aligned}$\\[3pt]
$S_{368}$ & $\begin{aligned}[t]&x_{276}-x_{549}^{2}-x_{550}^{2}-x_{551}-x_{551}^{2}\end{aligned}$\\[3pt]
$S_{369}$ & $\begin{aligned}[t]&-3-x_{0}-x_{65}+x_{277}-x_{552}^{2}-x_{553}^{2}-x_{554}-x_{554}^{2}\end{aligned}$\\[3pt]
$S_{370}$ & $\begin{aligned}[t]&-1+x_{175}-x_{555}^{2}-x_{556}^{2}-x_{557}-x_{557}^{2}-x_{558}^{2}-x_{559}^{2}-x_{560}-x_{560}^{2}\end{aligned}$\\[3pt]
$S_{371}$ & $\begin{aligned}[t]&x_{62}-x_{70}+x_{265}-x_{561}^{2}-x_{562}^{2}-x_{563}-x_{563}^{2}\end{aligned}$\\[3pt]
$S_{372}$ & $\begin{aligned}[t]&x_{285}-x_{564}^{2}-x_{565}^{2}-x_{566}-x_{566}^{2}\end{aligned}$\\[3pt]
$S_{373}$ & $\begin{aligned}[t]&-1-x_{29}+x_{286}-x_{567}^{2}-x_{568}^{2}-x_{569}-x_{569}^{2}\end{aligned}$\\[3pt]
$S_{374}$ & $\begin{aligned}[t]&-1+x_{176}-x_{570}^{2}-x_{571}^{2}-x_{572}-x_{572}^{2}-x_{573}^{2}-x_{574}^{2}-x_{575}-x_{575}^{2}\end{aligned}$\\[3pt]
$S_{375}$ & $\begin{aligned}[t]&x_{295}-x_{576}^{2}-x_{577}^{2}-x_{578}-x_{578}^{2}\end{aligned}$\\[3pt]
$S_{376}$ & $\begin{aligned}[t]&-1-x_{30}+x_{296}-x_{579}^{2}-x_{580}^{2}-x_{581}-x_{581}^{2}\end{aligned}$\\[3pt]
$S_{377}$ & $\begin{aligned}[t]&-1+x_{176}-x_{582}^{2}-x_{583}^{2}-x_{584}-x_{584}^{2}-x_{585}^{2}-x_{586}^{2}-x_{587}-x_{587}^{2}\end{aligned}$\\[3pt]
$S_{378}$ & $\begin{aligned}[t]&x_{304}-x_{588}^{2}-x_{589}^{2}-x_{590}-x_{590}^{2}\end{aligned}$\\[3pt]
$S_{379}$ & $\begin{aligned}[t]&-1-x_{31}+x_{305}-x_{591}^{2}-x_{592}^{2}-x_{593}-x_{593}^{2}\end{aligned}$\\[3pt]
$S_{380}$ & $\begin{aligned}[t]&-1+x_{176}-x_{594}^{2}-x_{595}^{2}-x_{596}-x_{596}^{2}-x_{597}^{2}-x_{598}^{2}-x_{599}-x_{599}^{2}\end{aligned}$\\[3pt]
$S_{381}$ & $\begin{aligned}[t]&x_{313}-x_{600}^{2}-x_{601}^{2}-x_{602}-x_{602}^{2}\end{aligned}$\\[3pt]
$S_{382}$ & $\begin{aligned}[t]&-1-x_{32}+x_{314}-x_{603}^{2}-x_{604}^{2}-x_{605}-x_{605}^{2}\end{aligned}$\\[3pt]
$S_{383}$ & $\begin{aligned}[t]&-1+x_{176}-x_{606}^{2}-x_{607}^{2}-x_{608}-x_{608}^{2}-x_{609}^{2}-x_{610}^{2}-x_{611}-x_{611}^{2}\end{aligned}$\\[3pt]
$S_{384}$ & $\begin{aligned}[t]&x_{322}-x_{612}^{2}-x_{613}^{2}-x_{614}-x_{614}^{2}\end{aligned}$\\[3pt]
$S_{385}$ & $\begin{aligned}[t]&-1-x_{33}+x_{323}-x_{615}^{2}-x_{616}^{2}-x_{617}-x_{617}^{2}\end{aligned}$\\[3pt]
$S_{386}$ & $\begin{aligned}[t]&-1+x_{177}-x_{618}^{2}-x_{619}^{2}-x_{620}-x_{620}^{2}-x_{621}^{2}-x_{622}^{2}-x_{623}-x_{623}^{2}\end{aligned}$\\[3pt]
$S_{387}$ & $\begin{aligned}[t]&x_{332}-x_{624}^{2}-x_{625}^{2}-x_{626}-x_{626}^{2}\end{aligned}$\\[3pt]
$S_{388}$ & $\begin{aligned}[t]&-1-x_{34}+x_{333}-x_{627}^{2}-x_{628}^{2}-x_{629}-x_{629}^{2}\end{aligned}$\\[3pt]
$S_{389}$ & $\begin{aligned}[t]&-1+x_{177}-x_{630}^{2}-x_{631}^{2}-x_{632}-x_{632}^{2}-x_{633}^{2}-x_{634}^{2}-x_{635}-x_{635}^{2}\end{aligned}$\\[3pt]
$S_{390}$ & $\begin{aligned}[t]&x_{341}-x_{636}^{2}-x_{637}^{2}-x_{638}-x_{638}^{2}\end{aligned}$\\[3pt]
$S_{391}$ & $\begin{aligned}[t]&-1-x_{35}+x_{342}-x_{639}^{2}-x_{640}^{2}-x_{641}-x_{641}^{2}\end{aligned}$\\[3pt]
$S_{392}$ & $\begin{aligned}[t]&-1+x_{178}-x_{642}^{2}-x_{643}^{2}-x_{644}-x_{644}^{2}-x_{645}^{2}-x_{646}^{2}-x_{647}-x_{647}^{2}\end{aligned}$\\[3pt]
$S_{393}$ & $\begin{aligned}[t]&x_{351}-x_{648}^{2}-x_{649}^{2}-x_{650}-x_{650}^{2}\end{aligned}$\\[3pt]
$S_{394}$ & $\begin{aligned}[t]&-1-x_{36}+x_{352}-x_{651}^{2}-x_{652}^{2}-x_{653}-x_{653}^{2}\end{aligned}$\\[3pt]
$S_{395}$ & $\begin{aligned}[t]&-1+x_{178}-x_{654}^{2}-x_{655}^{2}-x_{656}-x_{656}^{2}-x_{657}^{2}-x_{658}^{2}-x_{659}-x_{659}^{2}\end{aligned}$\\[3pt]
$S_{396}$ & $\begin{aligned}[t]&x_{374}-x_{660}^{2}-x_{661}^{2}-x_{662}-x_{662}^{2}\end{aligned}$\\[3pt]
\end{longtable}
\endgroup


\clearpage
\section{The complete \svars-variable small-coefficient equation}\label{app:small}

Again each coordinate ranges independently over $\Z$. The equation is
\[
 \boxed{\PS(x_0,\ldots,x_{\svarlast})=
       \sum_{j=0}^{\slast}T_j(x_0,\ldots,x_{\svarlast})^2=0.}
\]
Its fully collected coefficient height is \sheight. It trades a few
additional coordinates for fewer expanded terms and much smaller
coefficients than the principal integer variant.
\begingroup\footnotesize
\setlength{\tabcolsep}{3pt}
\begin{longtable}{r@{\hspace{0.8em}}l}
\toprule
Residual & Polynomial (set each to zero)\\
\midrule
\endfirsthead
\toprule
Residual & Polynomial (continued)\\
\midrule
\endhead
\bottomrule\endfoot
$T_{0}$ & $\begin{aligned}[t]&-x_{4}^{2}+x_{5}\end{aligned}$\\[3pt]
$T_{1}$ & $\begin{aligned}[t]&-x_{5}^{2}+x_{6}\end{aligned}$\\[3pt]
$T_{2}$ & $\begin{aligned}[t]&-x_{6}^{2}+x_{7}\end{aligned}$\\[3pt]
$T_{3}$ & $\begin{aligned}[t]&-x_{7}^{2}+x_{8}\end{aligned}$\\[3pt]
$T_{4}$ & $\begin{aligned}[t]&-x_{8}^{2}+x_{9}\end{aligned}$\\[3pt]
$T_{5}$ & $\begin{aligned}[t]&-x_{9}^{2}+x_{10}\end{aligned}$\\[3pt]
$T_{6}$ & $\begin{aligned}[t]&-x_{10}^{2}+x_{11}\end{aligned}$\\[3pt]
$T_{7}$ & $\begin{aligned}[t]&-x_{11}^{2}+x_{12}\end{aligned}$\\[3pt]
$T_{8}$ & $\begin{aligned}[t]&-x_{12}^{2}+x_{13}\end{aligned}$\\[3pt]
$T_{9}$ & $\begin{aligned}[t]&-x_{13}^{2}+x_{14}\end{aligned}$\\[3pt]
$T_{10}$ & $\begin{aligned}[t]&-x_{14}^{2}+x_{15}\end{aligned}$\\[3pt]
$T_{11}$ & $\begin{aligned}[t]&-x_{15}^{2}+x_{16}\end{aligned}$\\[3pt]
$T_{12}$ & $\begin{aligned}[t]&-x_{16}^{2}+x_{17}\end{aligned}$\\[3pt]
$T_{13}$ & $\begin{aligned}[t]&-x_{17}^{2}+x_{18}\end{aligned}$\\[3pt]
$T_{14}$ & $\begin{aligned}[t]&-x_{18}^{2}+x_{19}\end{aligned}$\\[3pt]
$T_{15}$ & $\begin{aligned}[t]&-x_{19}^{2}+x_{20}\end{aligned}$\\[3pt]
$T_{16}$ & $\begin{aligned}[t]&-x_{20}^{2}+x_{21}\end{aligned}$\\[3pt]
$T_{17}$ & $\begin{aligned}[t]&-x_{21}^{2}+x_{22}\end{aligned}$\\[3pt]
$T_{18}$ & $\begin{aligned}[t]&-x_{22}^{2}+x_{23}\end{aligned}$\\[3pt]
$T_{19}$ & $\begin{aligned}[t]&-x_{23}^{2}+x_{24}\end{aligned}$\\[3pt]
$T_{20}$ & $\begin{aligned}[t]&-x_{24}^{2}+x_{25}\end{aligned}$\\[3pt]
$T_{21}$ & $\begin{aligned}[t]&1+x_{3}-x_{25}^{2}\end{aligned}$\\[3pt]
$T_{22}$ & $\begin{aligned}[t]&-x_{0}x_{2}-2x_{2}+x_{28}\end{aligned}$\\[3pt]
$T_{23}$ & $\begin{aligned}[t]&-2+x_{1}-x_{2}x_{26}-x_{3}-x_{26}\end{aligned}$\\[3pt]
$T_{24}$ & $\begin{aligned}[t]&-1+x_{1}-x_{27}-x_{27}x_{28}\end{aligned}$\\[3pt]
$T_{25}$ & $\begin{aligned}[t]&-3-x_{0}-x_{1}-x_{2}-x_{3}+x_{38}\end{aligned}$\\[3pt]
$T_{26}$ & $\begin{aligned}[t]&-x_{38}^{2}+x_{39}\end{aligned}$\\[3pt]
$T_{27}$ & $\begin{aligned}[t]&-x_{39}^{2}+x_{40}\end{aligned}$\\[3pt]
$T_{28}$ & $\begin{aligned}[t]&-x_{40}^{2}+x_{41}\end{aligned}$\\[3pt]
$T_{29}$ & $\begin{aligned}[t]&-x_{41}^{2}+x_{42}\end{aligned}$\\[3pt]
$T_{30}$ & $\begin{aligned}[t]&-x_{42}^{2}+x_{43}\end{aligned}$\\[3pt]
$T_{31}$ & $\begin{aligned}[t]&-x_{43}^{2}+x_{44}\end{aligned}$\\[3pt]
$T_{32}$ & $\begin{aligned}[t]&-x_{44}^{2}+x_{45}\end{aligned}$\\[3pt]
$T_{33}$ & $\begin{aligned}[t]&-x_{45}^{2}+x_{46}\end{aligned}$\\[3pt]
$T_{34}$ & $\begin{aligned}[t]&-x_{46}^{2}+x_{47}\end{aligned}$\\[3pt]
$T_{35}$ & $\begin{aligned}[t]&-x_{47}^{2}+x_{48}\end{aligned}$\\[3pt]
$T_{36}$ & $\begin{aligned}[t]&-x_{48}^{2}+x_{49}\end{aligned}$\\[3pt]
$T_{37}$ & $\begin{aligned}[t]&-x_{49}^{2}+x_{50}\end{aligned}$\\[3pt]
$T_{38}$ & $\begin{aligned}[t]&-x_{50}^{2}+x_{51}\end{aligned}$\\[3pt]
$T_{39}$ & $\begin{aligned}[t]&-x_{51}^{2}+x_{52}\end{aligned}$\\[3pt]
$T_{40}$ & $\begin{aligned}[t]&-x_{52}^{2}+x_{53}\end{aligned}$\\[3pt]
$T_{41}$ & $\begin{aligned}[t]&-x_{53}^{2}+x_{54}\end{aligned}$\\[3pt]
$T_{42}$ & $\begin{aligned}[t]&-x_{54}^{2}+x_{55}\end{aligned}$\\[3pt]
$T_{43}$ & $\begin{aligned}[t]&-x_{55}^{2}+x_{56}\end{aligned}$\\[3pt]
$T_{44}$ & $\begin{aligned}[t]&-x_{56}^{2}+x_{57}\end{aligned}$\\[3pt]
$T_{45}$ & $\begin{aligned}[t]&-x_{57}^{2}+x_{58}\end{aligned}$\\[3pt]
$T_{46}$ & $\begin{aligned}[t]&-x_{58}^{2}+x_{59}\end{aligned}$\\[3pt]
$T_{47}$ & $\begin{aligned}[t]&-x_{59}^{2}+x_{60}\end{aligned}$\\[3pt]
$T_{48}$ & $\begin{aligned}[t]&-x_{60}^{2}+x_{222}-x_{225}\end{aligned}$\\[3pt]
$T_{49}$ & $\begin{aligned}[t]&-x_{0}x_{61}-2x_{61}+x_{65}\end{aligned}$\\[3pt]
$T_{50}$ & $\begin{aligned}[t]&-x_{67}x_{68}+x_{72}\end{aligned}$\\[3pt]
$T_{51}$ & $\begin{aligned}[t]&-x_{71}x_{72}+x_{73}x_{75}+x_{74}\end{aligned}$\\[3pt]
$T_{52}$ & $\begin{aligned}[t]&1-2x_{61}-x_{61}x_{66}+x_{73}\end{aligned}$\\[3pt]
$T_{53}$ & $\begin{aligned}[t]&-1-x_{37}x_{61}-x_{61}+x_{74}\end{aligned}$\\[3pt]
$T_{54}$ & $\begin{aligned}[t]&-6-5x_{29}-x_{29}^{2}-2x_{29}x_{30}-5x_{30}-x_{30}^{2}+x_{76}\end{aligned}$\\[3pt]
$T_{55}$ & $\begin{aligned}[t]&-72-68x_{30}-16x_{30}^{2}-16x_{30}x_{76}-34x_{76}-4x_{76}^{2}+x_{77}\end{aligned}$\\[3pt]
$T_{56}$ & $\begin{aligned}[t]&-30-11x_{29}-x_{29}^{2}-4x_{29}x_{77}-22x_{77}-4x_{77}^{2}+x_{78}\end{aligned}$\\[3pt]
$T_{57}$ & $\begin{aligned}[t]&-380-312x_{77}-64x_{77}^{2}-32x_{77}x_{78}-78x_{78}-4x_{78}^{2}+x_{79}\end{aligned}$\\[3pt]
$T_{58}$ & $\begin{aligned}[t]&-6-5x_{30}-x_{30}^{2}-2x_{30}x_{31}-5x_{31}-x_{31}^{2}+x_{80}\end{aligned}$\\[3pt]
$T_{59}$ & $\begin{aligned}[t]&-30-11x_{29}-x_{29}^{2}-8x_{29}x_{31}-4x_{29}x_{80}-44x_{31}-16x_{31}^{2}-16x_{31}x_{80}-22x_{80}-4x_{80}^{2}+x_{81}\end{aligned}$\\[3pt]
$T_{60}$ & $\begin{aligned}[t]&-6-5x_{29}-x_{29}^{2}-2x_{29}x_{31}-5x_{31}-x_{31}^{2}+x_{82}\end{aligned}$\\[3pt]
$T_{61}$ & $\begin{aligned}[t]&-12-4x_{30}-4x_{31}-2x_{76}-2x_{82}+x_{83}\end{aligned}$\\[3pt]
$T_{62}$ & $\begin{aligned}[t]&-42-32x_{31}-8x_{80}-2x_{81}-8x_{82}-2x_{83}-2x_{83}^{2}+x_{84}\end{aligned}$\\[3pt]
$T_{63}$ & $\begin{aligned}[t]&-96-64x_{31}-32x_{82}-8x_{83}-8x_{83}^{2}-2x_{84}-2x_{84}^{2}+x_{85}\end{aligned}$\\[3pt]
$T_{64}$ & $\begin{aligned}[t]&-10-4x_{29}-4x_{30}-2x_{80}-2x_{82}+x_{86}\end{aligned}$\\[3pt]
$T_{65}$ & $\begin{aligned}[t]&-27-20x_{30}-2x_{76}-8x_{80}-2x_{86}-2x_{86}^{2}+x_{87}\end{aligned}$\\[3pt]
$T_{66}$ & $\begin{aligned}[t]&-80-64x_{30}-32x_{80}-8x_{86}-8x_{86}^{2}-2x_{87}-2x_{87}^{2}+x_{88}\end{aligned}$\\[3pt]
$T_{67}$ & $\begin{aligned}[t]&-56-15x_{29}-x_{29}^{2}-8x_{29}x_{31}-4x_{29}x_{80}-60x_{31}-16x_{31}^{2}-16x_{31}x_{80}-30x_{80}-4x_{80}^{2}+x_{89}\end{aligned}$\\[3pt]
$T_{68}$ & $\begin{aligned}[t]&-56-15x_{30}-x_{30}^{2}-8x_{30}x_{31}-4x_{30}x_{82}-60x_{31}-16x_{31}^{2}-16x_{31}x_{82}-30x_{82}-4x_{82}^{2}+x_{90}\end{aligned}$\\[3pt]
$T_{69}$ & $\begin{aligned}[t]&-52-32x_{31}-8x_{80}-8x_{82}-2x_{89}-2x_{90}+x_{91}\end{aligned}$\\[3pt]
$T_{70}$ & $\begin{aligned}[t]&-104-64x_{31}-32x_{82}-8x_{90}-2x_{91}-2x_{91}^{2}+x_{92}\end{aligned}$\\[3pt]
$T_{71}$ & $\begin{aligned}[t]&-110-168x_{30}-64x_{30}^{2}-32x_{30}x_{76}-32x_{30}x_{80}-42x_{76}-4x_{76}^{2}-8x_{76}x_{80}-42x_{80}-4x_{80}^{2}+x_{93}\end{aligned}$\\[3pt]
$T_{72}$ & $\begin{aligned}[t]&-27-4x_{29}-16x_{30}-8x_{76}-2x_{82}-2x_{93}+x_{94}\end{aligned}$\\[3pt]
$T_{73}$ & $\begin{aligned}[t]&-80-64x_{30}-32x_{76}-8x_{93}-2x_{94}-2x_{94}^{2}+x_{95}\end{aligned}$\\[3pt]
$T_{74}$ & $\begin{aligned}[t]&-56-60x_{30}-16x_{30}^{2}-16x_{30}x_{76}-30x_{76}-4x_{76}^{2}+x_{96}\end{aligned}$\\[3pt]
$T_{75}$ & $\begin{aligned}[t]&-30-11x_{29}-x_{29}^{2}-4x_{29}x_{96}-22x_{96}-4x_{96}^{2}+x_{97}\end{aligned}$\\[3pt]
$T_{76}$ & $\begin{aligned}[t]&-380-312x_{96}-64x_{96}^{2}-32x_{96}x_{97}-78x_{97}-4x_{97}^{2}+x_{98}\end{aligned}$\\[3pt]
$T_{77}$ & $\begin{aligned}[t]&-30-11x_{30}-x_{30}^{2}-4x_{30}x_{98}-22x_{98}-4x_{98}^{2}+x_{99}\end{aligned}$\\[3pt]
$T_{78}$ & $\begin{aligned}[t]&-156-200x_{30}-64x_{30}^{2}-64x_{30}x_{76}-100x_{76}-16x_{76}^{2}+x_{100}\end{aligned}$\\[3pt]
$T_{79}$ & $\begin{aligned}[t]&-37-20x_{30}-8x_{76}-2x_{80}-2x_{100}+x_{101}\end{aligned}$\\[3pt]
$T_{80}$ & $\begin{aligned}[t]&-120-64x_{30}-32x_{76}-8x_{100}-2x_{101}-2x_{101}^{2}+x_{102}\end{aligned}$\\[3pt]
$T_{81}$ & $\begin{aligned}[t]&-56-60x_{29}-16x_{29}^{2}-16x_{29}x_{82}-8x_{29}x_{102}-30x_{82}-4x_{82}^{2}-4x_{82}x_{102}-15x_{102}-x_{102}^{2}+x_{103}\end{aligned}$\\[3pt]
$T_{82}$ & $\begin{aligned}[t]&-6-5x_{31}-x_{31}^{2}-2x_{31}x_{32}-5x_{32}-x_{32}^{2}+x_{104}\end{aligned}$\\[3pt]
$T_{83}$ & $\begin{aligned}[t]&-30-11x_{29}-x_{29}^{2}-8x_{29}x_{32}-4x_{29}x_{104}-44x_{32}-16x_{32}^{2}-16x_{32}x_{104}-22x_{104}-4x_{104}^{2}+x_{105}\end{aligned}$\\[3pt]
$T_{84}$ & $\begin{aligned}[t]&-72-136x_{31}-64x_{31}^{2}-32x_{31}x_{80}-32x_{31}x_{82}-34x_{80}-4x_{80}^{2}-8x_{80}x_{82}-34x_{82}-4x_{82}^{2}+x_{106}\end{aligned}$\\[3pt]
$T_{85}$ & $\begin{aligned}[t]&-20-4x_{29}-16x_{31}-8x_{80}-2x_{105}-2x_{106}+x_{107}\end{aligned}$\\[3pt]
$T_{86}$ & $\begin{aligned}[t]&-64-64x_{31}-32x_{80}-8x_{106}-2x_{107}-2x_{107}^{2}+x_{108}\end{aligned}$\\[3pt]
$T_{87}$ & $\begin{aligned}[t]&-21-16x_{30}-x_{31}-8x_{80}-2x_{86}-2x_{86}^{2}+x_{109}\end{aligned}$\\[3pt]
$T_{88}$ & $\begin{aligned}[t]&-80-64x_{30}-32x_{80}-8x_{86}-8x_{86}^{2}-2x_{109}-2x_{109}^{2}+x_{110}\end{aligned}$\\[3pt]
$T_{89}$ & $\begin{aligned}[t]&-42-52x_{31}-16x_{31}^{2}-16x_{31}x_{80}-26x_{80}-4x_{80}^{2}+x_{111}\end{aligned}$\\[3pt]
$T_{90}$ & $\begin{aligned}[t]&-30-11x_{31}-x_{31}^{2}-4x_{31}x_{111}-22x_{111}-4x_{111}^{2}+x_{112}\end{aligned}$\\[3pt]
$T_{91}$ & $\begin{aligned}[t]&-306-35x_{30}-x_{30}^{2}-16x_{30}x_{111}-4x_{30}x_{112}-280x_{111}-64x_{111}^{2}-32x_{111}x_{112}-70x_{112}-4x_{112}^{2}+x_{113}\end{aligned}$\\[3pt]
$T_{92}$ & $\begin{aligned}[t]&-4556-4320x_{111}-1024x_{111}^{2}-512x_{111}x_{112}-128x_{111}x_{113}-1080x_{112}-64x_{112}^{2}-32x_{112}x_{113}-270x_{113}\\ &-4x_{113}^{2}+x_{114}\end{aligned}$\\[3pt]
$T_{93}$ & $\begin{aligned}[t]&-30-11x_{31}-x_{31}^{2}-4x_{31}x_{114}-22x_{114}-4x_{114}^{2}+x_{115}\end{aligned}$\\[3pt]
$T_{94}$ & $\begin{aligned}[t]&-306-35x_{30}-x_{30}^{2}-16x_{30}x_{114}-4x_{30}x_{115}-280x_{114}-64x_{114}^{2}-32x_{114}x_{115}-70x_{115}-4x_{115}^{2}+x_{116}\end{aligned}$\\[3pt]
$T_{95}$ & $\begin{aligned}[t]&-68-x_{110}-32x_{114}-8x_{115}-2x_{116}+x_{117}\end{aligned}$\\[3pt]
$T_{96}$ & $\begin{aligned}[t]&-264-128x_{114}-32x_{115}-8x_{116}-2x_{117}-2x_{117}^{2}+x_{118}\end{aligned}$\\[3pt]
$T_{97}$ & $\begin{aligned}[t]&-2-3x_{29}-x_{29}^{2}-2x_{29}x_{118}-3x_{118}-x_{118}^{2}+x_{119}\end{aligned}$\\[3pt]
$T_{98}$ & $\begin{aligned}[t]&-12-28x_{118}-16x_{118}^{2}-16x_{118}x_{119}-14x_{119}-4x_{119}^{2}+x_{120}\end{aligned}$\\[3pt]
$T_{99}$ & $\begin{aligned}[t]&-30-11x_{31}-x_{31}^{2}-4x_{31}x_{120}-22x_{120}-4x_{120}^{2}+x_{121}\end{aligned}$\\[3pt]
$T_{100}$ & $\begin{aligned}[t]&-306-35x_{29}-x_{29}^{2}-16x_{29}x_{120}-4x_{29}x_{121}-280x_{120}-64x_{120}^{2}-32x_{120}x_{121}-70x_{121}-4x_{121}^{2}+x_{122}\end{aligned}$\\[3pt]
$T_{101}$ & $\begin{aligned}[t]&-4556-4320x_{120}-1024x_{120}^{2}-512x_{120}x_{121}-128x_{120}x_{122}-1080x_{121}-64x_{121}^{2}-32x_{121}x_{122}-270x_{122}\\ &-4x_{122}^{2}+x_{123}\end{aligned}$\\[3pt]
$T_{102}$ & $\begin{aligned}[t]&-90-76x_{31}-16x_{31}^{2}-16x_{31}x_{80}-16x_{31}x_{123}-38x_{80}-4x_{80}^{2}-8x_{80}x_{123}-38x_{123}-4x_{123}^{2}+x_{124}\end{aligned}$\\[3pt]
$T_{103}$ & $\begin{aligned}[t]&-306-35x_{30}-x_{30}^{2}-16x_{30}x_{123}-4x_{30}x_{124}-280x_{123}-64x_{123}^{2}-32x_{123}x_{124}-70x_{124}-4x_{124}^{2}+x_{125}\end{aligned}$\\[3pt]
$T_{104}$ & $\begin{aligned}[t]&-4556-4320x_{114}-1024x_{114}^{2}-512x_{114}x_{115}-128x_{114}x_{116}-1080x_{115}-64x_{115}^{2}-32x_{115}x_{116}-270x_{116}\\ &-4x_{116}^{2}+x_{126}\end{aligned}$\\[3pt]
$T_{105}$ & $\begin{aligned}[t]&-70-32x_{123}-8x_{124}-2x_{125}-2x_{126}+x_{127}\end{aligned}$\\[3pt]
$T_{106}$ & $\begin{aligned}[t]&-x_{127}-x_{127}^{2}+x_{128}\end{aligned}$\\[3pt]
$T_{107}$ & $\begin{aligned}[t]&-306-35x_{31}-x_{31}^{2}-16x_{31}x_{126}-4x_{31}x_{128}-280x_{126}-64x_{126}^{2}-32x_{126}x_{128}-70x_{128}-4x_{128}^{2}+x_{129}\end{aligned}$\\[3pt]
$T_{108}$ & $\begin{aligned}[t]&-42-52x_{30}-16x_{30}^{2}-8x_{30}x_{31}-16x_{30}x_{76}-13x_{31}-x_{31}^{2}-4x_{31}x_{76}-26x_{76}-4x_{76}^{2}+x_{130}\end{aligned}$\\[3pt]
$T_{109}$ & $\begin{aligned}[t]&-756-1100x_{30}-400x_{30}^{2}-400x_{30}x_{76}-80x_{30}x_{130}-550x_{76}-100x_{76}^{2}-40x_{76}x_{130}-110x_{130}-4x_{130}^{2}\\ &+x_{131}\end{aligned}$\\[3pt]
$T_{110}$ & $\begin{aligned}[t]&-88-64x_{30}-32x_{76}-8x_{130}-2x_{131}+x_{132}\end{aligned}$\\[3pt]
$T_{111}$ & $\begin{aligned}[t]&-66-32x_{126}-8x_{128}-2x_{129}-x_{132}+x_{133}\end{aligned}$\\[3pt]
$T_{112}$ & $\begin{aligned}[t]&-264-128x_{126}-32x_{128}-8x_{129}-2x_{133}-2x_{133}^{2}+x_{134}\end{aligned}$\\[3pt]
$T_{113}$ & $\begin{aligned}[t]&-2-3x_{134}-x_{134}^{2}+x_{135}\end{aligned}$\\[3pt]
$T_{114}$ & $\begin{aligned}[t]&-1-x_{3}-x_{3}x_{34}-x_{34}+x_{136}\end{aligned}$\\[3pt]
$T_{115}$ & $\begin{aligned}[t]&-x_{33}^{2}+2x_{33}x_{34}-x_{34}^{2}+x_{137}\end{aligned}$\\[3pt]
$T_{116}$ & $\begin{aligned}[t]&-2x_{79}x_{136}-4x_{136}+x_{138}\end{aligned}$\\[3pt]
$T_{117}$ & $\begin{aligned}[t]&-x_{33}^{2}+2x_{33}x_{34}+2x_{33}x_{138}-x_{34}^{2}-2x_{34}x_{138}-x_{138}^{2}+x_{139}\end{aligned}$\\[3pt]
$T_{118}$ & $\begin{aligned}[t]&-x_{85}x_{136}+x_{140}\end{aligned}$\\[3pt]
$T_{119}$ & $\begin{aligned}[t]&-x_{33}^{2}+2x_{33}x_{34}+2x_{33}x_{140}-x_{34}^{2}-2x_{34}x_{140}-x_{140}^{2}+x_{141}\end{aligned}$\\[3pt]
$T_{120}$ & $\begin{aligned}[t]&-x_{88}x_{136}+x_{142}\end{aligned}$\\[3pt]
$T_{121}$ & $\begin{aligned}[t]&-x_{33}^{2}+2x_{33}x_{34}+2x_{33}x_{142}-x_{34}^{2}-2x_{34}x_{142}-x_{142}^{2}+x_{143}\end{aligned}$\\[3pt]
$T_{122}$ & $\begin{aligned}[t]&-x_{92}x_{136}+x_{144}\end{aligned}$\\[3pt]
$T_{123}$ & $\begin{aligned}[t]&-x_{33}^{2}+2x_{33}x_{34}+2x_{33}x_{144}-x_{34}^{2}-2x_{34}x_{144}-x_{144}^{2}+x_{145}\end{aligned}$\\[3pt]
$T_{124}$ & $\begin{aligned}[t]&-x_{95}x_{136}+x_{146}\end{aligned}$\\[3pt]
$T_{125}$ & $\begin{aligned}[t]&-x_{33}^{2}+2x_{33}x_{34}+2x_{33}x_{146}-x_{34}^{2}-2x_{34}x_{146}-x_{146}^{2}+x_{147}\end{aligned}$\\[3pt]
$T_{126}$ & $\begin{aligned}[t]&-4x_{30}x_{136}-2x_{99}x_{136}-4x_{136}+x_{148}\end{aligned}$\\[3pt]
$T_{127}$ & $\begin{aligned}[t]&-x_{33}^{2}+2x_{33}x_{34}+2x_{33}x_{148}-x_{34}^{2}-2x_{34}x_{148}-x_{148}^{2}+x_{149}\end{aligned}$\\[3pt]
$T_{128}$ & $\begin{aligned}[t]&-4x_{102}x_{136}-2x_{103}x_{136}+x_{150}\end{aligned}$\\[3pt]
$T_{129}$ & $\begin{aligned}[t]&-x_{33}^{2}+2x_{33}x_{34}+2x_{33}x_{150}-x_{34}^{2}-2x_{34}x_{150}-x_{150}^{2}+x_{151}\end{aligned}$\\[3pt]
$T_{130}$ & $\begin{aligned}[t]&-x_{108}x_{136}+x_{152}\end{aligned}$\\[3pt]
$T_{131}$ & $\begin{aligned}[t]&-x_{33}^{2}+2x_{33}x_{34}+2x_{33}x_{152}-x_{34}^{2}-2x_{34}x_{152}-x_{152}^{2}+x_{153}\end{aligned}$\\[3pt]
$T_{132}$ & $\begin{aligned}[t]&-x_{29}^{2}+2x_{29}x_{31}-x_{31}^{2}+x_{154}\end{aligned}$\\[3pt]
$T_{133}$ & $\begin{aligned}[t]&-x_{30}^{2}+2x_{30}x_{31}-x_{31}^{2}+x_{155}\end{aligned}$\\[3pt]
$T_{134}$ & $\begin{aligned}[t]&-x_{154}x_{155}+x_{156}\end{aligned}$\\[3pt]
$T_{135}$ & $\begin{aligned}[t]&-4-4x_{32}-x_{32}^{2}+x_{157}\end{aligned}$\\[3pt]
$T_{136}$ & $\begin{aligned}[t]&-x_{156}^{2}+2x_{156}x_{157}-x_{157}^{2}+x_{158}\end{aligned}$\\[3pt]
$T_{137}$ & $\begin{aligned}[t]&-2x_{135}x_{136}-4x_{136}+x_{159}\end{aligned}$\\[3pt]
$T_{138}$ & $\begin{aligned}[t]&-x_{33}^{2}+2x_{33}x_{34}+2x_{33}x_{159}-x_{34}^{2}-2x_{34}x_{159}-x_{159}^{2}+x_{160}\end{aligned}$\\[3pt]
$T_{139}$ & $\begin{aligned}[t]&-1-x_{3}-x_{3}x_{30}-x_{30}+x_{161}\end{aligned}$\\[3pt]
$T_{140}$ & $\begin{aligned}[t]&-1-x_{34}-x_{34}x_{161}-x_{161}+x_{162}\end{aligned}$\\[3pt]
$T_{141}$ & $\begin{aligned}[t]&-1-x_{3}-x_{3}x_{29}-x_{29}+x_{163}\end{aligned}$\\[3pt]
$T_{142}$ & $\begin{aligned}[t]&-4-4x_{3}-4x_{3}x_{30}-2x_{3}x_{76}-4x_{30}-2x_{76}+x_{164}\end{aligned}$\\[3pt]
$T_{143}$ & $\begin{aligned}[t]&-1-x_{33}-x_{33}x_{163}-x_{163}+x_{165}\end{aligned}$\\[3pt]
$T_{144}$ & $\begin{aligned}[t]&-x_{164}x_{165}-x_{165}+x_{166}\end{aligned}$\\[3pt]
$T_{145}$ & $\begin{aligned}[t]&-x_{162}^{2}+2x_{162}x_{166}-x_{166}^{2}+x_{167}\end{aligned}$\\[3pt]
$T_{146}$ & $\begin{aligned}[t]&-6-6x_{3}-4x_{3}x_{30}-2x_{3}x_{76}-4x_{30}-2x_{76}+x_{168}\end{aligned}$\\[3pt]
$T_{147}$ & $\begin{aligned}[t]&-1-x_{34}-x_{34}x_{168}-x_{168}+x_{169}\end{aligned}$\\[3pt]
$T_{148}$ & $\begin{aligned}[t]&-1-x_{33}-x_{33}x_{161}-x_{161}+x_{170}\end{aligned}$\\[3pt]
$T_{149}$ & $\begin{aligned}[t]&-x_{169}^{2}+2x_{169}x_{170}-x_{170}^{2}+x_{171}\end{aligned}$\\[3pt]
$T_{150}$ & $\begin{aligned}[t]&-x_{137}x_{139}+x_{172}\end{aligned}$\\[3pt]
$T_{151}$ & $\begin{aligned}[t]&-x_{141}x_{172}+x_{173}\end{aligned}$\\[3pt]
$T_{152}$ & $\begin{aligned}[t]&-x_{143}x_{173}+x_{174}\end{aligned}$\\[3pt]
$T_{153}$ & $\begin{aligned}[t]&-x_{145}x_{174}+x_{175}\end{aligned}$\\[3pt]
$T_{154}$ & $\begin{aligned}[t]&-x_{147}x_{175}+x_{176}\end{aligned}$\\[3pt]
$T_{155}$ & $\begin{aligned}[t]&-x_{149}x_{176}+x_{177}\end{aligned}$\\[3pt]
$T_{156}$ & $\begin{aligned}[t]&-x_{151}x_{177}+x_{178}\end{aligned}$\\[3pt]
$T_{157}$ & $\begin{aligned}[t]&-x_{153}x_{178}+x_{179}\end{aligned}$\\[3pt]
$T_{158}$ & $\begin{aligned}[t]&-x_{158}x_{179}-x_{160}x_{179}+x_{180}\end{aligned}$\\[3pt]
$T_{159}$ & $\begin{aligned}[t]&-x_{167}x_{180}+x_{181}\end{aligned}$\\[3pt]
$T_{160}$ & $\begin{aligned}[t]&-x_{2}-x_{2}x_{37}+x_{182}\end{aligned}$\\[3pt]
$T_{161}$ & $\begin{aligned}[t]&-1-x_{35}-x_{35}x_{182}-x_{182}+x_{183}\end{aligned}$\\[3pt]
$T_{162}$ & $\begin{aligned}[t]&-1+2x_{1}-x_{1}^{2}+2x_{1}x_{33}+2x_{1}x_{183}-2x_{33}-x_{33}^{2}-2x_{33}x_{183}+x_{74}x_{186}-x_{171}x_{181}-2x_{183}-x_{183}^{2}-x_{185}\end{aligned}$\\[3pt]
$T_{163}$ & $\begin{aligned}[t]&-1-x_{2}-x_{2}x_{36}-x_{36}-x_{36}x_{182}-x_{182}+x_{184}\end{aligned}$\\[3pt]
$T_{164}$ & $\begin{aligned}[t]&-1+2x_{1}-x_{1}^{2}+2x_{1}x_{34}+2x_{1}x_{184}-2x_{34}-x_{34}^{2}-2x_{34}x_{184}-2x_{184}-x_{184}^{2}+x_{185}\end{aligned}$\\[3pt]
$T_{165}$ & $\begin{aligned}[t]&-4-x_{29}-x_{30}-x_{31}-x_{32}+x_{187}\end{aligned}$\\[3pt]
$T_{166}$ & $\begin{aligned}[t]&-4-x_{33}-x_{34}-x_{35}-x_{36}+8x_{74}-x_{187}-x_{480}^{2}-x_{481}^{2}-x_{482}-x_{482}^{2}\end{aligned}$\\[3pt]
$T_{167}$ & $\begin{aligned}[t]&1+x_{63}-x_{64}x_{74}-x_{74}+x_{284}\end{aligned}$\\[3pt]
$T_{168}$ & $\begin{aligned}[t]&-x_{63}x_{188}-x_{188}-x_{188}x_{284}+x_{193}\end{aligned}$\\[3pt]
$T_{169}$ & $\begin{aligned}[t]&-x_{194}-x_{194}x_{197}+x_{198}\end{aligned}$\\[3pt]
$T_{170}$ & $\begin{aligned}[t]&-x_{198}^{2}+x_{199}\end{aligned}$\\[3pt]
$T_{171}$ & $\begin{aligned}[t]&1-x_{195}^{2}-x_{199}+x_{199}x_{200}\end{aligned}$\\[3pt]
$T_{172}$ & $\begin{aligned}[t]&-1-2x_{194}-x_{194}^{2}+x_{200}\end{aligned}$\\[3pt]
$T_{173}$ & $\begin{aligned}[t]&1-2x_{63}x_{195}-2x_{63}x_{196}+x_{193}x_{201}-4x_{195}-2x_{195}x_{284}-4x_{196}-2x_{196}x_{284}+x_{202}\end{aligned}$\\[3pt]
$T_{174}$ & $\begin{aligned}[t]&-4-4x_{63}-x_{63}^{2}-2x_{63}x_{284}+x_{202}-4x_{284}-x_{284}^{2}\end{aligned}$\\[3pt]
$T_{175}$ & $\begin{aligned}[t]&-x_{204}^{2}+x_{207}\end{aligned}$\\[3pt]
$T_{176}$ & $\begin{aligned}[t]&1-x_{203}^{2}+x_{207}x_{230}\end{aligned}$\\[3pt]
$T_{177}$ & $\begin{aligned}[t]&1-x_{195}^{2}-2x_{195}x_{196}-x_{196}^{2}+x_{230}\end{aligned}$\\[3pt]
$T_{178}$ & $\begin{aligned}[t]&-4x_{204}-4x_{204}x_{205}+x_{208}\end{aligned}$\\[3pt]
$T_{179}$ & $\begin{aligned}[t]&1-x_{195}-x_{196}-x_{203}x_{206}+x_{208}\end{aligned}$\\[3pt]
$T_{180}$ & $\begin{aligned}[t]&-x_{210}^{2}+x_{213}\end{aligned}$\\[3pt]
$T_{181}$ & $\begin{aligned}[t]&1-x_{209}^{2}+x_{213}x_{230}\end{aligned}$\\[3pt]
$T_{182}$ & $\begin{aligned}[t]&-x_{212}^{2}+x_{214}\end{aligned}$\\[3pt]
$T_{183}$ & $\begin{aligned}[t]&1-x_{211}^{2}-x_{214}+x_{214}x_{215}\end{aligned}$\\[3pt]
$T_{184}$ & $\begin{aligned}[t]&-1-2x_{208}-x_{208}^{2}+x_{215}\end{aligned}$\\[3pt]
$T_{185}$ & $\begin{aligned}[t]&-x_{203}x_{216}-x_{209}+x_{211}\end{aligned}$\\[3pt]
$T_{186}$ & $\begin{aligned}[t]&-1-x_{0}-4x_{210}x_{217}+x_{212}\end{aligned}$\\[3pt]
$T_{187}$ & $\begin{aligned}[t]&5+4x_{0}+x_{0}^{2}-2x_{0}x_{195}-2x_{0}x_{196}-4x_{195}-4x_{196}+x_{218}\end{aligned}$\\[3pt]
$T_{188}$ & $\begin{aligned}[t]&-x_{0}x_{210}+x_{4}-x_{209}-2x_{210}+x_{218}x_{219}+x_{220}\end{aligned}$\\[3pt]
$T_{189}$ & $\begin{aligned}[t]&-x_{195}x_{210}-x_{196}x_{210}+x_{220}\end{aligned}$\\[3pt]
$T_{190}$ & $\begin{aligned}[t]&-x_{222}^{2}+x_{226}\end{aligned}$\\[3pt]
$T_{191}$ & $\begin{aligned}[t]&1-x_{221}^{2}+x_{226}x_{230}\end{aligned}$\\[3pt]
$T_{192}$ & $\begin{aligned}[t]&-x_{224}^{2}+x_{227}\end{aligned}$\\[3pt]
$T_{193}$ & $\begin{aligned}[t]&1+x_{215}x_{227}-x_{223}^{2}-x_{227}\end{aligned}$\\[3pt]
$T_{194}$ & $\begin{aligned}[t]&-x_{203}x_{228}-x_{221}+x_{223}\end{aligned}$\\[3pt]
$T_{195}$ & $\begin{aligned}[t]&-x_{222}-4x_{222}x_{229}+x_{224}+x_{225}\end{aligned}$\\[3pt]
$T_{196}$ & $\begin{aligned}[t]&-2x_{195}x_{222}-2x_{196}x_{222}+x_{231}\end{aligned}$\\[3pt]
$T_{197}$ & $\begin{aligned}[t]&-2x_{195}x_{221}-2x_{196}x_{221}+x_{232}\end{aligned}$\\[3pt]
$T_{198}$ & $\begin{aligned}[t]&-x_{222}-2x_{222}x_{230}+x_{233}\end{aligned}$\\[3pt]
$T_{199}$ & $\begin{aligned}[t]&1-4x_{195}-2x_{195}x_{222}+2x_{195}x_{225}-4x_{196}-2x_{196}x_{222}+2x_{196}x_{225}+x_{234}+x_{235}\end{aligned}$\\[3pt]
$T_{200}$ & $\begin{aligned}[t]&-4-4x_{222}-x_{222}^{2}+2x_{222}x_{225}+4x_{225}-x_{225}^{2}+x_{235}\end{aligned}$\\[3pt]
$T_{201}$ & $\begin{aligned}[t]&-x_{222}x_{232}-x_{222}x_{233}+x_{225}x_{232}+x_{225}x_{233}-2x_{232}-2x_{233}+x_{237}\end{aligned}$\\[3pt]
$T_{202}$ & $\begin{aligned}[t]&-x_{62}-x_{195}x_{232}-x_{195}x_{233}-x_{196}x_{232}-x_{196}x_{233}+x_{221}+2x_{221}x_{230}+x_{230}x_{231}-x_{234}x_{236}+x_{237}\end{aligned}$\\[3pt]
$T_{203}$ & $\begin{aligned}[t]&2+2x_{62}+x_{62}^{2}-2x_{62}x_{195}-2x_{62}x_{196}-2x_{195}-2x_{196}+x_{238}\end{aligned}$\\[3pt]
$T_{204}$ & $\begin{aligned}[t]&-x_{62}x_{222}+x_{64}+x_{195}x_{222}+x_{196}x_{222}-x_{221}-x_{222}+x_{238}x_{239}\end{aligned}$\\[3pt]
$T_{205}$ & $\begin{aligned}[t]&2+2x_{65}+x_{65}^{2}-2x_{65}x_{195}-2x_{65}x_{196}-2x_{195}-2x_{196}+x_{240}\end{aligned}$\\[3pt]
$T_{206}$ & $\begin{aligned}[t]&-x_{65}x_{210}+x_{66}-x_{209}-x_{210}+x_{220}+x_{240}x_{241}\end{aligned}$\\[3pt]
$T_{207}$ & $\begin{aligned}[t]&1+x_{61}^{2}-2x_{61}x_{195}-2x_{61}x_{196}+x_{242}\end{aligned}$\\[3pt]
$T_{208}$ & $\begin{aligned}[t]&-x_{61}x_{210}+x_{67}-x_{209}+x_{220}+x_{242}x_{243}\end{aligned}$\\[3pt]
$T_{209}$ & $\begin{aligned}[t]&-2x_{195}x_{210}-2x_{196}x_{210}+x_{244}\end{aligned}$\\[3pt]
$T_{210}$ & $\begin{aligned}[t]&-2x_{195}x_{209}-2x_{196}x_{209}+x_{245}\end{aligned}$\\[3pt]
$T_{211}$ & $\begin{aligned}[t]&-x_{210}-2x_{210}x_{230}+x_{246}\end{aligned}$\\[3pt]
$T_{212}$ & $\begin{aligned}[t]&10+6x_{0}+x_{0}^{2}-2x_{0}x_{195}-2x_{0}x_{196}-6x_{195}-6x_{196}+x_{247}\end{aligned}$\\[3pt]
$T_{213}$ & $\begin{aligned}[t]&x_{0}x_{245}+x_{0}x_{246}-x_{69}-x_{195}x_{245}-x_{195}x_{246}-x_{196}x_{245}-x_{196}x_{246}+x_{209}+2x_{209}x_{230}+x_{230}x_{244}+3x_{245}\\ &+3x_{246}-x_{247}x_{248}\end{aligned}$\\[3pt]
$T_{214}$ & $\begin{aligned}[t]&2+2x_{69}+x_{69}^{2}-2x_{69}x_{195}-2x_{69}x_{196}-2x_{195}-2x_{196}+x_{249}\end{aligned}$\\[3pt]
$T_{215}$ & $\begin{aligned}[t]&-x_{69}x_{210}-x_{209}-x_{210}+x_{220}+x_{249}x_{250}+x_{465}^{2}+x_{466}^{2}+x_{467}+x_{467}^{2}\end{aligned}$\\[3pt]
$T_{216}$ & $\begin{aligned}[t]&1-2x_{63}x_{195}-2x_{63}x_{196}-2x_{195}-2x_{195}x_{284}-2x_{196}-2x_{196}x_{284}+x_{251}+x_{252}\end{aligned}$\\[3pt]
$T_{217}$ & $\begin{aligned}[t]&-1-2x_{63}-x_{63}^{2}-2x_{63}x_{284}+x_{252}-2x_{284}-x_{284}^{2}\end{aligned}$\\[3pt]
$T_{218}$ & $\begin{aligned}[t]&-x_{63}x_{222}+x_{188}+x_{195}x_{222}+x_{196}x_{222}-x_{221}-x_{222}-x_{222}x_{284}+x_{251}x_{253}\end{aligned}$\\[3pt]
$T_{219}$ & $\begin{aligned}[t]&-x_{63}x_{210}+x_{189}-x_{209}-x_{210}-x_{210}x_{284}+x_{220}+x_{251}x_{254}\end{aligned}$\\[3pt]
$T_{220}$ & $\begin{aligned}[t]&-x_{63}x_{209}-x_{63}x_{220}-x_{209}-x_{209}x_{284}-x_{220}-x_{220}x_{284}+x_{256}\end{aligned}$\\[3pt]
$T_{221}$ & $\begin{aligned}[t]&-x_{190}-x_{195}x_{220}-x_{196}x_{220}+x_{210}x_{230}-x_{251}x_{255}+x_{256}\end{aligned}$\\[3pt]
$T_{222}$ & $\begin{aligned}[t]&-x_{258}^{2}+x_{261}\end{aligned}$\\[3pt]
$T_{223}$ & $\begin{aligned}[t]&1+x_{230}x_{261}-x_{257}^{2}\end{aligned}$\\[3pt]
$T_{224}$ & $\begin{aligned}[t]&-x_{260}^{2}+x_{262}\end{aligned}$\\[3pt]
$T_{225}$ & $\begin{aligned}[t]&1+x_{215}x_{262}-x_{259}^{2}-x_{262}\end{aligned}$\\[3pt]
$T_{226}$ & $\begin{aligned}[t]&-x_{203}x_{263}-x_{257}+x_{259}\end{aligned}$\\[3pt]
$T_{227}$ & $\begin{aligned}[t]&-1-x_{2}-4x_{258}x_{264}+x_{260}\end{aligned}$\\[3pt]
$T_{228}$ & $\begin{aligned}[t]&-x_{63}x_{258}+x_{191}+x_{195}x_{258}+x_{196}x_{258}+x_{251}x_{265}-x_{257}-x_{258}-x_{258}x_{284}\end{aligned}$\\[3pt]
$T_{229}$ & $\begin{aligned}[t]&-x_{267}^{2}+x_{270}\end{aligned}$\\[3pt]
$T_{230}$ & $\begin{aligned}[t]&1+x_{230}x_{270}-x_{266}^{2}\end{aligned}$\\[3pt]
$T_{231}$ & $\begin{aligned}[t]&-x_{269}^{2}+x_{271}\end{aligned}$\\[3pt]
$T_{232}$ & $\begin{aligned}[t]&1+x_{215}x_{271}-x_{268}^{2}-x_{271}\end{aligned}$\\[3pt]
$T_{233}$ & $\begin{aligned}[t]&-x_{203}x_{272}-x_{266}+x_{268}\end{aligned}$\\[3pt]
$T_{234}$ & $\begin{aligned}[t]&-1-x_{1}-4x_{267}x_{273}+x_{269}\end{aligned}$\\[3pt]
$T_{235}$ & $\begin{aligned}[t]&-x_{63}x_{267}+x_{192}+x_{195}x_{267}+x_{196}x_{267}+x_{251}x_{274}-x_{266}-x_{267}-x_{267}x_{284}\end{aligned}$\\[3pt]
$T_{236}$ & $\begin{aligned}[t]&-x_{276}^{2}+x_{279}\end{aligned}$\\[3pt]
$T_{237}$ & $\begin{aligned}[t]&1+x_{230}x_{279}-x_{275}^{2}\end{aligned}$\\[3pt]
$T_{238}$ & $\begin{aligned}[t]&-x_{278}^{2}+x_{280}\end{aligned}$\\[3pt]
$T_{239}$ & $\begin{aligned}[t]&1+x_{215}x_{280}-x_{277}^{2}-x_{280}\end{aligned}$\\[3pt]
$T_{240}$ & $\begin{aligned}[t]&-x_{203}x_{281}-x_{275}+x_{277}\end{aligned}$\\[3pt]
$T_{241}$ & $\begin{aligned}[t]&-x_{62}-x_{222}+x_{225}-4x_{276}x_{282}+x_{278}\end{aligned}$\\[3pt]
$T_{242}$ & $\begin{aligned}[t]&-x_{63}x_{188}+x_{63}x_{276}-x_{193}x_{283}-x_{195}x_{276}-x_{196}x_{276}+x_{275}+2x_{276}+x_{276}x_{284}-x_{549}^{2}-x_{550}^{2}-x_{551}\\ &-x_{551}^{2}\end{aligned}$\\[3pt]
$T_{243}$ & $\begin{aligned}[t]&-x_{286}^{2}+x_{289}\end{aligned}$\\[3pt]
$T_{244}$ & $\begin{aligned}[t]&1+x_{230}x_{289}-x_{285}^{2}\end{aligned}$\\[3pt]
$T_{245}$ & $\begin{aligned}[t]&-x_{288}^{2}+x_{290}\end{aligned}$\\[3pt]
$T_{246}$ & $\begin{aligned}[t]&1+x_{215}x_{290}-x_{287}^{2}-x_{290}\end{aligned}$\\[3pt]
$T_{247}$ & $\begin{aligned}[t]&-x_{203}x_{291}-x_{285}+x_{287}\end{aligned}$\\[3pt]
$T_{248}$ & $\begin{aligned}[t]&-1-x_{0}-x_{69}-4x_{286}x_{292}+x_{288}\end{aligned}$\\[3pt]
$T_{249}$ & $\begin{aligned}[t]&-x_{63}x_{189}-x_{189}-x_{189}x_{284}+x_{294}\end{aligned}$\\[3pt]
$T_{250}$ & $\begin{aligned}[t]&x_{63}x_{286}-x_{70}x_{189}-x_{195}x_{286}-x_{196}x_{286}+x_{284}x_{286}+x_{285}+2x_{286}-x_{293}x_{294}-x_{564}^{2}-x_{565}^{2}-x_{566}-x_{566}^{2}\end{aligned}$\\[3pt]
$T_{251}$ & $\begin{aligned}[t]&-x_{296}^{2}+x_{299}\end{aligned}$\\[3pt]
$T_{252}$ & $\begin{aligned}[t]&1+x_{230}x_{299}-x_{295}^{2}\end{aligned}$\\[3pt]
$T_{253}$ & $\begin{aligned}[t]&-x_{298}^{2}+x_{300}\end{aligned}$\\[3pt]
$T_{254}$ & $\begin{aligned}[t]&1+x_{215}x_{300}-x_{297}^{2}-x_{300}\end{aligned}$\\[3pt]
$T_{255}$ & $\begin{aligned}[t]&-x_{203}x_{301}-x_{295}+x_{297}\end{aligned}$\\[3pt]
$T_{256}$ & $\begin{aligned}[t]&-3-x_{0}-x_{66}-4x_{296}x_{302}+x_{298}\end{aligned}$\\[3pt]
$T_{257}$ & $\begin{aligned}[t]&x_{63}x_{296}-x_{71}x_{189}-x_{195}x_{296}-x_{196}x_{296}+x_{284}x_{296}-x_{294}x_{303}+x_{295}+2x_{296}-x_{579}^{2}-x_{580}^{2}-x_{581}-x_{581}^{2}\end{aligned}$\\[3pt]
$T_{258}$ & $\begin{aligned}[t]&-x_{305}^{2}+x_{308}\end{aligned}$\\[3pt]
$T_{259}$ & $\begin{aligned}[t]&1+x_{230}x_{308}-x_{304}^{2}\end{aligned}$\\[3pt]
$T_{260}$ & $\begin{aligned}[t]&-x_{307}^{2}+x_{309}\end{aligned}$\\[3pt]
$T_{261}$ & $\begin{aligned}[t]&1+x_{215}x_{309}-x_{306}^{2}-x_{309}\end{aligned}$\\[3pt]
$T_{262}$ & $\begin{aligned}[t]&-x_{203}x_{310}-x_{304}+x_{306}\end{aligned}$\\[3pt]
$T_{263}$ & $\begin{aligned}[t]&-1-x_{29}-4x_{305}x_{311}+x_{307}\end{aligned}$\\[3pt]
$T_{264}$ & $\begin{aligned}[t]&-x_{74}x_{190}+x_{313}\end{aligned}$\\[3pt]
$T_{265}$ & $\begin{aligned}[t]&-x_{63}x_{305}+x_{195}x_{305}+x_{196}x_{305}-x_{284}x_{305}-x_{304}-2x_{305}+x_{312}x_{313}+x_{594}^{2}+x_{595}^{2}+x_{596}+x_{596}^{2}\end{aligned}$\\[3pt]
$T_{266}$ & $\begin{aligned}[t]&-x_{315}^{2}+x_{318}\end{aligned}$\\[3pt]
$T_{267}$ & $\begin{aligned}[t]&1+x_{230}x_{318}-x_{314}^{2}\end{aligned}$\\[3pt]
$T_{268}$ & $\begin{aligned}[t]&-x_{317}^{2}+x_{319}\end{aligned}$\\[3pt]
$T_{269}$ & $\begin{aligned}[t]&1+x_{215}x_{319}-x_{316}^{2}-x_{319}\end{aligned}$\\[3pt]
$T_{270}$ & $\begin{aligned}[t]&-x_{203}x_{320}-x_{314}+x_{316}\end{aligned}$\\[3pt]
$T_{271}$ & $\begin{aligned}[t]&-1-x_{30}-4x_{315}x_{321}+x_{317}\end{aligned}$\\[3pt]
$T_{272}$ & $\begin{aligned}[t]&-x_{63}x_{315}+x_{195}x_{315}+x_{196}x_{315}-x_{284}x_{315}+x_{313}x_{322}-x_{314}-2x_{315}+x_{606}^{2}+x_{607}^{2}+x_{608}+x_{608}^{2}\end{aligned}$\\[3pt]
$T_{273}$ & $\begin{aligned}[t]&-x_{324}^{2}+x_{327}\end{aligned}$\\[3pt]
$T_{274}$ & $\begin{aligned}[t]&1+x_{230}x_{327}-x_{323}^{2}\end{aligned}$\\[3pt]
$T_{275}$ & $\begin{aligned}[t]&-x_{326}^{2}+x_{328}\end{aligned}$\\[3pt]
$T_{276}$ & $\begin{aligned}[t]&1+x_{215}x_{328}-x_{325}^{2}-x_{328}\end{aligned}$\\[3pt]
$T_{277}$ & $\begin{aligned}[t]&-x_{203}x_{329}-x_{323}+x_{325}\end{aligned}$\\[3pt]
$T_{278}$ & $\begin{aligned}[t]&-1-x_{31}-4x_{324}x_{330}+x_{326}\end{aligned}$\\[3pt]
$T_{279}$ & $\begin{aligned}[t]&-x_{63}x_{324}+x_{195}x_{324}+x_{196}x_{324}-x_{284}x_{324}+x_{313}x_{331}-x_{323}-2x_{324}+x_{618}^{2}+x_{619}^{2}+x_{620}+x_{620}^{2}\end{aligned}$\\[3pt]
$T_{280}$ & $\begin{aligned}[t]&-x_{333}^{2}+x_{336}\end{aligned}$\\[3pt]
$T_{281}$ & $\begin{aligned}[t]&1+x_{230}x_{336}-x_{332}^{2}\end{aligned}$\\[3pt]
$T_{282}$ & $\begin{aligned}[t]&-x_{335}^{2}+x_{337}\end{aligned}$\\[3pt]
$T_{283}$ & $\begin{aligned}[t]&1+x_{215}x_{337}-x_{334}^{2}-x_{337}\end{aligned}$\\[3pt]
$T_{284}$ & $\begin{aligned}[t]&-x_{203}x_{338}-x_{332}+x_{334}\end{aligned}$\\[3pt]
$T_{285}$ & $\begin{aligned}[t]&-1-x_{32}-4x_{333}x_{339}+x_{335}\end{aligned}$\\[3pt]
$T_{286}$ & $\begin{aligned}[t]&-x_{63}x_{333}+x_{195}x_{333}+x_{196}x_{333}-x_{284}x_{333}+x_{313}x_{340}-x_{332}-2x_{333}+x_{630}^{2}+x_{631}^{2}+x_{632}+x_{632}^{2}\end{aligned}$\\[3pt]
$T_{287}$ & $\begin{aligned}[t]&-x_{342}^{2}+x_{345}\end{aligned}$\\[3pt]
$T_{288}$ & $\begin{aligned}[t]&1+x_{230}x_{345}-x_{341}^{2}\end{aligned}$\\[3pt]
$T_{289}$ & $\begin{aligned}[t]&-x_{344}^{2}+x_{346}\end{aligned}$\\[3pt]
$T_{290}$ & $\begin{aligned}[t]&1+x_{215}x_{346}-x_{343}^{2}-x_{346}\end{aligned}$\\[3pt]
$T_{291}$ & $\begin{aligned}[t]&-x_{203}x_{347}-x_{341}+x_{343}\end{aligned}$\\[3pt]
$T_{292}$ & $\begin{aligned}[t]&-1-x_{33}-4x_{342}x_{348}+x_{344}\end{aligned}$\\[3pt]
$T_{293}$ & $\begin{aligned}[t]&-x_{74}x_{191}+x_{350}\end{aligned}$\\[3pt]
$T_{294}$ & $\begin{aligned}[t]&-x_{63}x_{342}+x_{195}x_{342}+x_{196}x_{342}-x_{284}x_{342}-x_{341}-2x_{342}+x_{349}x_{350}+x_{642}^{2}+x_{643}^{2}+x_{644}+x_{644}^{2}\end{aligned}$\\[3pt]
$T_{295}$ & $\begin{aligned}[t]&-x_{352}^{2}+x_{355}\end{aligned}$\\[3pt]
$T_{296}$ & $\begin{aligned}[t]&1+x_{230}x_{355}-x_{351}^{2}\end{aligned}$\\[3pt]
$T_{297}$ & $\begin{aligned}[t]&-x_{354}^{2}+x_{356}\end{aligned}$\\[3pt]
$T_{298}$ & $\begin{aligned}[t]&1+x_{215}x_{356}-x_{353}^{2}-x_{356}\end{aligned}$\\[3pt]
$T_{299}$ & $\begin{aligned}[t]&-x_{203}x_{357}-x_{351}+x_{353}\end{aligned}$\\[3pt]
$T_{300}$ & $\begin{aligned}[t]&-1-x_{34}-4x_{352}x_{358}+x_{354}\end{aligned}$\\[3pt]
$T_{301}$ & $\begin{aligned}[t]&-x_{63}x_{352}+x_{195}x_{352}+x_{196}x_{352}-x_{284}x_{352}+x_{350}x_{359}-x_{351}-2x_{352}+x_{654}^{2}+x_{655}^{2}+x_{656}+x_{656}^{2}\end{aligned}$\\[3pt]
$T_{302}$ & $\begin{aligned}[t]&-x_{361}^{2}+x_{364}\end{aligned}$\\[3pt]
$T_{303}$ & $\begin{aligned}[t]&1+x_{230}x_{364}-x_{360}^{2}\end{aligned}$\\[3pt]
$T_{304}$ & $\begin{aligned}[t]&-x_{363}^{2}+x_{365}\end{aligned}$\\[3pt]
$T_{305}$ & $\begin{aligned}[t]&1+x_{215}x_{365}-x_{362}^{2}-x_{365}\end{aligned}$\\[3pt]
$T_{306}$ & $\begin{aligned}[t]&-x_{203}x_{366}-x_{360}+x_{362}\end{aligned}$\\[3pt]
$T_{307}$ & $\begin{aligned}[t]&-1-x_{35}-4x_{361}x_{367}+x_{363}\end{aligned}$\\[3pt]
$T_{308}$ & $\begin{aligned}[t]&-x_{74}x_{192}+x_{369}\end{aligned}$\\[3pt]
$T_{309}$ & $\begin{aligned}[t]&-x_{63}x_{361}+x_{195}x_{361}+x_{196}x_{361}-x_{284}x_{361}-x_{360}-2x_{361}+x_{368}x_{369}+x_{666}^{2}+x_{667}^{2}+x_{668}+x_{668}^{2}\end{aligned}$\\[3pt]
$T_{310}$ & $\begin{aligned}[t]&-x_{371}^{2}+x_{374}\end{aligned}$\\[3pt]
$T_{311}$ & $\begin{aligned}[t]&1+x_{230}x_{374}-x_{370}^{2}\end{aligned}$\\[3pt]
$T_{312}$ & $\begin{aligned}[t]&-x_{373}^{2}+x_{375}\end{aligned}$\\[3pt]
$T_{313}$ & $\begin{aligned}[t]&1+x_{215}x_{375}-x_{372}^{2}-x_{375}\end{aligned}$\\[3pt]
$T_{314}$ & $\begin{aligned}[t]&-x_{203}x_{376}-x_{370}+x_{372}\end{aligned}$\\[3pt]
$T_{315}$ & $\begin{aligned}[t]&-1-x_{36}-4x_{371}x_{377}+x_{373}\end{aligned}$\\[3pt]
$T_{316}$ & $\begin{aligned}[t]&-x_{63}x_{371}+x_{195}x_{371}+x_{196}x_{371}-x_{284}x_{371}+x_{369}x_{378}-x_{370}-2x_{371}+x_{678}^{2}+x_{679}^{2}+x_{680}+x_{680}^{2}\end{aligned}$\\[3pt]
$T_{317}$ & $\begin{aligned}[t]&-x_{213}x_{226}+x_{379}\end{aligned}$\\[3pt]
$T_{318}$ & $\begin{aligned}[t]&-x_{261}x_{379}+x_{380}\end{aligned}$\\[3pt]
$T_{319}$ & $\begin{aligned}[t]&-x_{270}x_{380}+x_{381}\end{aligned}$\\[3pt]
$T_{320}$ & $\begin{aligned}[t]&-x_{279}x_{381}+x_{382}\end{aligned}$\\[3pt]
$T_{321}$ & $\begin{aligned}[t]&-x_{289}x_{382}+x_{383}\end{aligned}$\\[3pt]
$T_{322}$ & $\begin{aligned}[t]&-x_{299}x_{383}+x_{384}\end{aligned}$\\[3pt]
$T_{323}$ & $\begin{aligned}[t]&-x_{308}x_{384}+x_{385}\end{aligned}$\\[3pt]
$T_{324}$ & $\begin{aligned}[t]&-x_{318}x_{385}+x_{386}\end{aligned}$\\[3pt]
$T_{325}$ & $\begin{aligned}[t]&-x_{327}x_{386}+x_{387}\end{aligned}$\\[3pt]
$T_{326}$ & $\begin{aligned}[t]&-x_{336}x_{387}+x_{388}\end{aligned}$\\[3pt]
$T_{327}$ & $\begin{aligned}[t]&-x_{345}x_{388}+x_{389}\end{aligned}$\\[3pt]
$T_{328}$ & $\begin{aligned}[t]&-x_{355}x_{389}+x_{390}\end{aligned}$\\[3pt]
$T_{329}$ & $\begin{aligned}[t]&-x_{364}x_{390}+x_{391}\end{aligned}$\\[3pt]
$T_{330}$ & $\begin{aligned}[t]&-x_{374}x_{391}+x_{392}\end{aligned}$\\[3pt]
$T_{331}$ & $\begin{aligned}[t]&x_{204}-x_{392}-x_{392}x_{393}\end{aligned}$\\[3pt]
$T_{332}$ & $\begin{aligned}[t]&-9-3x_{0}-x_{2}-x_{4}-x_{222}+x_{225}+x_{394}\end{aligned}$\\[3pt]
$T_{333}$ & $\begin{aligned}[t]&-2-x_{61}-x_{62}-x_{222}+x_{225}-x_{394}+x_{395}\end{aligned}$\\[3pt]
$T_{334}$ & $\begin{aligned}[t]&-2-x_{62}-x_{64}-x_{65}-x_{395}+x_{396}\end{aligned}$\\[3pt]
$T_{335}$ & $\begin{aligned}[t]&-1-x_{66}-x_{67}-2x_{69}-x_{396}+x_{397}\end{aligned}$\\[3pt]
$T_{336}$ & $\begin{aligned}[t]&-2-x_{1}-x_{63}-x_{284}-x_{397}+x_{398}-x_{465}^{2}-x_{466}^{2}-x_{467}-x_{467}^{2}\end{aligned}$\\[3pt]
$T_{337}$ & $\begin{aligned}[t]&-x_{188}-x_{189}-x_{190}-x_{191}-x_{192}+x_{194}-x_{398}\end{aligned}$\\[3pt]
$T_{338}$ & $\begin{aligned}[t]&x_{0}-x_{399}^{2}-x_{400}^{2}-x_{401}-x_{401}^{2}\end{aligned}$\\[3pt]
$T_{339}$ & $\begin{aligned}[t]&x_{3}-x_{402}^{2}-x_{403}^{2}-x_{404}-x_{404}^{2}\end{aligned}$\\[3pt]
$T_{340}$ & $\begin{aligned}[t]&x_{4}-x_{405}^{2}-x_{406}^{2}-x_{407}-x_{407}^{2}\end{aligned}$\\[3pt]
$T_{341}$ & $\begin{aligned}[t]&-2+x_{2}-x_{3}-x_{408}^{2}-x_{409}^{2}-x_{410}-x_{410}^{2}\end{aligned}$\\[3pt]
$T_{342}$ & $\begin{aligned}[t]&x_{27}-x_{411}^{2}-x_{412}^{2}-x_{413}-x_{413}^{2}\end{aligned}$\\[3pt]
$T_{343}$ & $\begin{aligned}[t]&x_{29}-x_{414}^{2}-x_{415}^{2}-x_{416}-x_{416}^{2}\end{aligned}$\\[3pt]
$T_{344}$ & $\begin{aligned}[t]&x_{30}-x_{417}^{2}-x_{418}^{2}-x_{419}-x_{419}^{2}\end{aligned}$\\[3pt]
$T_{345}$ & $\begin{aligned}[t]&x_{31}-x_{420}^{2}-x_{421}^{2}-x_{422}-x_{422}^{2}\end{aligned}$\\[3pt]
$T_{346}$ & $\begin{aligned}[t]&x_{32}-x_{423}^{2}-x_{424}^{2}-x_{425}-x_{425}^{2}\end{aligned}$\\[3pt]
$T_{347}$ & $\begin{aligned}[t]&x_{33}-x_{426}^{2}-x_{427}^{2}-x_{428}-x_{428}^{2}\end{aligned}$\\[3pt]
$T_{348}$ & $\begin{aligned}[t]&x_{34}-x_{429}^{2}-x_{430}^{2}-x_{431}-x_{431}^{2}\end{aligned}$\\[3pt]
$T_{349}$ & $\begin{aligned}[t]&x_{35}-x_{432}^{2}-x_{433}^{2}-x_{434}-x_{434}^{2}\end{aligned}$\\[3pt]
$T_{350}$ & $\begin{aligned}[t]&x_{36}-x_{435}^{2}-x_{436}^{2}-x_{437}-x_{437}^{2}\end{aligned}$\\[3pt]
$T_{351}$ & $\begin{aligned}[t]&x_{37}-x_{438}^{2}-x_{439}^{2}-x_{440}-x_{440}^{2}\end{aligned}$\\[3pt]
$T_{352}$ & $\begin{aligned}[t]&x_{62}-x_{441}^{2}-x_{442}^{2}-x_{443}-x_{443}^{2}\end{aligned}$\\[3pt]
$T_{353}$ & $\begin{aligned}[t]&x_{64}-x_{444}^{2}-x_{445}^{2}-x_{446}-x_{446}^{2}\end{aligned}$\\[3pt]
$T_{354}$ & $\begin{aligned}[t]&-1-x_{61}x_{63}+x_{63}+x_{64}-x_{447}^{2}-x_{448}^{2}-x_{449}-x_{449}^{2}\end{aligned}$\\[3pt]
$T_{355}$ & $\begin{aligned}[t]&x_{61}x_{63}-x_{64}-x_{450}^{2}-x_{451}^{2}-x_{452}-x_{452}^{2}\end{aligned}$\\[3pt]
$T_{356}$ & $\begin{aligned}[t]&x_{66}-x_{453}^{2}-x_{454}^{2}-x_{455}-x_{455}^{2}\end{aligned}$\\[3pt]
$T_{357}$ & $\begin{aligned}[t]&x_{67}-x_{456}^{2}-x_{457}^{2}-x_{458}-x_{458}^{2}\end{aligned}$\\[3pt]
$T_{358}$ & $\begin{aligned}[t]&x_{68}-x_{459}^{2}-x_{460}^{2}-x_{461}-x_{461}^{2}\end{aligned}$\\[3pt]
$T_{359}$ & $\begin{aligned}[t]&x_{69}-x_{462}^{2}-x_{463}^{2}-x_{464}-x_{464}^{2}\end{aligned}$\\[3pt]
$T_{360}$ & $\begin{aligned}[t]&-1-x_{68}x_{70}+x_{70}+x_{465}^{2}+x_{466}^{2}+x_{467}+x_{467}^{2}-x_{468}^{2}-x_{469}^{2}-x_{470}-x_{470}^{2}\end{aligned}$\\[3pt]
$T_{361}$ & $\begin{aligned}[t]&x_{68}x_{70}-x_{465}^{2}-x_{466}^{2}-x_{467}-x_{467}^{2}-x_{471}^{2}-x_{472}^{2}-x_{473}-x_{473}^{2}\end{aligned}$\\[3pt]
$T_{362}$ & $\begin{aligned}[t]&x_{71}-x_{474}^{2}-x_{475}^{2}-x_{476}-x_{476}^{2}\end{aligned}$\\[3pt]
$T_{363}$ & $\begin{aligned}[t]&-1+x_{73}-x_{74}-x_{477}^{2}-x_{478}^{2}-x_{479}-x_{479}^{2}\end{aligned}$\\[3pt]
$T_{364}$ & $\begin{aligned}[t]&x_{195}-x_{483}^{2}-x_{484}^{2}-x_{485}-x_{485}^{2}\end{aligned}$\\[3pt]
$T_{365}$ & $\begin{aligned}[t]&x_{196}-x_{486}^{2}-x_{487}^{2}-x_{488}-x_{488}^{2}\end{aligned}$\\[3pt]
$T_{366}$ & $\begin{aligned}[t]&x_{197}-x_{489}^{2}-x_{490}^{2}-x_{491}-x_{491}^{2}\end{aligned}$\\[3pt]
$T_{367}$ & $\begin{aligned}[t]&x_{205}-x_{492}^{2}-x_{493}^{2}-x_{494}-x_{494}^{2}\end{aligned}$\\[3pt]
$T_{368}$ & $\begin{aligned}[t]&x_{209}-x_{495}^{2}-x_{496}^{2}-x_{497}-x_{497}^{2}\end{aligned}$\\[3pt]
$T_{369}$ & $\begin{aligned}[t]&-1-x_{0}+x_{210}-x_{498}^{2}-x_{499}^{2}-x_{500}-x_{500}^{2}\end{aligned}$\\[3pt]
$T_{370}$ & $\begin{aligned}[t]&x_{218}-x_{501}^{2}-x_{502}^{2}-x_{503}-x_{503}^{2}\end{aligned}$\\[3pt]
$T_{371}$ & $\begin{aligned}[t]&x_{221}-x_{504}^{2}-x_{505}^{2}-x_{506}-x_{506}^{2}\end{aligned}$\\[3pt]
$T_{372}$ & $\begin{aligned}[t]&x_{225}-x_{507}^{2}-x_{508}^{2}-x_{509}-x_{509}^{2}\end{aligned}$\\[3pt]
$T_{373}$ & $\begin{aligned}[t]&x_{234}-x_{510}^{2}-x_{511}^{2}-x_{512}-x_{512}^{2}\end{aligned}$\\[3pt]
$T_{374}$ & $\begin{aligned}[t]&x_{238}-x_{513}^{2}-x_{514}^{2}-x_{515}-x_{515}^{2}\end{aligned}$\\[3pt]
$T_{375}$ & $\begin{aligned}[t]&x_{240}-x_{516}^{2}-x_{517}^{2}-x_{518}-x_{518}^{2}\end{aligned}$\\[3pt]
$T_{376}$ & $\begin{aligned}[t]&x_{242}-x_{519}^{2}-x_{520}^{2}-x_{521}-x_{521}^{2}\end{aligned}$\\[3pt]
$T_{377}$ & $\begin{aligned}[t]&x_{247}-x_{522}^{2}-x_{523}^{2}-x_{524}-x_{524}^{2}\end{aligned}$\\[3pt]
$T_{378}$ & $\begin{aligned}[t]&x_{249}-x_{525}^{2}-x_{526}^{2}-x_{527}-x_{527}^{2}\end{aligned}$\\[3pt]
$T_{379}$ & $\begin{aligned}[t]&x_{251}-x_{528}^{2}-x_{529}^{2}-x_{530}-x_{530}^{2}\end{aligned}$\\[3pt]
$T_{380}$ & $\begin{aligned}[t]&x_{257}-x_{531}^{2}-x_{532}^{2}-x_{533}-x_{533}^{2}\end{aligned}$\\[3pt]
$T_{381}$ & $\begin{aligned}[t]&-1-x_{2}+x_{258}-x_{534}^{2}-x_{535}^{2}-x_{536}-x_{536}^{2}\end{aligned}$\\[3pt]
$T_{382}$ & $\begin{aligned}[t]&x_{266}-x_{537}^{2}-x_{538}^{2}-x_{539}-x_{539}^{2}\end{aligned}$\\[3pt]
$T_{383}$ & $\begin{aligned}[t]&-1-x_{1}+x_{267}-x_{540}^{2}-x_{541}^{2}-x_{542}-x_{542}^{2}\end{aligned}$\\[3pt]
$T_{384}$ & $\begin{aligned}[t]&x_{275}-x_{543}^{2}-x_{544}^{2}-x_{545}-x_{545}^{2}\end{aligned}$\\[3pt]
$T_{385}$ & $\begin{aligned}[t]&-x_{62}-x_{222}+x_{225}+x_{276}-x_{546}^{2}-x_{547}^{2}-x_{548}-x_{548}^{2}\end{aligned}$\\[3pt]
$T_{386}$ & $\begin{aligned}[t]&-1+x_{188}-x_{549}^{2}-x_{550}^{2}-x_{551}-x_{551}^{2}-x_{552}^{2}-x_{553}^{2}-x_{554}-x_{554}^{2}\end{aligned}$\\[3pt]
$T_{387}$ & $\begin{aligned}[t]&x_{284}-x_{555}^{2}-x_{556}^{2}-x_{557}-x_{557}^{2}\end{aligned}$\\[3pt]
$T_{388}$ & $\begin{aligned}[t]&x_{285}-x_{558}^{2}-x_{559}^{2}-x_{560}-x_{560}^{2}\end{aligned}$\\[3pt]
$T_{389}$ & $\begin{aligned}[t]&-1-x_{0}-x_{69}+x_{286}-x_{561}^{2}-x_{562}^{2}-x_{563}-x_{563}^{2}\end{aligned}$\\[3pt]
$T_{390}$ & $\begin{aligned}[t]&-1+x_{189}-x_{564}^{2}-x_{565}^{2}-x_{566}-x_{566}^{2}-x_{567}^{2}-x_{568}^{2}-x_{569}-x_{569}^{2}\end{aligned}$\\[3pt]
$T_{391}$ & $\begin{aligned}[t]&x_{63}-x_{70}+x_{284}-x_{570}^{2}-x_{571}^{2}-x_{572}-x_{572}^{2}\end{aligned}$\\[3pt]
$T_{392}$ & $\begin{aligned}[t]&x_{295}-x_{573}^{2}-x_{574}^{2}-x_{575}-x_{575}^{2}\end{aligned}$\\[3pt]
$T_{393}$ & $\begin{aligned}[t]&-3-x_{0}-x_{66}+x_{296}-x_{576}^{2}-x_{577}^{2}-x_{578}-x_{578}^{2}\end{aligned}$\\[3pt]
$T_{394}$ & $\begin{aligned}[t]&-1+x_{189}-x_{579}^{2}-x_{580}^{2}-x_{581}-x_{581}^{2}-x_{582}^{2}-x_{583}^{2}-x_{584}-x_{584}^{2}\end{aligned}$\\[3pt]
$T_{395}$ & $\begin{aligned}[t]&x_{63}-x_{71}+x_{284}-x_{585}^{2}-x_{586}^{2}-x_{587}-x_{587}^{2}\end{aligned}$\\[3pt]
$T_{396}$ & $\begin{aligned}[t]&x_{304}-x_{588}^{2}-x_{589}^{2}-x_{590}-x_{590}^{2}\end{aligned}$\\[3pt]
$T_{397}$ & $\begin{aligned}[t]&-1-x_{29}+x_{305}-x_{591}^{2}-x_{592}^{2}-x_{593}-x_{593}^{2}\end{aligned}$\\[3pt]
$T_{398}$ & $\begin{aligned}[t]&-1+x_{190}-x_{594}^{2}-x_{595}^{2}-x_{596}-x_{596}^{2}-x_{597}^{2}-x_{598}^{2}-x_{599}-x_{599}^{2}\end{aligned}$\\[3pt]
$T_{399}$ & $\begin{aligned}[t]&x_{314}-x_{600}^{2}-x_{601}^{2}-x_{602}-x_{602}^{2}\end{aligned}$\\[3pt]
$T_{400}$ & $\begin{aligned}[t]&-1-x_{30}+x_{315}-x_{603}^{2}-x_{604}^{2}-x_{605}-x_{605}^{2}\end{aligned}$\\[3pt]
$T_{401}$ & $\begin{aligned}[t]&-1+x_{190}-x_{606}^{2}-x_{607}^{2}-x_{608}-x_{608}^{2}-x_{609}^{2}-x_{610}^{2}-x_{611}-x_{611}^{2}\end{aligned}$\\[3pt]
$T_{402}$ & $\begin{aligned}[t]&x_{323}-x_{612}^{2}-x_{613}^{2}-x_{614}-x_{614}^{2}\end{aligned}$\\[3pt]
$T_{403}$ & $\begin{aligned}[t]&-1-x_{31}+x_{324}-x_{615}^{2}-x_{616}^{2}-x_{617}-x_{617}^{2}\end{aligned}$\\[3pt]
$T_{404}$ & $\begin{aligned}[t]&-1+x_{190}-x_{618}^{2}-x_{619}^{2}-x_{620}-x_{620}^{2}-x_{621}^{2}-x_{622}^{2}-x_{623}-x_{623}^{2}\end{aligned}$\\[3pt]
$T_{405}$ & $\begin{aligned}[t]&x_{332}-x_{624}^{2}-x_{625}^{2}-x_{626}-x_{626}^{2}\end{aligned}$\\[3pt]
$T_{406}$ & $\begin{aligned}[t]&-1-x_{32}+x_{333}-x_{627}^{2}-x_{628}^{2}-x_{629}-x_{629}^{2}\end{aligned}$\\[3pt]
$T_{407}$ & $\begin{aligned}[t]&-1+x_{190}-x_{630}^{2}-x_{631}^{2}-x_{632}-x_{632}^{2}-x_{633}^{2}-x_{634}^{2}-x_{635}-x_{635}^{2}\end{aligned}$\\[3pt]
$T_{408}$ & $\begin{aligned}[t]&x_{341}-x_{636}^{2}-x_{637}^{2}-x_{638}-x_{638}^{2}\end{aligned}$\\[3pt]
$T_{409}$ & $\begin{aligned}[t]&-1-x_{33}+x_{342}-x_{639}^{2}-x_{640}^{2}-x_{641}-x_{641}^{2}\end{aligned}$\\[3pt]
$T_{410}$ & $\begin{aligned}[t]&-1+x_{191}-x_{642}^{2}-x_{643}^{2}-x_{644}-x_{644}^{2}-x_{645}^{2}-x_{646}^{2}-x_{647}-x_{647}^{2}\end{aligned}$\\[3pt]
$T_{411}$ & $\begin{aligned}[t]&x_{351}-x_{648}^{2}-x_{649}^{2}-x_{650}-x_{650}^{2}\end{aligned}$\\[3pt]
$T_{412}$ & $\begin{aligned}[t]&-1-x_{34}+x_{352}-x_{651}^{2}-x_{652}^{2}-x_{653}-x_{653}^{2}\end{aligned}$\\[3pt]
$T_{413}$ & $\begin{aligned}[t]&-1+x_{191}-x_{654}^{2}-x_{655}^{2}-x_{656}-x_{656}^{2}-x_{657}^{2}-x_{658}^{2}-x_{659}-x_{659}^{2}\end{aligned}$\\[3pt]
$T_{414}$ & $\begin{aligned}[t]&x_{360}-x_{660}^{2}-x_{661}^{2}-x_{662}-x_{662}^{2}\end{aligned}$\\[3pt]
$T_{415}$ & $\begin{aligned}[t]&-1-x_{35}+x_{361}-x_{663}^{2}-x_{664}^{2}-x_{665}-x_{665}^{2}\end{aligned}$\\[3pt]
$T_{416}$ & $\begin{aligned}[t]&-1+x_{192}-x_{666}^{2}-x_{667}^{2}-x_{668}-x_{668}^{2}-x_{669}^{2}-x_{670}^{2}-x_{671}-x_{671}^{2}\end{aligned}$\\[3pt]
$T_{417}$ & $\begin{aligned}[t]&x_{370}-x_{672}^{2}-x_{673}^{2}-x_{674}-x_{674}^{2}\end{aligned}$\\[3pt]
$T_{418}$ & $\begin{aligned}[t]&-1-x_{36}+x_{371}-x_{675}^{2}-x_{676}^{2}-x_{677}-x_{677}^{2}\end{aligned}$\\[3pt]
$T_{419}$ & $\begin{aligned}[t]&-1+x_{192}-x_{678}^{2}-x_{679}^{2}-x_{680}-x_{680}^{2}-x_{681}^{2}-x_{682}^{2}-x_{683}-x_{683}^{2}\end{aligned}$\\[3pt]
$T_{420}$ & $\begin{aligned}[t]&x_{393}-x_{684}^{2}-x_{685}^{2}-x_{686}-x_{686}^{2}\end{aligned}$\\[3pt]
\end{longtable}
\endgroup

\end{document}
